\documentclass[10pt,twoside,openany]{book}

% --- Page Geometry (6.5 x 9.5 in Royal/Academic Monograph Format) ---
\usepackage[
  paperwidth=6.5in,
  paperheight=9.5in,
  top=0.85in,
  bottom=0.90in,
  inner=0.85in,
  outer=0.75in,
  headheight=14pt,
  footskip=28pt
]{geometry}

% --- Core Typography & Math Packages ---
\usepackage{amsmath,amssymb,amsfonts,amsthm,mathtools}
\usepackage{microtype}
\usepackage{palatino}
\usepackage{mathpazo}
\linespread{1.06}
\emergencystretch=1.5em

% --- Colors, Graphics, TikZ (for Vector Cover & Architectural Diagrams) ---
\usepackage[dvipsnames,table]{xcolor}
\usepackage{tikz}
\usetikzlibrary{calc,positioning,arrows.meta,shapes.geometric,backgrounds,fit}

\definecolor{NavyPrimary}{HTML}{0B192C}
\definecolor{SlateDark}{HTML}{1E3E62}
\definecolor{CrimsonAccent}{HTML}{A61C1C}
\definecolor{GoldAccent}{HTML}{D4AF37}
\definecolor{TealAccent}{HTML}{0F766E}
\definecolor{SoftBg}{HTML}{F8FAFC}
\definecolor{CodeBg}{HTML}{F1F5F9}
\definecolor{BorderGray}{HTML}{CBD5E1}
\definecolor{RuleGray}{HTML}{94A3B8}
\tikzset{
  dbox/.style={rectangle, rounded corners=2pt, draw=NavyPrimary, fill=SoftBg, thick, align=flush center, font=\tiny, inner sep=2.5pt},
  dhead/.style={dbox, fill=NavyPrimary!10!white, font=\tiny\bfseries},
  dcrypto/.style={dbox, draw=TealAccent, fill=TealAccent!8!white},
  dinst/.style={dbox, draw=GoldAccent!80!black, fill=GoldAccent!14!white},
  dwarn/.style={dbox, draw=CrimsonAccent, fill=CrimsonAccent!7!white},
  darr/.style={-{Stealth[length=1.7mm]}, thick, color=SlateDark},
  dlab/.style={font=\tiny, fill=white, inner sep=1pt, align=flush center},
  dlife/.style={RuleGray, dashed, thin}
}

% --- Tables, Lists, Code Listings & Callout Boxes ---
\usepackage{booktabs}
\usepackage{longtable}
\usepackage{tabularx}
\usepackage{placeins}
\usepackage{array}
\usepackage{multirow}
\usepackage{enumitem}
\usepackage{listings}
\usepackage{fancyhdr}
\usepackage{titlesec}
\usepackage{xurl}
\usepackage[bottom,hang,flushmargin]{footmisc}
\usepackage[
  colorlinks=true,
  linkcolor=SlateDark,
  citecolor=CrimsonAccent,
  urlcolor=TealAccent,
  bookmarks=true,
  bookmarksnumbered=true,
  pdftitle={Distributed Participatory Economics Protocol (DPEP) — Working Draft Monograph},
  pdfauthor={Post-Capitalist Labs (https://github.com/Post-Capitalist-Labs)},
  pdfsubject={Participatory Economics, Post-Capitalist Labs, Distributed Systems, Post-Quantum Cryptography}
]{hyperref}

% --- Footnote Formatting ---
\setlength{\footnotemargin}{1.1em}
\renewcommand{\footnoterule}{\kern-3pt \hrule width 0.35\textwidth height 0.4pt \kern 2.6pt}

% --- Custom Table Column Types ---
\newcolumntype{L}[1]{>{\raggedright\arraybackslash}p{#1}}
\newcolumntype{C}[1]{>{\centering\arraybackslash}p{#1}}
\newcolumntype{R}[1]{>{\raggedleft\arraybackslash}p{#1}}
\newcolumntype{Y}{>{\raggedright\arraybackslash}X}
\renewcommand{\arraystretch}{1.22}

% --- Code Listing Styling ---
\lstdefinestyle{dpepcode}{
  backgroundcolor=\color{CodeBg},
  basicstyle=\ttfamily\scriptsize,
  breakatwhitespace=false,
  breaklines=true,
  captionpos=b,
  commentstyle=\color{TealAccent}\itshape,
  keywordstyle=\color{CrimsonAccent}\bfseries,
  numberstyle=\tiny\color{RuleGray},
  numbers=left,
  numbersep=6pt,
  frame=single,
  rulecolor=\color{BorderGray},
  showspaces=false,
  showstringspaces=false,
  showtabs=false,
  tabsize=2,
  xleftmargin=12pt,
  xrightmargin=4pt,
  aboveskip=8pt,
  belowskip=8pt
}
\lstset{style=dpepcode}

% --- Theorem & Box Environments ---
\newtheoremstyle{dpepstyle}
  {8pt}{8pt}{\normalfont}{0pt}{\bfseries\color{NavyPrimary}}{.}{0.5em}{}
\theoremstyle{dpepstyle}
\newtheorem{definition}{Definition}[chapter]
\newtheorem{theorem}{Theorem}[chapter]
\newtheorem{invariant}{Protocol Invariant}[chapter]
\newtheorem{designrule}{Design Principle}[chapter]

% --- Epigraph Environment ---
\newcommand{\chapterepigraph}[2]{%
  \vspace{-0.4em}
  \begin{FlushRight}
  \begin{minipage}{0.82\textwidth}
    \itshape\small\color{SlateDark} ``#1''\par
    \vspace{3pt}
    \upshape\footnotesize\textbf{--- #2}
  \end{minipage}
  \end{FlushRight}
  \vspace{1em}
}
\newenvironment{FlushRight}{\begin{flushright}}{\end{flushright}}

% --- Audience Callout Macro (For Advocates, Economists, Programmers) ---
\newcommand{\audiencebox}[3]{%
  \par\smallskip\noindent
  \colorbox{SoftBg}{%
    \parbox{\dimexpr\textwidth-2\fboxsep\relax}{%
      \vspace{5pt}
      {\footnotesize\textbf{\textcolor{CrimsonAccent}{\MakeUppercase{#1}} --- \textcolor{NavyPrimary}{#2}}}\par\vspace{3pt}
      {\small #3}
      \vspace{5pt}
    }%
  }\par\smallskip
}

% --- Headers and Footers ---
\pagestyle{fancy}
\fancyhf{}
\fancyhead[LE]{\footnotesize\scshape\textcolor{SlateDark}{Distributed Participatory Economics Protocol}}
\fancyhead[RO]{\footnotesize\scshape\textcolor{SlateDark}{\nouppercase{\rightmark}}}
\fancyfoot[LE,RO]{\small\textbf{\thepage}}
\fancyfoot[LO,RE]{\scriptsize\textcolor{RuleGray}{Post-Capitalist Labs Imprint $\cdot$ Working Draft Monograph}}
\renewcommand{\headrulewidth}{0.4pt}
\renewcommand{\footrulewidth}{0pt}

% --- Chapter & Section Title Styling ---
\titleformat{\chapter}[display]
  {\normalfont\bfseries\color{NavyPrimary}}
  {\filright\footnotesize\scshape\textcolor{CrimsonAccent}{\MakeUppercase{\chaptertitlename} \thechapter}}
  {4pt}
  {\titlerule[0.8pt]\vspace{6pt}\filright\LARGE}
  [\vspace{4pt}\titlerule]
\titlespacing*{\chapter}{0pt}{-15pt}{18pt}

\titleformat{\section}
  {\normalfont\large\bfseries\color{NavyPrimary}}
  {\textcolor{CrimsonAccent}{\thesection}}{0.6em}{}
\titleformat{\subsection}
  {\normalfont\normalsize\bfseries\color{SlateDark}}
  {\thesubsection}{0.5em}{}
\titleformat{\subsubsection}
  {\normalfont\normalsize\itshape\color{SlateDark}}
  {\thesubsubsection}{0.5em}{}

% --- Float placement: avoid near-empty float pages ---
\renewcommand{\topfraction}{0.9}
\renewcommand{\bottomfraction}{0.8}
\renewcommand{\textfraction}{0.07}
\renewcommand{\floatpagefraction}{0.7}
\setcounter{topnumber}{3}
\setcounter{totalnumber}{5}
% --- Start the main matter on the next page (not forced recto; digital edition) ---
\makeatletter
\renewcommand{\mainmatter}{\clearpage\@mainmattertrue\pagenumbering{arabic}}
\makeatother

\setcounter{tocdepth}{2}
\setcounter{secnumdepth}{3}

\begin{document}
\frontmatter

% ==============================================================================
% FULL-BLEED VECTOR FRONT COVER DESIGN (PDF PAGE 1)
% ==============================================================================
\thispagestyle{empty}
\noindent\null
\begin{tikzpicture}[remember picture, overlay]
  % Deep Navy Background
  \fill[NavyPrimary] (current page.south west) rectangle (current page.north east);

  % Top Crimson Banner Accent
  \fill[CrimsonAccent] (current page.north west) rectangle ($(current page.north east)+(0,-0.28in)$);
  \fill[GoldAccent] ($(current page.north west)+(0,-0.28in)$) rectangle ($(current page.north east)+(0,-0.32in)$);

  % Left Vertical Spine Accent Bar
  \fill[SlateDark] ($(current page.south west)+(0.45in,0)$) rectangle ($(current page.north west)+(0.52in,-0.32in)$);
  \fill[GoldAccent] ($(current page.south west)+(0.52in,0)$) rectangle ($(current page.north west)+(0.54in,-0.32in)$);

  % Decorative Polycentric Federated & Distributed Mesh Constellation (Centered in middle band)
  \begin{scope}[shift={($(current page.south west)+(3.48in,4.16in)$)}]
    % Outer Ecological & Planetary Commons Boundary
    \draw[SlateDark!85!white, line width=1.1pt, dashed] (0,0) circle (1.38in);

    % Background Peer-to-Peer Geodesic Chords (No Central Hub — Horizontal Polycentric Mesh)
    \foreach \ang in {0, 60, 120, 180, 240, 300} {
      % Cross-federation horizontal peer chords (connecting non-adjacent regional/industry federations)
      \draw[GoldAccent!55!SlateDark, line width=0.85pt, opacity=0.75] (\ang:0.74in) -- (\ang+120:0.74in);
      % Perimeter ring connecting the 6 regional/industry federation hubs
      \draw[GoldAccent!85!white, line width=1.15pt, opacity=0.9] (\ang:0.74in) -- (\ang+60:0.74in);
      % Diagonal cross-ties between regional hubs and neighbouring city/sector nodes
      \draw[TealAccent!65!white, thin, opacity=0.6] (\ang:0.74in) -- (\ang+30:1.06in);
      \draw[TealAccent!65!white, thin, opacity=0.6] (\ang:0.74in) -- (\ang-30:1.06in);
      % Links from each federation hub out to its local council cluster
      \draw[TealAccent!80!white, line width=0.8pt, opacity=0.8] (\ang:0.74in) -- (\ang-18:1.30in);
      \draw[TealAccent!80!white, line width=0.8pt, opacity=0.8] (\ang:0.74in) -- (\ang:1.34in);
      \draw[TealAccent!80!white, line width=0.8pt, opacity=0.8] (\ang:0.74in) -- (\ang+18:1.30in);
      % Intra-cluster horizontal council mesh links
      \draw[white!70!TealAccent, very thin, opacity=0.65] (\ang-18:1.30in) -- (\ang:1.34in) -- (\ang+18:1.30in);
    }

    % Secondary Inter-Cluster Mesh Ring (12 City / Sector Federation Relays)
    \foreach \ang in {30, 90, 150, 210, 270, 330} {
      % Cross-ties between city/sector nodes
      \draw[SlateDark!65!TealAccent, thin, opacity=0.55] (\ang:1.06in) -- (\ang+60:1.06in);
      % Links from city/sector relays to adjacent local council leaves
      \draw[GoldAccent!60!white, very thin, opacity=0.55] (\ang:1.06in) -- (\ang-12:1.30in);
      \draw[GoldAccent!60!white, very thin, opacity=0.55] (\ang:1.06in) -- (\ang+12:1.30in);
      \fill[TealAccent] (\ang:1.06in) circle (2.6pt);
      \draw[white, very thin] (\ang:1.06in) circle (2.6pt);
    }

    % Outer Local Worker & Consumer Council Nodes (18 Polycentric Leaf Councils)
    \foreach \ang in {0, 60, 120, 180, 240, 300} {
      \fill[white!92!GoldAccent] (\ang-18:1.30in) circle (1.9pt);
      \fill[TealAccent!85!white] (\ang:1.34in) circle (2.2pt);
      \draw[white, very thin] (\ang:1.34in) circle (2.2pt);
      \fill[white!92!GoldAccent] (\ang+18:1.30in) circle (1.9pt);
    }

    % 6 Polycentric Regional & Industry Federation Hubs (No Single Master Center)
    \foreach \ang in {0, 60, 120, 180, 240, 300} {
      \fill[CrimsonAccent] (\ang:0.74in) circle (3.5pt);
      \draw[GoldAccent, line width=0.7pt] (\ang:0.74in) circle (3.5pt);
    }

    % Open Shared Protocol Invariant Halo (Transparent Center — Shared Commons, Not a Central Office)
    \fill[NavyPrimary, opacity=0.82] (0,0) circle (0.36in);
    \draw[GoldAccent!75!white, line width=0.9pt, densely dashed] (0,0) circle (0.38in);
    \node[text=GoldAccent, font=\scriptsize\bfseries, align=center] at (0,0) {Two-Root\\[-1.5pt]\tiny$SB/SC \ge 1$};
  \end{scope}

  % Top Classification Header & Working Draft Badge
  \node[anchor=north west, text=GoldAccent, font=\footnotesize\scshape\bfseries]
    at ($(current page.north west)+(0.8in,-0.65in)$)
    {Post-Capitalist Labs Imprint $\cdot$ Working Draft Monograph};

  % Main Book Title Block
  \node[anchor=north west, text=white, font=\fontsize{25}{29}\selectfont\bfseries, align=left, text width=5.1in]
    at ($(current page.north west)+(0.8in,-1.05in)$)
    {DISTRIBUTED\\PARTICIPATORY\\ECONOMICS\\PROTOCOL};

  % Gold Rule Under Title
  \draw[GoldAccent, line width=1.8pt]
    ($(current page.north west)+(0.8in,-2.85in)$) -- ($(current page.north west)+(4.2in,-2.85in)$);

  % Subtitle in Opaque Framed Callout Box (zero background bleed-through)
  \node[anchor=north west, fill=NavyPrimary!95!black, draw=GoldAccent!60!SlateDark, line width=0.7pt, rounded corners=3pt, inner sep=7pt, text=white!92!GoldAccent, font=\small\itshape, align=left, text width=4.86in]
    at ($(current page.north west)+(0.8in,-2.98in)$)
    {A Rigorous, Source-Grounded Feasibility Study and Post-Quantum Technical Architecture for Democratic Planning, Self-Management, and the Productive Commons};

  % Audience Badges near Bottom
  \node[anchor=south west, fill=SlateDark!80!black, draw=GoldAccent!70!black, rounded corners=3pt, inner sep=8pt, text=white, text width=4.9in]
    at ($(current page.south west)+(0.8in,1.25in)$)
    {\scriptsize\textbf{\textcolor{GoldAccent}{DESIGNED FOR THREE INTERDISCIPLINARY AUDIENCES:}}\\[3pt]
     $\bullet$ \textbf{Advocates \& Cooperative Practitioners:} Council democracy, balanced job complexes, and historical lessons from Catalonia (1936), the Paris Commune (1871), and Classical Athens.\\[1pt]
     $\bullet$ \textbf{Economists \& Mechanism Designers:} Formal welfare theorems, Sandstr\"om's double-entry accounting, pollution demand revelation (CAPs), and candidate ADMM regularisation under convex restrictions.\\[1pt]
     $\bullet$ \textbf{Systems Programmers \& Cryptographers:} Local-first CRDTs, Two-Root BLAKE3 aggregation, deterministic WASM IFB engines, and NIST FIPS 203/204/205 post-quantum primitives.};

  % Footer Attribution (Post-Capitalist Labs Imprint)
  \node[anchor=south west, text=white!90!black, text width=5.06in]
    at ($(current page.south west)+(0.8in,0.52in)$)
    {{\fontsize{6.6}{8}\selectfont\hbox to 5.06in{\textbf{\textcolor{GoldAccent}{POST-CAPITALIST LABS IMPRINT}} $\cdot$ \texttt{https://github.com/Post-Capitalist-Labs} \hfill \textcolor{GoldAccent}{\textbf{Working Draft $\cdot$ Nov 2026}}}}\\[2.5pt]
     {\fontsize{5.4}{7}\selectfont\color{white!75!black} Leveraging technology to amplify social movements $\cdot$ Decentralised consensus, mechanism design \& democratic coordination}};
\end{tikzpicture}
\clearpage

% ==============================================================================
% FRONT MATTER (PDF Page 2: Imprint Colophon; Page 3: Reader's Guide; Page 4: TOC)
% ==============================================================================

% --- PDF PAGE 2: TITLE-VERSO / PUBLISHER IMPRINT & COLOPHON PAGE ---
\thispagestyle{empty}
\noindent{\footnotesize\scshape\bfseries\textcolor{CrimsonAccent}{Post-Capitalist Labs Imprint $\cdot$ Open-Source Research Monograph Series}}\\[0.35em]
{\Large\bfseries\color{NavyPrimary} Distributed Participatory Economics Protocol (DPEP)}\\[0.25em]
{\normalsize\itshape\color{SlateDark} A Rigorous, Source-Grounded Feasibility Study and Post-Quantum Technical Architecture for Democratic Planning, Self-Management, and the Productive Commons}\\[0.6em]
\rule{\textwidth}{0.6pt}\\[0.8em]

\noindent\colorbox{SoftBg}{%
  \parbox{\dimexpr\textwidth-2\fboxsep\relax}{%
    \vspace{6pt}
    {\normalsize\bfseries\color{NavyPrimary} About Post-Capitalist Labs (PCL)}\par\vspace{4pt}
    {\small\color{SlateDark}
    At \textbf{Post-Capitalist Labs}, we harness the transformative power of technology to bolster social movements and build directly democratic alternatives to capitalism. Our work spans decentralised consensus, applied cryptography, operations research, and democratic mechanism design.\par\vspace{4pt}
    \textbf{Our Vision:} We envisage a world where solidarity, care, classlessness, diversity, self-reflection, and direct democracy are not just ideals, but the foundations of everyday life. Our goal is to make this vision a reality through the thoughtful application of open-source technology, where people remain in control.\par\vspace{4pt}
    \textbf{Our Approach:}
    \begin{itemize}[leftmargin=1.5em,itemsep=2pt,topsep=2pt]
      \item \textbf{Learning from History:} Analyzing past and present democratic movements and self-managed collectives to guide institutional design.
      \item \textbf{Emancipatory Modeling:} Developing rigorous mathematical and computational models for equitable, non-market resource coordination.
      \item \textbf{Open-Source Engineering:} Turning theoretical models into verifiable, accessible, community-driven proofs of concept and protocols.
    \end{itemize}
    \vspace{2pt}
    \textbf{Organization Profile:} \url{https://github.com/Post-Capitalist-Labs}\\
    \textbf{Repository:} \url{https://github.com/Post-Capitalist-Labs/distributed-participatory-economics-protocol}}
    \vspace{6pt}
  }%
}

\vfill
\noindent\rule{\textwidth}{0.4pt}\\[0.4em]
{\scriptsize\color{SlateDark}
\textbf{Publication Metadata:}\\
\textbf{Imprint:} Post-Capitalist Labs Imprint (\url{https://github.com/Post-Capitalist-Labs})\\
\textbf{Status:} \textbf{Working Draft Monograph (November 2026)} --- Released for community feedback and peer technical review; separates demonstrated results, analytical estimates and open hypotheses ahead of gated Phase 1--3 experiments.\\[0.35em]
\textbf{Source Access \& Copyright Notice:} Sources are classified by how they were accessed;\footnote{Full list: \path{research/SOURCE_AUDIT.md} in the project repository.} this is not a claim that every page was re-read. \textbf{Tier 1 (full text or code):} Albert \& Hahnel (1991a, 1991b), Albert (2003), Dolgoff (1974), Szczepanczyk (2023), all 70 \texttt{pequod-cljs} generators and 150 CSV logs @ \texttt{8bd46b5c}, the Winter 2024 \emph{EJPE} issue, Deng et al.\ (2017), VDAF draft~22, the \texttt{prio} 0.18.1 crate (executed), the BLAKE3 specification, and Leval (1975, Ch.~V). \textbf{Tier 2 (preview, abstract or summary):} Hahnel (2021, 2022), Sandstr\"om (2020, 2025), Laibman (2014). \textbf{Tier 3 (standards or bibliographic record only):} NIST FIPS 202--205 and the remaining historical, economic and political works. \textbf{Tier 4:} \texttt{pequod-cljs} audited by script, not run. Short quotations and epigraphs are reproduced with attribution under fair dealing / fair use.}
\clearpage

% --- PDF PAGE 3: INTERDISCIPLINARY READER'S GUIDE ---
\thispagestyle{empty}
\noindent{\footnotesize\scshape\bfseries\textcolor{CrimsonAccent}{Interdisciplinary Navigation Guide}}\\[0.3em]
{\Large\bfseries\color{NavyPrimary} How to Read This Book Across Disciplines}\\[0.4em]
\rule{\textwidth}{0.6pt}\\[0.6em]
{\small\color{SlateDark} This monograph bridges libertarian-socialist institutional theory, formal welfare and planning economics, and modern distributed systems cryptography. Depending on your primary background, we recommend the following reading paths:}\\[0.6em]

\audiencebox{For Advocates, Cooperators \& Organizers}{Institutional Self-Management Without Hierarchy}{
Start with \textbf{Chapter 1 (Executive Summary)} and \textbf{Chapter 2 (Historical \& Institutional Inspirations)} to see how lessons from the 1936--1939 Spanish Revolution, the 1871 Paris Commune, Classical Athens, Occupy Wall Street, Mondrag\'on, Porto Alegre Participatory Budgeting, Rojava, and Elinor Ostrom's commons principles are translated into concrete protocol safeguards against bureaucratic capture, meeting fatigue, and inter-regional inequality. Then read \textbf{Chapter 3} for the complete traceability matrix linking Michael Albert and Robin Hahnel's 1991 and 2021 books to our design.}

\vspace{4pt}
\audiencebox{For Economists \& Operations Researchers}{Formal Models, Accounting \& Non-Convex Convergence}{
Focus on \textbf{Chapter 4 (Formal Economic Model)}, which unites Albert and Hahnel's Formal Model of a Participatory Economy (FMPE) with Anders Sandstr\"om's double-entry \emph{Society Account} and Robin Hahnel's Pollution Demand-Revealing Mechanism (PDRM). See \textbf{Section 4.4} for our source-code audit of \texttt{pequod-cljs} and the four research-candidate solvers (Algorithms 3a--3d), alongside \textbf{Chapter 8} and \textbf{Chapter 10} for engagement with John O'Neill (2024) and the non-convex experiments still required.}

\vspace{4pt}
\audiencebox{For Systems Programmers \& Cryptographers}{Local-First P2P, Two-Root Aggregation \& Post-Quantum Security}{
Jump directly to \textbf{Chapter 4 (Section 4.5 State Machine \& Canonical CBOR Schemas)}, \textbf{Chapter 5 (Four-Layer Architecture D \& 6-Role Mobile Taxonomy)}, \textbf{Chapter 6 (Weighted Stack Comparison)}, \textbf{Chapter 7 (NIST FIPS 203/204/205, Sybil Resistance \& \texttt{Prio3} VDAF Aggregation)}, and \textbf{Chapter 9 (Experimental Programme, Kill Gates \& Multi-Scale Estimates)}.}
\clearpage

\tableofcontents
\clearpage

% ==============================================================================
% FRONT MATTER: LIST OF ABBREVIATIONS AND NOTATION
% ==============================================================================
\chapter*{List of Abbreviations and Mathematical Notation}
\addcontentsline{toc}{chapter}{List of Abbreviations and Mathematical Notation}
\markboth{List of Abbreviations and Notation}{List of Abbreviations and Notation}

{\small\color{SlateDark} Because this monograph integrates participatory economics, labor and cooperative history, operations research, distributed systems, and post-quantum cryptography, this reference table defines every abbreviation and mathematical symbol used across Chapters 1--10.}

\subsection*{A. Participatory Economics, Accounting, and Bibliographic Abbreviations}
{\scriptsize\setlength{\tabcolsep}{4pt}\renewcommand{\arraystretch}{1.18}
\begin{longtable}{L{0.92in} L{1.75in} L{1.95in}}
\toprule
\textbf{Abbreviation} & \textbf{Full Term} & \textbf{Definition \& Role in Monograph} \\
\midrule
\endfirsthead
\multicolumn{3}{l}{\scriptsize\itshape Continued from previous page} \\
\toprule
\textbf{Abbreviation} & \textbf{Full Term} & \textbf{Definition \& Role in Monograph} \\
\midrule
\endhead
\bottomrule
\endlastfoot
\textbf{BJC} & Balanced Job Complex & Bundle of workplace tasks balanced for empowerment and desirability within and across workplaces (PEPE Ch.~2). \\
\textbf{CAP} & Community of Affected Parties & Ecological/geographic collective of residents impacted by a specific pollutant in an airshed or watershed (DEP Ch.~7). \\
\textbf{CC} ($CC_h$) & Consumers' Council & Democratic neighbourhood consumer council or nested ward/city/regional consumer federation. \\
\textbf{CFB} & Consumption Facilitation Board & Informational board assisting consumer councils and federations in evaluating collective consumption projects (LF Ch.~6). \\
\textbf{CLT} & Community Land Trust & Member-controlled trust holding land or long-lived assets outside market sale (Chapter~11). \\
\textbf{CLUEA} & \emph{Consejo Levantino Unificado de Exportaci\'on Agr\'icola} & Unified Levantine Council for Agricultural Export organised by the CNT and UGT in Valencia (1936--1939). \\
\textbf{CNT} & \emph{Confederaci\'on Nacional del Trabajo} & Anarcho-syndicalist National Confederation of Labour in the 1936--1939 Spanish Revolution. \\
\textbf{CPR} & Common-Pool Resource & Shared natural or human-made resource system governed by collective-choice rules (Ostrom 1990). \\
\textbf{DAANES} & Democratic Autonomous Administration of North and East Syria & Rojava's confederal administration, discussed as a historical precedent (Chapter~2). \\
\textbf{DEP} & \emph{Democratic Economic Planning} & Primary reference volume by Robin Hahnel et al.\ (Routledge, 2021). \\
\textbf{DPEP} & Distributed Participatory Economics Protocol & The open-source decentralised coordination protocol specified in this monograph. \\
\textbf{EFB} & Employment Facilitation Board & Facilitation board assisting workers transferring between councils and balancing cross-workplace BJCs (LF Ch.~6). \\
\textbf{EJPE} & \emph{Erasmus Journal for Philosophy and Economics} & Academic journal hosting the Winter 2024 symposium on Hahnel's \emph{Democratic Economic Planning}. \\
\textbf{FMPE} & Formal Model of a Participatory Economy & Mathematical formalisation of participatory planning analyzed in PEPE Chapter 5 (pp.~73--106). \\
\textbf{FTE} & Full-Time Equivalent & Unit of engineering or research staffing effort across roadmap scenarios. \\
\textbf{GA} & General Assembly & Direct plenary deliberative assembly of a workplace, village collective, or horizontal movement. \\
\textbf{GDP} & Gross Domestic Product & Aggregate social value of net final output ($Y^{\text{GDP}}$) cleared in the Society Account. \\
\textbf{GDPR} & General Data Protection Regulation & EU data-protection law governing pilot data handling, retention and erasure. \\
\textbf{GHG} & Greenhouse Gas & Atmospheric emissions ($\text{CO}_2$-equivalent) subject to hard biophysical ceilings ($E_z^{\max}$). \\
\textbf{HFB} & Household Facilitation Board & Facilitation board assisting individuals and families with housing and neighbourhood membership (LF Ch.~6). \\
\textbf{IF} ($IF_k$) & Industry Federation & Federation of worker councils within industry $k$ managing sector capacity, R\&D, and $SB/SC$ appeals. \\
\textbf{IFB} & Iteration Facilitation Board & Informational board/engine that aggregates supply and demand and updates indicative prices via public rules. \\
\textbf{LF} & \emph{Looking Forward} & Primary reference volume by Michael Albert and Robin Hahnel (South End Press, 1991). \\
\textbf{Parecon / PE} & Participatory Economics & Post-capitalist economic model developed by Michael Albert and Robin Hahnel. \\
\textbf{PB} & Participatory Budgeting / Planning Bureau & Municipal democratic budgeting (Porto Alegre, 1989); also denotes the abstract Planning Bureau in FMPE Ch.~5. \\
\textbf{PCL} & Post-Capitalist Labs & Open-source research and engineering collective publishing this monograph (\url{https://github.com/Post-Capitalist-Labs}). \\
\textbf{PDRM} & Pollution Demand-Revealing Mechanism & Iterative procedure pricing local pollutant damages via CAP emission tolerances and credits (DEP Ch.~7). \\
\textbf{PEP} & Participatory Economy Project & Research collective maintaining \texttt{participatoryeconomy.org} and the 2023 User Story Mapping. \\
\textbf{PEPE} & \emph{The Political Economy of Participatory Economics} & Primary technical reference volume by Michael Albert and Robin Hahnel (Princeton University Press, 1991). \\
\textbf{PFB} & Production Facilitation Board & Board assisting worker councils with production externalities, scale economies, and technology choice (LF Ch.~6). \\
\textbf{PIPA} & Practically Identical Production Activities & Equivalence classes $\{PIPA\}_i$ partitioning workplace tasks by qualitative empowerment effect (PEPE Ch.~2, p.~26). \\
\textbf{PLAC} & \emph{Parecon: Life After Capitalism} & Comprehensive institutional and daily-life synthesis of Participatory Economics by Michael Albert (Verso Books, 2003). \\
\textbf{SB / SC} & Social Benefit / Social Cost & Ratio $SB_j(P)/SC_j(P)$ measuring the social value of a worker council's outputs relative to the opportunity cost of its inputs and endowments ($\ge 1$ constraint). \\
\textbf{SDF} & Syrian Democratic Forces & Military coalition of North and East Syria; party to the 2026 integration agreement (Chapter~2). \\
\textbf{SKU} & Stock Keeping Unit & Granular product variant (size, color, style) resolved at intra-year fulfilment from coarse planning categories. \\
\textbf{TFEU} & Treaty on the Functioning of the European Union & EU treaty whose Article~101 governs coordination between independent enterprises (Chapter~11). \\
\textbf{UFB} & Updating Facilitation Board & Facilitation board mediating intra-year plan adjustments when demand shifts or supply shocks exceed slack (LF Ch.~6). \\
\textbf{UGT} & \emph{Uni\'on General de Trabajadores} & Socialist General Union of Workers that co-organised industrial and agrarian collectives in Spain (1936--1939). \\
\textbf{WC} ($WC_j$) & Workers' Council & Democratic workplace collective where each worker has one vote and holds a Balanced Job Complex. \\
\end{longtable}}

\subsection*{B. Operations Research, AI, Distributed Systems, and Cryptography Abbreviations}
{\scriptsize\setlength{\tabcolsep}{4pt}\renewcommand{\arraystretch}{1.18}
\begin{longtable}{L{0.92in} L{1.75in} L{1.95in}}
\toprule
\textbf{Abbreviation} & \textbf{Full Term} & \textbf{Definition \& Role in Monograph} \\
\midrule
\endfirsthead
\multicolumn{3}{l}{\scriptsize\itshape Continued from previous page} \\
\toprule
\textbf{Abbreviation} & \textbf{Full Term} & \textbf{Definition \& Role in Monograph} \\
\midrule
\endhead
\bottomrule
\endlastfoot
\textbf{ADMM} & Alternating Direction Method of Multipliers & Augmented-Lagrangian primal-dual method; Algorithm~3b is the proximal Jacobian variant of Deng et al.\ (2017), a research candidate. \\
\textbf{AI} & Artificial Intelligence & Machine-learned models; restricted in DPEP to advisory and accessibility roles (Chapter~8). \\
\textbf{AP / CP} & Availability / Consistency under Partition & CAP theorem properties: local-first drafting is AP (available offline), while Epoch BFT price finality is CP. \\
\textbf{ARIMA / ETS} & AutoRegressive Integrated Moving Average / Error-Trend-Seasonality & Classical statistical time-series forecasting models used for baseline demand and SKU disaggregation. \\
\textbf{BBS} & Boneh--Boyen--Shacham signatures & Pairing-based signatures used for anonymous credentials; not post-quantum. \\
\textbf{BFT} & Byzantine Fault Tolerance & Distributed consensus property guaranteeing safety even when up to $f < n/3$ validator nodes act maliciously. \\
\textbf{BHT} & Brassard--H{\o}yer--Tapp algorithm & Quantum collision-finding algorithm; relevant to hash security margins. \\
\textbf{BLAKE3} & BLAKE3 Cryptographic Hash & Hash function with 256-bit output and a 128-bit security target; used for the Two-Root aggregation trees and content addressing. \\
\textbf{BLS} & Boneh--Lynn--Shacham signatures & Pairing-based aggregatable signatures; used in threshold randomness beacons; not post-quantum. \\
\textbf{BSP} & Bulk-Synchronous Parallel & Execution model alternating parallel local computation rounds with synchronised global aggregation checkpoints. \\
\textbf{CapCert} & Capability Certificate & Chained, signed certificates encoding council membership and revocable delegation. Format to be chosen (UCAN 1.0 and Biscuit are candidates); signed under the active \texttt{CryptoSuiteId} (Section~\ref{sec:pq-urgency}). \\
\textbf{CBOR (CDE)} & Concise Binary Object Representation & RFC 8949 Canonical Deterministic Encoding used for all signed DPEP messages with fixed $10^{-6}$ micro-unit scale. \\
\textbf{CFRG} & Crypto Forum Research Group & IRTF group that develops cryptographic specifications such as VDAF. \\
\textbf{CID} & Content Identifier & Self-certifying cryptographic hash address referencing a redactable qualitative \texttt{List} or plan blob. \\
\textbf{CPU / RAM} & Central Processing Unit / Random-Access Memory & Processor time and working memory, used in device and server cost estimates. \\
\textbf{CRDT} & Conflict-Free Replicated Data Type & Data structure (Loro / Automerge) allowing offline concurrent edits in council workspaces that merge deterministically. \\
\textbf{CRQC} & Cryptographically Relevant Quantum Computer & Future fault-tolerant quantum computer capable of breaking classical elliptic-curve and RSA cryptography via Shor's algorithm. \\
\textbf{CSV} & Comma-Separated Values & Plain-text table format of the \texttt{pequod-cljs} run logs. \\
\textbf{DA} & Data Availability & Verification that all leaf proposals of a round root are downloadable by auditing nodes (0.24--1.46~GB raw per round at $N=10^6$; Table~\ref{tab:proposal-sizes}). \\
\textbf{DAO} & Decentralised Autonomous Organisation & Blockchain governance construct critiqued in Chapter 2 for conflating token voting with council democracy. \\
\textbf{DAP} & Distributed Aggregation Protocol & IETF protocol for running VDAFs between clients and aggregators. \\
\textbf{DP} & Differential Privacy & Calibrated noise added to released totals to bound what can be learned about one person. \\
\textbf{DTN} & Delay-Tolerant Networking & Store-and-forward networking (including LoRa mesh and USB data-mules) for intermittently connected regions. \\
\textbf{EEIO / LCA} & Environmentally Extended Input-Output / Lifecycle Assessment & Matrix accounting methods estimating cradle-to-gate resource and emission footprints across supply chains. \\
\textbf{FIPS} & Federal Information Processing Standards & NIST post-quantum cryptographic standards finalised in August 2024: FIPS 203, FIPS 204, and FIPS 205. \\
\textbf{FLP} & Fully Linear Proof & VDAF-22 proof on secret shares; binds report validity, not aggregator honesty. \\
\textbf{FN-DSA} & FFT over NTRU-Lattice Digital Signature Algorithm & Falcon-based signature standard (FIPS~206, draft) with smaller signatures than ML-DSA. \\
\textbf{FOC} & First-Order Conditions & Optimality conditions from setting derivatives of the objective to zero. \\
\textbf{FRI / STARK} & Fast Reed-Solomon IOP / Scalable Transparent ARgument of Knowledge & Discarded household-budget design; not part of Architecture D. \\
\textbf{GC} & Garbage Collection & Automatic memory reclamation in managed language runtimes. \\
\textbf{HCI} & Human--Computer Interaction & Study of interface usability, used in the Phase~1H probe. \\
\textbf{HPKE} & Hybrid Public Key Encryption & IETF scheme (RFC~9180) used by DAP to encrypt report shares to aggregators. \\
\textbf{HQC} & Hamming Quasi-Cyclic & Code-based key-encapsulation mechanism selected by NIST in 2025 as a backup to ML-KEM. \\
\textbf{HRU} & Household Reporting Unit & Per-category integer unit for household \texttt{Prio3} reports, converted exactly to protocol micro-units (Section~\ref{sec:prio3-cost}). \\
\textbf{IDL} & Interface Definition Language & Language-neutral description of message schemas (Listings~4.1--4.2). \\
\textbf{IEEE-754} & IEEE Floating-Point Standard & Hardware floating-point representation replaced in DPEP core rules by deterministic 64-bit fixed-point integers. \\
\textbf{IETF / IRTF} & Internet Engineering / Research Task Force & Bodies that publish Internet standards (RFCs) and research drafts. \\
\textbf{IIS} & Irreducible Infeasible Subsystem & Minimal set of constraints that makes an optimisation problem infeasible. \\
\textbf{KEM} & Key-Encapsulation Mechanism & Public-key method for establishing a shared secret key (e.g.\ ML-KEM). \\
\textbf{KKT} & Karush-Kuhn-Tucker Conditions & First-order necessary and sufficient optimality conditions for constrained nonlinear/conic optimisation. \\
\textbf{LAN / WAN} & Local / Wide Area Network & Council-premises network versus links between councils and federations. \\
\textbf{LFS} & Large File Storage & Git extension storing large data files outside the repository history. \\
\textbf{LLM} & Large Language Model & Generative text model; restricted in DPEP to translation, accessibility and summaries. \\
\textbf{LP / QP / MILP} & Linear / Quadratic / Mixed-Integer Linear Programming & Mathematical optimisation paradigms solved locally via \texttt{Clarabel.rs} and \texttt{HiGHS}. \\
\textbf{LWE / SIS} & Learning With Errors / Short Integer Solution & Lattice problems whose module variants underlie ML-KEM and ML-DSA. \\
\textbf{M0--M5} & Maturity levels & Component maturity scale from unspecified (M0) to ubiquitous (M5) (Section~\ref{sec:maturity}). \\
\textbf{MAPE} & Mean Absolute Percentage Error & Accuracy metric used to audit and backtest statistical demand forecasts. \\
\textbf{MCP} & Mixed Complementarity Problem & Joint formulation of optimality and market-clearing conditions solved by PATH (Algorithm~3d). \\
\textbf{ML-DSA-65} & Module-Lattice Digital Signature Algorithm & NIST FIPS 204 (Dilithium3) post-quantum digital signature scheme combined in hybrid mode with Ed25519. \\
\textbf{ML-KEM-768} & Module-Lattice Key-Encapsulation Mechanism & NIST FIPS 203 (Kyber-768) post-quantum key-establishment scheme combined in hybrid mode with X25519. \\
\textbf{MLS} & Messaging Layer Security & IETF RFC 9420 protocol providing end-to-end encrypted group messaging inside councils. \\
\textbf{MPC} & Secure Multiparty Computation & Candidate hardening of aggregator correctness (Table~\ref{tab:agg-options}); not the Phase-1 construction. \\
\textbf{NASA-TLX} & NASA Task Load Index & Standardised workload questionnaire proposed for measuring participation fatigue in Phase 1H and Phase 3. \\
\textbf{NIST} & National Institute of Standards and Technology & US standards body that published FIPS 203--205 (post-quantum cryptography). \\
\textbf{NMT / TTS} & Neural Machine Translation / Text-to-Speech & Local accessibility models translating qualitative \texttt{Lists} across languages and assisting visually impaired users. \\
\textbf{OR} & Operations Research & Discipline of mathematical optimisation, linear programming, and supply-chain clearing. \\
\textbf{P2P} & Peer-to-Peer & Decentralised networking architecture without central client-server gatekeepers. \\
\textbf{PoC} & Proof of Concept & A proposed, unbuilt reference implementation (Phases 1--3, Chapter~9); no implementation exists. \\
\textbf{PQ} & Post-Quantum & Prefix for schemes designed to resist quantum attacks (see PQC). \\
\textbf{PQC} & Post-Quantum Cryptography & Cryptographic algorithms resistant to both classical and quantum (Shor/Grover) attacks. \\
\textbf{PRG} & Pseudorandom Generator & Deterministic expansion of a short random seed into a long pseudorandom string. \\
\textbf{Prio3 / VDAF} & Verifiable Distributed Aggregation Function & IETF CFRG \texttt{draft-irtf-cfrg-vdaf-22} (Informational). Privacy needs one honest aggregator; correct totals need all of them. \\
\textbf{PWA} & Progressive Web Application & Offline-capable cross-platform web client running the Rust/WASM \texttt{Role-LP} engine on smartphones and desktops. \\
\textbf{QUIC} & Quick UDP Internet Connections & Modern encrypted, multiplexed transport protocol (\texttt{quinn} / \texttt{iroh}) used for Layer 2 federation sync. \\
\textbf{RFC} & Request for Comments & IETF document series that includes Internet standards. \\
\textbf{Role-LP / CC / SR / PCW / BFT / GW} & DPEP Node Roles 1--6 & Lightweight Participant (\texttt{LP}), Council Coordinator (\texttt{CC}), Storage Replica (\texttt{SR}), Planning Compute Worker (\texttt{PCW}), Epoch Validator (\texttt{BFT}), and Delay-Tolerant Gateway (\texttt{GW}). \\
\textbf{SDK} & Software Development Kit & Libraries and tools for building on a platform (e.g.\ Cosmos SDK). \\
\textbf{SHA-2 / SHA-3} & Secure Hash Algorithm 2 / 3 & NIST hash-function families (e.g.\ SHA-256, SHA3-256; FIPS 180-4, 202). \\
\textbf{SLH-DSA} & Stateless Hash-Based Digital Signature Algorithm & NIST FIPS 205 (SPHINCS+) signature standard relying solely on hash functions, used for long-term root anchors. \\
\textbf{SQL / REST} & Structured Query Language / Representational State Transfer & Relational database queries and web APIs of the conventional Architecture~A. \\
\textbf{TLS} & Transport Layer Security & Standard protocol for encrypted network channels; used inside QUIC. \\
\textbf{TPS} & Transactions Per Second & Throughput metric shown in Chapter 5 to be irrelevant to epoch-synchronised participatory planning. \\
\textbf{TTL} & Time-To-Live & Hard expiration timer on delegate capability certificates and emergency replanning playbooks. \\
\textbf{Two-Root} & Two-Root BLAKE3 aggregation & Proposal Merkle root plus category-totals root (Section~\ref{sec:two-root}). Replaces internal-vector Sparse Merkle-Sum Trees. \\
\textbf{UCAN} & User Controlled Authorization Network & Capability-token format; a candidate format for CapCerts. \\
\textbf{UI / UX} & User Interface / User Experience & Screens and interaction design of participant clients. \\
\textbf{UniFFI} & Universal Foreign Function Interface & Mozilla binding generator exposing the Rust \texttt{dpep-core} library to native Kotlin (Android) and Swift (iOS). \\
\textbf{VRF / VDF} & Verifiable Random / Delay Function & Primitives for unbiasable public randomness; candidate post-quantum beacon building blocks. \\
\textbf{WASM} & WebAssembly & Deterministic, sandboxed binary instruction format (\texttt{Wasmtime}) executing IFB price-update rules. \\
\textbf{X25519 / Ed25519} & Curve25519 key exchange / signatures & Classical elliptic-curve halves of DPEP's hybrid schemes. \\
\textbf{XOF} & Extendable-Output Function & Hash function with output of any requested length (e.g.\ TurboSHAKE128 in VDAF). \\
\end{longtable}}

\subsection*{C. Summary of Core Mathematical Notation (Chapter 4)}
{\scriptsize\setlength{\tabcolsep}{4pt}\renewcommand{\arraystretch}{1.18}
\begin{longtable}{L{0.92in} L{3.74in}}
\toprule
\textbf{Symbol} & \textbf{Mathematical Meaning} \\
\midrule
$\Omega$ & Master catalogue of all $G = |\Omega|$ priced economic categories: $\mathcal{G}^{\text{priv}} \sqcup \mathcal{G}^{\text{pub}} \sqcup \mathcal{G}^{\text{int}} \sqcup \mathcal{R} \sqcup \mathcal{L} \sqcup \mathcal{K} \sqcup \mathcal{Z}$. \\
$P^{(t)} \in \mathbb{R}_+^{|\Omega|}$ & Vector of indicative social opportunity-cost prices at planning iteration $t$. \\
$\bar{a}_j, \underline{a}_j, r_j, e_j$ & Worker council $WC_j$'s proposed outputs ($\bar{a}_j$), intermediate inputs ($\underline{a}_j$), natural resources ($r_j$), and emissions ($e_j$). \\
$l_j^0, k_j^0, \Delta l_j, \Delta k_j$ & Initial labour ($l_j^0$) and capital ($k_j^0$) endowments of $WC_j$, and proposed net additions/releases ($\Delta l_j, \Delta k_j$). \\
$r^*, l^*, k^*, E_z^{\max}$ & Economy-wide available natural resources ($r^*$), labour supplies ($l^*$), capital stocks ($k^*$), and hard biophysical emission caps ($E_z^{\max}$). \\
$w_{\text{base}}, \epsilon_u, B_u(P)$ & Base hourly remuneration rate ($w_{\text{base}}$), worker effort multiplier ($\epsilon_u$), and citizen $u$'s warranted annual consumption budget ($B_u$). \\
$S_i^{(t)}, D_i^{(t)}, SL_i^{(t)}$ & Aggregate supply ($S_i^{(t)}$), aggregate demand ($D_i^{(t)}$), and planned slack buffer ($SL_i^{(t)} = \eta_i S_i^{(t)}$) for category $i$. \\
$\rho_i^{(t)}, \tau$ & Normalized relative excess demand imbalance $\rho_i^{(t)} = \frac{2|S_i^{(t)} - D_i^{(t)}|}{S_i^{(t)} + D_i^{(t)}}$ and two-sided clearing tolerance $\tau \in [0.01, 0.05]$ (Definition~\ref{def:success}). \\
\bottomrule
\end{longtable}}
\clearpage

\mainmatter

% ==============================================================================
% CHAPTER 1: EXECUTIVE SUMMARY
% ==============================================================================
\chapter{Executive Summary: Feasibility, Architecture, and Strategic Roadmap}

\chapterepigraph{No central agency composes the final plan. No market competition generates the final plan. A decentralised, social, iterative, communication process allows all actors to democratically formulate the final plan.}{Michael Albert and Robin Hahnel, \emph{Looking Forward: Participatory Economics for the Twenty First Century} (1991), Ch.~6\protect\footnotemark}
\footnotetext{Michael Albert and Robin Hahnel, \emph{Looking Forward: Participatory Economics for the Twenty First Century} (Boston, MA: South End Press, 1991), Ch.~6 (``Participatory Allocation''), accessed via ZNetwork authorized online edition (\url{https://znetwork.org/wp-content/uploads/zbooks/htdocs/ParEcon/lookingforward/lf6_1.htm}).}

\section{What Is Being Proposed}

This monograph sets out a source-grounded feasibility assessment, a gated research programme, and a technical architecture proposal for the \textbf{Distributed Participatory Economics Protocol (DPEP)}: an open-source, local-first, cryptographically verifiable coordination protocol intended to enable democratic worker councils, neighbourhood consumer councils, and nested federations to formulate, iterate, ratify, and execute comprehensive economic plans without capitalist markets or authoritarian central planning.

The intellectual foundation of this protocol is \textbf{Participatory Economics (Parecon)}, originally developed by Michael Albert and Robin Hahnel in \emph{The Political Economy of Participatory Economics} (1991) and \emph{Looking Forward} (1991),\footnote{Michael Albert and Robin Hahnel, \emph{The Political Economy of Participatory Economics} (Princeton, NJ: Princeton University Press, 1991), \url{https://doi.org/10.2307/j.ctv131bvr2}; Albert and Hahnel, \emph{Looking Forward} (1991).} elaborated across daily-life institutions and critic rebuttals in Michael Albert's \emph{Parecon: Life After Capitalism} (2003),\footnote{Michael Albert, \emph{Parecon: Life After Capitalism} (London \& New York: Verso Books, 2003), full text archived at ZCommunications/ZNetwork (\url{http://books.zcommunications.org/books/pareconv/}).} and mathematically extended in Robin Hahnel's \emph{Democratic Economic Planning} (2021),\footnote{Robin Hahnel (with contributions by Peter Bohmer, Savvina Chowdhury, Allison Kerkhoff, Mitchell Szczepanczyk, and Michael Weisdorf), \emph{Democratic Economic Planning} (Abingdon \& New York: Routledge, 2021), \url{https://doi.org/10.4324/9781003173700}.} Anders Sandstr\"om's \emph{Anarchist Accounting} (2020),\footnote{Anders Sandstr\"om, \emph{Anarchist Accounting: Accounting Principles for a Democratic Economy} (Abingdon \& New York: Routledge, 2020), \url{https://doi.org/10.4324/9781003036005}; see also Sandstr\"om, ``An Accounting System for a Libertarian Socialist Economy,'' Participatory Economy Project (Feb. 14, 2023), \url{https://participatoryeconomy.org/an-accounting-system-for-a-libertarian-socialist-economy/}.} and Mitchell Szczepanczyk's computer simulation work (\texttt{pequod-cljs}, 2023).\footnote{Mitchell Szczepanczyk, ``Pseudocode and Algorithms for Computer Simulations of Democratically Planned Economies,'' \emph{Journal of Information Economics} 1, no. 3 (2023): 43--54, \url{https://doi.org/10.58567/jie01030004}; open-source repository at \url{https://github.com/msszczep/pequod-cljs}.}

DPEP provides the computational substrate for four core institutional pillars:
\begin{enumerate}[leftmargin=1.8em,itemsep=2pt]
  \item \textbf{Self-Managed Work and Consumption:} Democratic worker councils ($WC_j$) and consumer councils ($CC_h$) in which participants exercise decision-making input proportional to the degree they are affected by an economic choice.\footnote{Albert and Hahnel, \emph{Political Economy of Participatory Economics}, Ch.~1, pp.~8, 16; Ch.~2, p.~24; Hahnel, \emph{Democratic Economic Planning}, Ch.~1, pp.~14--16.}
  \item \textbf{Equitable Remuneration and Balanced Job Complexes (BJCs):} Non-transferable consumption allowances calibrated to the duration, intensity, and onerousness (effort and sacrifice) of socially valued labour, supplemented by universal free public goods and council-approved special-need allowances, paired with job complexes balanced for empowerment and desirability within and across workplaces.\footnote{Albert and Hahnel, \emph{Political Economy of Participatory Economics}, Ch.~2, pp.~25--30; Ch.~3, pp.~44--55.}
  \item \textbf{Decentralised Iterative Planning:} A multi-round coordination protocol in which councils submit production and consumption proposals, federated aggregators compute excess demands across seven economic categories (private goods, public goods, intermediate inputs, natural resources, labour categories, capital stocks, and pollutant emissions), and a deterministic Iteration Facilitation Board (IFB) engine updates indicative social opportunity-cost prices until a feasible, ecologically bounded plan is ratified.\footnote{Albert and Hahnel, \emph{Political Economy of Participatory Economics}, Ch.~4, pp.~57--72; Ch.~5, pp.~73--106; Hahnel, \emph{Democratic Economic Planning}, Ch.~7, pp.~130--162.}
  \item \textbf{Intra-Year Operational Clearing and Scoped Replanning:} Continuous tracking of deliveries and inventory drawdowns against the ratified plan, absorbing routine consumer preference shifts locally within neighbourhood and city slack buffers while mediating larger supply shocks across affected industry federations via Updating Facilitation Boards (UFBs).\footnote{Albert and Hahnel, \emph{Political Economy of Participatory Economics}, Ch.~4, pp.~66--67; \emph{Looking Forward}, Ch.~6 (\texttt{lf6\_1.htm}); Hahnel, \emph{Democratic Economic Planning}, Ch.~8, pp.~167--170.}
\end{enumerate}

\section{Why It Matters: Breaking the ``No Third Way'' Dogma}

For over a century, mainstream and socialist economists alike have been haunted by Alec Nove's famous assertion in \emph{The Economics of Feasible Socialism} (1983):
\begin{quote}
\itshape\small ``In a complex industrial economy the interrelation between its parts can be based in principle either on freely chosen negotiated contracts [i.e., markets], or on a system of binding instructions from planning offices [i.e., command planning]. There is no third way.''\footnote{Alec Nove, \emph{The Economics of Feasible Socialism} (London: George Allen \& Unwin, 1983), pp.~ix--x, 44; quoted directly at the opening of Chapter 4 in Albert and Hahnel, \emph{Political Economy of Participatory Economics}, p.~57.}
\end{quote}

Yet historical experiments in democratic self-management---most notably the industrial and agrarian collectives of the \textbf{Spanish Revolution (1936--1939)}, the \textbf{Paris Commune (1871)}, federated cooperatives like \textbf{Mondrag\'on}, and horizontal movements like \textbf{Occupy Wall Street (2011)}---repeatedly encountered two structural barriers when attempting to coordinate beyond single workplaces or rural villages:
\begin{itemize}[leftmargin=1.6em,itemsep=2pt]
  \item \textbf{The Accounting and Valuation Barrier:} Without a unified social accounting and indicative-price system, wealthier collectives hoarded superior inherited machinery (what Gaston Leval and Frank Mintz termed ``collective egoism'' or \emph{neocapitalismo colectivo} in 1936 Barcelona),\footnote{Gaston Leval, \emph{Collectives in the Spanish Revolution}, trans. Vernon Richards (London: Freedom Press, 1975); Frank Mintz, \emph{Anarchism and Workers' Self-Management in Revolutionary Spain}, trans. Paul Sharkey (Oakland, CA: AK Press, 2013); Sam Dolgoff, ed., \emph{The Anarchist Collectives: Workers' Self-Management in the Spanish Revolution, 1936--1939} (New York: Free Life Editions, 1974), Ch.~6.} while inter-industry coordination of complex intermediate inputs broke down into cumbersome bilateral barter.
  \item \textbf{The Deliberative Scaling Barrier:} Without structured, asynchronous computational protocols, horizontal deliberation tended towards exhausting multi-hour plenary meetings (speaking time grows at least linearly with attendance, and possible pairwise exchanges quadratically; this is arithmetic, not an empirical threshold, Section~2.5) or unaccountable informal cliques (Jo Freeman's ``tyranny of structurelessness'').\footnote{Jo Freeman, ``The Tyranny of Structurelessness,'' \emph{The Second Wave} 2, no. 1 (1972) / \emph{Berkeley Journal of Sociology} 17 (1972--1973): 151--164; David Graeber, \emph{The Democracy Project: A History, a Crisis, a Movement} (New York: Spiegel \& Grau, 2013).}
\end{itemize}

Contemporary distributed computing, post-quantum cryptography, and mathematical optimisation make it \emph{conceivable} to attack both barriers at once. Whether millions of participants can coordinate industrial supply chains and ecological boundaries democratically, and at what wall-clock cost per round, is an experimental claim: Category B estimates in Table~\ref{tab:exec-feasibility} and Table~\ref{tab:scale-bench} are schema-derived, not measured at that scale.

\section{Technical Feasibility: Demonstrated Results, Estimates, and Open Targets}

A core mandate of this monograph is to separate what has been \textbf{shown}, what has been \textbf{estimated}, and what is \textbf{unknown}.

\paragraph{What is established, estimated and unknown (plain statement).}
\emph{Established:} the 1991 model and its later formalisations exist and are internally coherent at the level of accounting identities; the \texttt{pequod-cljs} Cobb--Douglas simulations reach a 5\% imbalance in about 12 rounds (fixed incomes, smooth convex technology); the unmodified price rule fails on Scarf's cycle and on a single indivisible plant; the NIST post-quantum primitives are standardised, and every function in the stack except threshold signatures, anonymous credentials, the sortition beacon, capability-certificate formats and post-quantum HPKE has a standardised, deployed primitive (Table~\ref{tab:maturity}), although almost every way DPEP combines them is original; and \texttt{Prio3SumVec} report sizes and desktop timings have been measured with \texttt{libprio}. \emph{Estimated (not measured at scale):} round latency, archive sizes, mobile bandwidth, hosting cost, and aggregator CPU at $10^6$ participants. \emph{Unknown:} whether any candidate solver clears economies with Leontief inputs, lumpy plants, endogenous incomes and strategic reporting; whether ordinary people can and will take part in multi-round planning without a technical elite forming around the software; how household budgets can be checked privately when prices change between rounds; and whether a credential scheme can enforce one report per person without a population register. Phase~1 and Phase~1H (Chapter~9) are designed to settle the first two unknowns, or to stop the programme if they cannot be settled.

\begin{table}[htbp]
\centering
\caption{Three-Tier Evidence Classification: Demonstrated Results, Analytical Estimates, and Research Targets}
\label{tab:exec-feasibility}
\scriptsize
\setlength{\tabcolsep}{3.5pt}
\begin{tabularx}{\textwidth}{L{1.45in} L{1.55in} Y}
\toprule
\textbf{Category A: Empirically Demonstrated (Code / Data Audited)} & \textbf{Category B: Analytical Estimates (Schema / Complexity Models)} & \textbf{Category C: Open Research Targets (Require Phase 1--3 Validation)} \\
\midrule
\textbf{1. Smooth Cobb-Douglas Convergence (\texttt{pequod-cljs} CSVs):} Across the 40 large-run logs the book cites (\texttt{dep1ex61--100.csv}),\footnotemark{} Year 1 reaches $\rho_i \le 5\%$ in \textbf{11.85 rounds} and $\rho_i \le 3\%$ in \textbf{19.20 rounds}; Year 2 reaches 5\% in \textbf{6.58 rounds}. The 50 small-run files (\texttt{dep1ex01--50.csv}) average \textbf{11.18} Year-1 5\% rounds. All 70 \texttt{dep1ex*.clj} generators (20 stored via Git LFS) were audited. & \textbf{1. Distributed Round Latency ($N=10^6$):} A round moves a 0.24--1.5~GB raw archive of 40k council proposals (Table~\ref{tab:proposal-sizes}) plus a BFT checkpoint, and verifying $10^6$ household reports costs each aggregator $\approx 4$ (coarse) to $\approx 61$ (dense) desktop core-minutes. No round-latency figure is claimed; Kill Gate~2 measures it against a pre-registered threshold. & \textbf{1. Non-Convex \& Leontief Convergence (Kill Gate 1):} Whether any of Algorithms 3a--3d \emph{clears} (not merely satisfies one-sided feasibility; Definition~\ref{def:success}) Leontief + lumpy setup-cost economies in $\le 25$ rounds on pre-registered instances, including Scarf cycles, infeasible plans and endogenous income. \\
\addlinespace[4pt]
\textbf{2. Refutation of Unregularised T\^atonnement on Scarf \& Single-Plant Jumps:} Fails Scarf's 3-good Leontief cycle ($\min \rho = 10.5\%$ over 3,000 rounds) and single indivisible plants ($\min \rho = 50.4\%$). & \textbf{2. One-Good Shapley--Folkman Stress Test ($K=20$):} The \emph{minimum} imbalance across 500 rounds is \textbf{0.40\%}, but the maximum in the last 100 rounds is \textbf{8.4\%} at $a=70$; a sweep $a\in[60,80]$ has median last-100-round max \textbf{6.0\%} and worst \textbf{9.5\%}. This says nothing about interdependent Leontief markets. & \textbf{2. Human Override \& Agency (Kill Gate 1H):} Whether 30--60 self-selected volunteers modify solver drafts, and what that can and cannot show about democratic agency. \\
\addlinespace[4pt]
\textbf{3. NIST PQC Standards and measured \texttt{Prio3} sizes:} FIPS 203/204/205 are standardised. \texttt{draft-irtf-cfrg-vdaf-22} is an Informational IRTF draft with no post-quantum claim. \texttt{Prio3SumVec} sizes and desktop timings were measured with \texttt{libprio} 0.18.1 (Table~\ref{tab:prio3-bench}): a dense 3,000-category, 20-bit report is $975{,}760\text{ B}$ to the leader and costs $\approx 3.5\text{ ms}$ to verify per aggregator on one desktop core. & \textbf{3. Mobile Round Bandwidth (\texttt{Role-LP}):} Quorum signatures at $n=128$ are $\approx 290\text{ kB}$ ($86 \times 3{,}373\text{ B}$). A proposed coarse 200-category, 16-bit household report is $55\text{ kB}$; a dense 3,000-category report does not fit the budget. Kill Gate 3 sets separate light-client and full-verification caps (Table~\ref{tab:gates}). & \textbf{3. Incentive \& Soft-Budget Resistance:} Mitigating capacity understatement (Weitzman 1980), soft-budget appeals (Kornai 1980), and carousel fraud (Hurwicz 1972) remains an open mechanism-design problem.\footnotemark \\
\bottomrule
\end{tabularx}
\end{table}
\addtocounter{footnote}{-1}\footnotetext{Rows 1 and 2 of Category~A are computed by \path{research/analysis/verify_r2.py} (Scarf's cycle by \path{research/analysis/verify_economics.py}) in the project repository.}
\stepcounter{footnote}\footnotetext{Leonid Hurwicz, ``On Informationally Decentralized Systems,'' in \emph{Decision and Organization}, ed.\ C.~B.\ McGuire and R.\ Radner (Amsterdam: North-Holland, 1972), pp.~297--336; Hurwicz, ``The Design of Mechanisms for Resource Allocation,'' \emph{American Economic Review} 63, no. 2 (1973): 1--30.}

\section{The Most Promising Architecture: Why We Reject a Global Blockchain}

\textbf{We explicitly reject building a per-transaction global blockchain (Architecture C), a centralised cloud database (Architecture A), or an unstructured gossip network without finality (Architecture B).}

Participatory planning is not a stream of high-frequency financial trades; it is an \textbf{epoch-synchronised, hierarchical bulk-synchronous coordination process} followed by localized intra-year fulfilment. Furthermore, putting personal consumption requests on a public ledger violates personal privacy, and broadcasting $3.3\text{ KB}$ post-quantum signatures for every household action across all global nodes wastes bandwidth and energy.

Instead, our weighted evaluation (Chapter 6: \textbf{4.07 / 5.00} for D against \textbf{3.10} for a federated SQL service and \textbf{2.54} for a blockchain; expert judgements, not measurements, and A overtakes D when simplicity, throughput and cost carry 70\% of the weight) selects \textbf{Architecture D --- A Four-Layer Hybrid Local-First, Federated P2P, and Selective Epoch-BFT Architecture}:
\begin{enumerate}[leftmargin=1.8em,itemsep=2pt]
  \item \textbf{Layer 4 (Human Deliberation \& Local Solver Edge):} Accessible mobile/desktop Progressive Web App (PWA) and native clients running an embedded interior-point conic/quadratic solver (\texttt{Clarabel.rs} / \texttt{HiGHS} in WASM) for instant local proposal drafting and ``what-if'' collective-good exploration.
  \item \textbf{Layer 3 (Local-First Council \& Household State):} Encrypted local \textbf{SQLite (SQLCipher)} for structured vectors and accounting ledgers + \textbf{Loro/Automerge CRDTs} for offline-capable workplace and neighbourhood council deliberation, authorised via revocable capability certificates (\textbf{CapCert}; format to be chosen, signed under the active \texttt{CryptoSuiteId}).
  \item \textbf{Layer 2 (Federated Aggregation \& Content-Addressed Data Plane):} Peer-to-peer \textbf{QUIC (\texttt{quinn} / \texttt{iroh})} transport organising councils into nested Ward $\to$ City $\to$ Region and Industry Federation \textbf{Two-Root BLAKE3 authenticated aggregation} (proposal Merkle root + category-totals root; Section~\ref{sec:two-root}), paired with content-addressed storage for qualitative council narratives (\texttt{Lists}).
  \item \textbf{Layer 1 (Selective Epoch-BFT Checkpoint \& Deterministic IFB Engine):} A rotating BFT committee of $n=128$ validators (drawn by sortition; an interim classical threshold beacon is adopted and a post-quantum beacon is open, Sections~\ref{sec:pq-urgency} and~\ref{sec:sortition}) that commits one checkpoint per planning round, plus daily operational rollups and scoped replans---\textbf{425 blocks/year at a 10-round plan} ($10+365+50$), or 435 at 20 rounds---anchoring the two aggregation roots, executing the deterministic WASM price-update rule, and settling Sandstr\"om's macro \textbf{Society Account}. ``10--20 checkpoint blocks'' refers only to the annual \emph{planning} rounds, not the Layer-1 archive.
\end{enumerate}

\section{Where AI Adds Value---and Where It Must Be Excluded}

In strict adherence to Parecon's anti-coordinator principles,\footnote{Albert and Hahnel, \emph{Political Economy of Participatory Economics}, Ch.~1, pp.~15--17; \emph{Looking Forward}, Ch.~9 (\texttt{lf9\_2.htm}) and Ch.~10 (\texttt{lf10\_1.htm}).} \textbf{AI must never become an unaccountable planning authority}:
\begin{itemize}[leftmargin=1.6em,itemsep=2pt]
  \item \textbf{Deterministic Core (by design, not yet implemented):} Indicative price adjustments, $SB/SC$ ratios, and Society Account clearing are specified as deterministic fixed-point arithmetic and convex operations-research solvers with no neural-network weights. Local IEEE-754 drafts are advisory only (Section~\ref{sec:fixed-point-bounds}).
  \item \textbf{Accountable Edge Support:} Classical statistical models (hierarchical ETS/ARIMA) forecast SKU size/color splits; local conic QP solvers (\texttt{Clarabel.rs}) automatically draft price-adjusted council revisions; and optional small local open-weight LLMs are restricted to multi-lingual translation, screen-reader accessibility, and summarising qualitative \texttt{Lists} with mandatory links to raw council source text.
\end{itemize}

\section{Estimated Effort, Budget Scenarios, and the Five Key Decisions}

\begin{table}[htbp]
\centering
\caption{Proposed budget scenarios at one assumed loaded rate (for author ratification)}
\label{tab:exec-budget}
\footnotesize
\setlength{\tabcolsep}{4.5pt}
\begin{tabular}{L{1.10in} L{0.86in} L{1.05in} L{1.42in}}
\toprule
\textbf{Phase} & \textbf{Staff \& Time} & \textbf{Non-labour} & \textbf{Estimate \& Gate} \\
\midrule
\textbf{Phase 1}\newline Economic simulator & 1.0 FTE\newline 4--6 months & Compute \$1k--\$3k & \textbf{\$41k--\$93k}. Pre-registration in month~1. Kill Gate 1. \\
\addlinespace[3pt]
\textbf{Phase 1H}\newline Low-tech human probe & 1.0 FTE\newline 3--4 months (in parallel) & Ethics review, consent, compensation, data protection: \$8k--\$25k & \textbf{\$38k--\$85k}. Kill Gate 1H. \\
\addlinespace[3pt]
\textbf{Phase 2}\newline Distributed coordination & 2.5 FTE\newline 4--6 months & WAN cluster \$2k--\$5k & \textbf{\$102k--\$230k}. Kill Gate 2. \\
\addlinespace[3pt]
\textbf{Phase 3}\newline PQC, \texttt{Prio3}, co-op pilot & 4.0 FTE\newline 9--12 months & Audit \$50k--\$150k; hosting \$10k--\$40k; pilot ethics and data protection \$20k--\$40k & \textbf{\$440k--\$950k} for Phase~3 alone. Kill Gate 3. \\
\midrule
\textbf{Programme} & & & \textbf{\$0.62M--\$1.36M} if every gate passes. \\
\bottomrule
\end{tabular}
\par\smallskip
{\scriptsize\raggedright \emph{Basis:} an assumed fully loaded cost of \$120k--\$180k per FTE-year (low end of each range uses the shorter duration and lower rate). This is a planning assumption, not a sourced market rate; it should be replaced by the author's actual pay scale, which in a cooperative may be flat and much lower. Figures from \path{research/analysis/verify_r3.py}. Two to three months is not realistic for building a simulator, pre-registering instances and implementing an endogenous-income solver, so Phase~1 is planned at 4--6 months.\par}
\end{table}

\subsection*{The Five Most Important Decisions for Project Leadership}
\begin{enumerate}[leftmargin=1.8em,itemsep=2pt]
  \item \textbf{Treat Architecture D as the leading long-term hypothesis} (local-first + Two-Root aggregation + epoch checkpoints) and reject a per-transaction blockchain. Use Architecture A (a disclosed trusted server) for Phase 1H and the earliest pilots: it is simpler to operate and explain, and it wins whenever simplicity, throughput and cost dominate the weighting (Section~\ref{sec:arch-sensitivity}).
  \item \textbf{Replace unconstrained Cobb-Douglas with constrained local solvers}, and treat Algorithms 3a--3c as research candidates to be benchmarked in Phase 1 rather than as a committed production solver.
  \item \textbf{Enforce Hard Biophysical Caps ($r^*, E_z^{\max}$) Lexicographically Above PDRM Pollution Pricing.}
  \item \textbf{Design for a ``3-Milestone Hybrid Deliberation Schedule''} separating automated intermediate input-clearing sub-iterations from human value-deliberation milestones.
  \item \textbf{Target Dual Utility:} Support both full-scale Parecon simulation research and immediate deployment for federated worker cooperatives and participatory budgeting communities (``Experiment 2: Developing a Parallel Economy'').\footnote{Albert and Hahnel, \emph{Political Economy of Participatory Economics}, Ch.~6, pp.~129--130.}
\end{enumerate}

\textbf{Final Recommendation:} \textbf{PROCEED WITH PHASE 1 AND PHASE 1H AS KILL GATES}---the Python-first Economic Coordination \& Non-Convexity Simulator (Phase 1) and the Low-Technology Human Participation Probe (Phase 1H) first, each with pre-registered pass and fail conditions (Table~\ref{tab:gates}), holding Phase 2 (distributed Rust nodes) and Phase 3 (post-quantum and \texttt{Prio3} field pilots, with a mandatory \$50k--\$150k external audit) contingent on passing Kill Gates 1 and 1H. A failed gate stops or narrows the programme; it is not reported as a partial success.

% ==============================================================================
% CHAPTER 2: HISTORICAL AND INSTITUTIONAL INSPIRATIONS
% ==============================================================================
\chapter{Historical and Institutional Inspirations: From Council Practice to Protocol Design}

\chapterepigraph{In almost all the industries, factories, mills, workshops, transportation services, public services, and utilities, the rank and file workers, their revolutionary committees, and their syndicates reorganized and administered production, distribution, and public services without capitalists, high salaried managers, or the authority of the state\dots\ They instituted not bourgeois formal democracy but genuine grass roots functional libertarian democracy, where each individual participated directly in the revolutionary reorganization of social life.}{Gaston Leval, \emph{Espagne Libertaire} (1971), p.~11, translated in Sam Dolgoff, ed., \emph{The Anarchist Collectives: Workers' Self-Management in the Spanish Revolution, 1936--1939} (1974), Ch.~1, p.~6\protect\footnotemark}
\footnotetext{Gaston Leval, \emph{Espagne Libertaire (1936--1939)} (Paris: \'Editions du Monde Libertaire, 1971), p.~11; verbatim English translation in Sam Dolgoff, ed., \emph{The Anarchist Collectives: Workers' Self-Management in the Spanish Revolution, 1936--1939} (New York: Free Life Editions, 1974), Ch.~1, p.~6.}

\section{Epistemological Premise: Institutions Before Cryptography}

A recurring failure mode of contemporary ``decentralised governance'' (DAO) architectures is the conflation of \textbf{cryptographic trustlessness} with \textbf{democratic self-management}. Recording a tally on an immutable ledger does not render an economy democratic if:
\begin{enumerate}[leftmargin=1.8em,itemsep=2pt]
  \item Proposal drafting and agenda-setting are monopolised by an unaccountable technical or managerial elite---what Albert and Hahnel diagnose as the \textbf{coordinator class}.\footnote{Albert and Hahnel, \emph{Political Economy of Participatory Economics}, Ch.~1, pp.~14--17, 20--22; \emph{Looking Forward}, Prologue (``Coordinatorism'' and ``The Origins of Coordinatorism'').}
  \item Decision-making influence is proportional to token wealth or meeting endurance rather than proportional to the degree participants are affected.
  \item Wealthier production collectives hoard superior inherited productive endowments and refuse to equalise resources across regions.
\end{enumerate}

In this chapter, we examine five historical and contemporary institutional traditions as empirical laboratories of democratic coordination---analyzing both their institutional breakthroughs and their structural failures---to derive binding architectural rules for DPEP.

\section{The Spanish Revolution (1936--1939): Industrial and Agrarian Collectives}

Following the defeat of the Nationalist military coup in July 1936 across Catalonia, Aragon, Valencia, and parts of Castile and Andalusia, workers and peasants affiliated with the anarchist CNT (\emph{Confederaci\'on Nacional del Trabajo}) and socialist UGT (\emph{Uni\'on General de Trabajadores}) collectivised a large share of Catalan industry (figures of 70--80\% are commonly repeated in the secondary literature; this monograph has not traced them to a primary count) and several hundred agrarian collectives in Aragon. Estimates of participation vary widely and come mostly from participants and sympathisers: Frank Mintz gives a minimum of about $1.84$ million collectivists ($1.08$ million industrial, $0.76$ million agrarian), revised downwards in 1977 from his earlier $2.44$--$3.2$ million, while the eyewitness and CNT militant Gaston Leval gave figures of up to $3$ million. Mintz's figures are cited here from the bibliographic record and were not checked against the volume.\footnote{Frank Mintz, \emph{Anarchism and Workers' Self-Management in Revolutionary Spain}, trans. Paul Sharkey (Oakland, CA: AK Press, 2013 [1970/1977]); Gaston Leval, \emph{Collectives in the Spanish Revolution} (London: Freedom Press, 1975); Sam Dolgoff, ed., \emph{The Anarchist Collectives: Workers' Self-Management in the Spanish Revolution, 1936--1939} (New York: Free Life Editions, 1974), available in full at \url{https://theanarchistlibrary.org/library/sam-dolgoff-editor-the-anarchist-collectives}.}

\subsection{Internal Self-Management, Delegation, and Recall}
Across both Barcelona factories (tramways, textiles, metalworks, water/gas/electric utilities) and Aragonese villages (Alca\~niz, Calanda, Bin\'efar, Graus, Mas de las Matas), supreme authority resided in the weekly or biweekly \textbf{General Assembly} (\emph{Asamblea General}) of all workers or villagers.\footnote{Dolgoff, \emph{The Anarchist Collectives}, Ch.~2 and Ch.~9; Murray Bookchin, ``Introductory Essay,'' in Dolgoff, \emph{The Anarchist Collectives}, pp.~xxvii--xxx.} Assemblies elected specialized administrative committees (\emph{Comit\'es de Administraci\'on}) to manage daily coordination, input procurement, and statistical reporting. As documented by Bookchin, Peirats, and Leval, these delegates operated under strict anti-bureaucratic constraints:
\begin{quote}
\itshape\small ``Where such popular assemblies were possible, their decisions were executed by a committee elected from the assembly. Apparently, the right to recall committee members was taken for granted and they certainly enjoyed no privileges, emoluments, or institutional power\dots\ It remained a cardinal principle of Spanish anarchists never to pay their delegates, even when the CNT numbered a million members. Normally, the responsibilities of elected delegates had to be discharged after working hours.''\footnote{Bookchin, ``Introductory Essay,'' in Dolgoff, \emph{The Anarchist Collectives}, p.~xxviii.}
\end{quote}

\subsection{Inter-Collective Coordination, Reorganization, and Accounting}
\begin{itemize}[leftmargin=1.6em,itemsep=2pt]
  \item \textbf{Regional Equalisation in Aragon:} The Regional Defence Council of Aragon was created at a plenary of unions and militia columns in Bujaraloz on 6 October 1936. The \textbf{Aragon Federation of Collectives} was constituted later, at the Caspe congress of 14--15 February 1937. Accounts of its size differ by a factor of five, and all come from sympathetic sources: Leval, who attended, lists 24 cantonal federations covering 275 villages and 141,430 member families; Peirats, citing Souchy, gives 500 collectives with 433,000 members; a widely reproduced summary gives 500 delegates representing 80,000 collectivists. The federation pooled surpluses at cantonal depots, moved machinery and seed between villages, and traded directly with Catalan syndicates. The Council was dissolved by decree on 10--11 August 1937, after Republican troops under L\'ister entered Aragon.\footnote{Leval, \emph{Collectives in the Spanish Revolution}, Ch.~V (Caspe congress); Jos\'e Peirats, in Dolgoff, \emph{The Anarchist Collectives}, Ch.~9 (giving Souchy's estimate); the 80,000 figure appears without a source in the Wikipedia article ``Regional Defence Council of Aragon'' and is reported only as a range bound. The October 1936 date belongs to the Council, not the Federation.}
  \item \textbf{Unified Export Coordination in Valencia (CLUEA):} In the Levante, the CNT and UGT established the \emph{Consejo Levantino Unificado de Exportaci\'on Agr\'icola} (CLUEA), federating $\sim$270 local orange and citrus committees to standardize grading, rail/maritime transport, and export clearing.\footnote{Mintz, \emph{Anarchism and Workers' Self-Management in Revolutionary Spain}, Ch.~4.}
  \item \textbf{Industrial Recomposition in Catalonia:} In Barcelona, the CNT's industrial unions (\emph{Sindicatos \'Unicos de Ramo}) consolidated hundreds of small, unsanitary workshops in woodworking, baking, tanning, and optic/glass production into modern integrated plants, eliminating redundant middlemen while preserving full employment.\footnote{Augustin Souchy, ``Collectivizations in Catalonia,'' in Dolgoff, \emph{The Anarchist Collectives}, Ch.~7.}
  \item \textbf{Accounting and the ``Family Wage'':} Many Aragonese villages abolished official \emph{pesetas} internally, distributing abundant staples free according to need (Maxim 4) and issuing a family consumption booklet (\emph{libreta de consumo}) or local points calibrated to family size for scarce goods.\footnote{Leval, \emph{Collectives in the Spanish Revolution}, pp.~165--210.} Inter-collective trade relied on standardised valuation tables (indexed in wheat equivalents or shadow pesetas) cleared through cantonal ledgers.
\end{itemize}

\subsection{Historical Limitations and Structural Failures}
Despite remarkable wartime productivity in munitions, transport, and agriculture, the Spanish collectives suffered four critical weaknesses that directly inform DPEP's design:
\begin{enumerate}[leftmargin=1.8em,itemsep=2pt]
  \item \textbf{``Collective Egoism'' (\emph{Neocapitalismo Colectivo}) Between Rich and Poor Enterprises:} In urban Catalonia, prior to the socialisation decrees of late 1937, collectives remained financially autonomous inside a quasi-market environment. Wealthy collectives that inherited modern machinery or lucrative monopolies (such as the Barcelona tramways, power utilities, and luxury hotels) enjoyed large surpluses, while raw-material-starved textile and construction collectives struggled to pay base wages.\footnote{Jos\'e Peirats, in Dolgoff, \emph{The Anarchist Collectives}, Ch.~6; Mintz, \emph{Anarchism and Workers' Self-Management}.} \textbf{Protocol Lesson:} This failure is why Albert and Hahnel (1991a, Ch.~5, pp.~84--87) and Sandstr\"om (2020) insist that productive assets belong to the \emph{Productive Commons}, and that every Worker Council be charged an opportunity-cost rental fee ($P_k \cdot k_j$) on its inherited capital stock.
  \item \textbf{Burnout from Unpaid Evening Governance:} Requiring delegates to perform all federation administration after a full day of physical factory labour caused severe exhaustion and favoured childless militants over workers with caregiving duties. \textbf{Protocol Lesson:} Governance and facilitation work must be counted inside paid work hours as part of every citizen's Balanced Job Complex.
  \item \textbf{Lack of Unified Shadow Pricing for Complex Intermediate Chains:} While cantonal barter cleared simple agricultural staples, urban industries lacked an iterative mechanism to compute social opportunity costs across hundreds of interdependent intermediate inputs once pre-war market prices grew stale.
  \item \textbf{Wartime militia coercion, smallholder conflict, and ``workers against work'' (Casanova 2005; Seidman 1991):} Internal weaknesses were not merely accounting bugs. Casanova (2005) and Seidman (1991, 2002), consulted at bibliographic level only, are the standard critical accounts; they and other secondary sources describe Aragonese collectives formed under armed militia columns, conflict with ``individualist'' smallholders, absenteeism in Barcelona factories, credit starvation by the central government, dissolution in Aragon by L\'ister's 11th Division in August 1937, and Franco's victory in 1939. No page-level claim is made for either author.
\end{enumerate}

\section{The Paris Commune (March 18 -- May 28, 1871)}

During its 72 days of existence before its bloody suppression during \emph{La Semaine Sanglante}, the Paris Commune established foundational principles of council democracy celebrated in Karl Marx's \emph{The Civil War in France} (1871).\footnote{Karl Marx, \emph{The Civil War in France} (London: General Council of the International Workingmen's Association, 1871); John Merriman, \emph{Massacre: The Life and Death of the Paris Commune of 1871} (New York: Basic Books, 2014).} A rigorous institutional analysis must distinguish \textbf{what the Commune actually implemented} from \textbf{unimplemented proposals and later theoretical extrapolations}.

\begin{table}[htbp]
\centering
\caption{Paris Commune (1871): Actually Implemented Institutions vs. Proposals \& Extrapolations}
\label{tab:paris-commune}
\footnotesize
\setlength{\tabcolsep}{3.5pt}
\begin{tabularx}{\textwidth}{L{0.95in} L{1.85in} Y}
\toprule
\textbf{Dimension} & \textbf{Actually Implemented (March--May 1871)} & \textbf{Proposed on Paper or Later Extrapolation} \\
\midrule
\textbf{Representation \& Recall} & Municipal elections across 20 \emph{arrondissements} (March 26, 1871) by universal \emph{male} suffrage; councilors declared revocable at any time; merged legislative and executive work across 10 Commissions.\footnotemark & Formal recall was rarely exercised during the 72-day military siege; women were excluded from formal voting despite leading grass-roots defence and cooperative organising via the \emph{Union des Femmes} (\'Elisabeth Dmitrieff, Nathalie Lemel). \\
\addlinespace[4pt]
\textbf{Remuneration Cap} & \textbf{Decree of April 2, 1871:} Capped maximum salaries of all Commune officials and municipal administrators at \textbf{6,000 francs/year} (equivalent to a skilled Parisian workman's wage) and banned multiple stipends. & Did not abolish wage hierarchies inside private Parisian enterprises or establish balanced job complexes inside workshops. \\
\addlinespace[4pt]
\textbf{Workplace Self-Management} & \textbf{Decree of April 16, 1871} (drafted by L\'eo Frankel's Commission of Labour and Exchange): Ordered a census of workshops abandoned by fleeing owners and their transfer to worker cooperative associations; abolished night baking (April 20) and employer wage fines (April 27).\footnotemark & Only a few dozen cooperative workshops (e.g., the Louvre armaments workshop electing its own council and foremen; women's uniform cooperatives) actually began operating before Versailles troops entered Paris on May 21. \\
\addlinespace[4pt]
\textbf{National Federation \& Credit} & Operated strictly within besieged Paris; left the \emph{Banque de France} intact (negotiating advances via Charles Beslay rather than socializing credit). & \textbf{Declaration to the French People (April 19, 1871)} (drafted by Proudhonian Pierre Denis) and Marx envisioned a national federation of autonomous communes coordinating via delegated national assemblies---never realized due to the crushing of provincial communes in Lyon and Marseille. \\
\addlinespace[4pt]
\textbf{Crisis Centralisation} & On May 1, 1871, under bombardment, the Jacobin/Blanquist majority voted (45 to 23) to create a 5-member \textbf{Committee of Public Safety} with sweeping powers over the objection of the Federalist minority (Frankel, Varlin, Longuet, Courbet). & Demonstrates the acute vulnerability of council democracies during crisis to revert to emergency command authority unless crisis-mode bounds (TTLs and post-incident audits) are constitutionally hard-coded. \\
\bottomrule
\end{tabularx}
\end{table}
\addtocounter{footnote}{-1}
\footnotetext{Marx, \emph{The Civil War in France}, Section III.}
\stepcounter{footnote}
\footnotetext{Merriman, \emph{Massacre: The Life and Death of the Paris Commune of 1871}, Ch.~5.}

\section{Classical Athenian Democracy ($\sim$508/7 -- 322 BCE)}

Over nearly two centuries from the reforms of Kleisthenes (508/7 BCE) to the Macedonian conquest (322 BCE), classical Athens developed the most sophisticated institutional machinery for direct citizen self-governance in antiquity, analyzed definitively by Mogens Herman Hansen.\footnote{Mogens Herman Hansen, \emph{The Athenian Democracy in the Age of Demosthenes: Structure, Principles, and Ideology} (Oxford: Blackwell / Norman: University of Oklahoma Press, 1991).}

\subsection{Core Institutional Mechanisms Adaptable to DPEP}
\begin{enumerate}[leftmargin=1.8em,itemsep=2pt]
  \item \textbf{Direct Assembly (\emph{Ekklesia}) vs.\ Norm-Setting (\emph{Nomothetai}):} Open to all adult male citizens ($\sim$30,000--40,000 eligible in the 4th century BCE, with $\sim$6,000 attending the 40 annual meetings on the Pnyx), the \emph{Ekklesia} voted on decrees (\emph{psephismata}), while permanent constitutional laws (\emph{nomoi}) were ratified by a separate panel of 1,000 \emph{nomothetai} selected by lot.
  \item \textbf{Sortition (\emph{Kleroterion}) and Strict Rotation:} Roughly 600 of the $\sim$700 domestic magistrates (\emph{archai}) and the 6,000 annual jurors (\emph{dikastai}) were selected by \textbf{lot} using bronze allotment machines (\emph{kleroteria}) rather than election---following Aristotle's principle (\emph{Politics} 1294b) that elections are intrinsically oligarchic (favoring the wealthy, famous, and rhetorically polished) whereas sortition is democratic.\footnote{Hansen, \emph{Athenian Democracy in the Age of Demosthenes}, Ch.~9, pp.~230--245.} Almost all offices were restricted to a non-repeatable one-year term (or repeatable only once in a lifetime for the \emph{Boule}), preventing a permanent coordinator class.
  \item \textbf{Probouleutic Agenda Facilitation (\emph{Boule} of 500):} Five hundred citizens (50 from each of the 10 Kleisthenic tribes, chosen by lot from local \emph{demes}) served for one year on the Council (\emph{Boule}), with each tribe serving as the standing executive committee (\emph{prytaneis}) for $1/10$th of the year and a new chair (\emph{epistates}) drawn by lot every 24 hours. Crucially, the \emph{Boule}'s role was \emph{probouleutic}: it prepared structured proposals (\emph{probouleumata}) for the \emph{Ekklesia} to debate, amend, or reject, but could not impose major policies on its own.\footnote{Hansen, \emph{Athenian Democracy in the Age of Demosthenes}, Ch.~10, pp.~246--265.} This is the closest historical \emph{analogy} (not a direct institutional ancestor) of Parecon's \textbf{Facilitation Boards}.
  \item \textbf{Public Scrutiny and Mandatory Audit (\emph{Dokimasia}, \emph{Euthynai}, \emph{Graphe Paranomon}):} Every magistrate underwent pre-office eligibility scrutiny (\emph{dokimasia}) and a mandatory end-of-term financial and conduct audit (\emph{euthynai}) before 10 \emph{logistai} (accountants) and 10 \emph{euthynoi} chosen by lot.\footnote{Hansen, \emph{Athenian Democracy in the Age of Demosthenes}, Ch.~9, pp.~218--224.}
  \item \textbf{Civic Remuneration (\emph{Misthophoria}):} Introduced by Pericles for jurors and councilors and expanded in the 4th century BCE to assembly attendance (\emph{misthos ekklesiastikos}, 3 to 6 obols per day), ensuring working citizens (\emph{thetes}) could afford to participate.
\end{enumerate}

\subsection{Explicit Examination of Structural Exclusions}
Athenian democracy excluded \textbf{women} (who lacked political citizenship and independent property control), \textbf{metics} (resident non-citizens, $\sim$20,000--30,000 artisans and merchants), and \textbf{enslaved people} ($\sim$80,000--150,000 forced laborers in the Laurion silver mines, agriculture, crafts, and households), leaving only $\sim$10--15\% of the total population with full political rights.\footnote{Hansen, \emph{Athenian Democracy in the Age of Demosthenes}, Ch.~4, pp.~86--124.} Furthermore, as Moses Finley and G.~E.~M.\ de Ste.\ Croix demonstrate, Athens' fiscal capacity to fund \emph{misthophoria} and public works depended heavily on the \textbf{Laurion silver mines (worked by tens of thousands of enslaved miners)} and 5th-century \textbf{Delian League imperial tribute (\emph{phoros})}.\footnote{M.~I.\ Finley, \emph{Economy and Society in Ancient Greece} (London: Chatto \& Windus, 1981), Ch.~5; G.~E.~M.\ de Ste.\ Croix, \emph{The Class Struggle in the Ancient Greek World} (Ithaca, NY: Cornell University Press, 1981).} In a universal, non-extractive participatory economy, every hour spent in deliberation and \emph{euthynai} audits carries a real opportunity cost in productive labour; governance time must therefore be budgeted inside paid Balanced Job Complexes and economised carefully, while sortition should draw from all citizens regardless of hardware ownership (the beacon design is open; Section~\ref{sec:sortition}).

\section{Occupy Wall Street (2011) and Additional Precedents}

\begin{itemize}[leftmargin=1.5em,itemsep=1pt,topsep=2pt]
  \item \textbf{Occupy Wall Street (2011), Freeman's ``Structurelessness,'' and Assembly Scaling:} In small groups without formal rules, Jo Freeman (1972) showed that power concentrates in unaccountable informal friendship cliques (``the tyranny of structurelessness'')---a sociological dynamic that asynchronous software alone does \emph{not} eliminate.\footnote{Graeber, \emph{The Democracy Project}, Chs.~1--3; Freeman, ``The Tyranny of Structurelessness,'' pp.~151--164; Jane J.\ Mansbridge, \emph{Beyond Adversary Democracy} (New York: Basic Books, 1980); Robert A.\ Dahl and Edward R.\ Tufte, \emph{Size and Democracy} (Stanford: Stanford University Press, 1973); Albert, \emph{Parecon: Life After Capitalism} (2003), Ch.~5.} Mansbridge (1980) and Dahl \& Tufte (1973) are the classic treatments of assembly scale; they were not consulted directly, and no specific size threshold is attributed to them. The $\sim$50--200 neighbourhood-council bound and the 50--150 speaker figure in Table~\ref{tab:hist-matrix} are \textbf{working design heuristics} from speaking-time arithmetic (one hour shared among 60 speakers is one minute each), not sourced empirical findings. Albert (2003, Ch.~5) argues why unanimity vetoes fit small-team hiring while majority/supermajority rules with minority reconsideration delays fit workplace investments.
  \item \textbf{Mondrag\'on Cooperative Corporation (1956--Present):} Across roughly 90 cooperatives and a group workforce commonly reported at 70,000--80,000, of whom a substantially smaller number are cooperative worker-owners (\emph{socios}) because much of the retail (Eroski) and overseas workforce are non-member employees (membership figures vary by year and source and are not independently checked here), Mondrag\'on separates the elected \textbf{Audit / Watchdog Committee (\emph{Comisi\'on de Vigilancia})} from the frontline workers' \textbf{Social Council (\emph{Consejo Social})},\footnote{William Foote Whyte and Kathleen King Whyte, \emph{Making Mondrag\'on}, 2nd ed. (Ithaca, NY: ILR Press, 1991), consulted at bibliographic level.} illustrating both federated industrial self-management and the erosion of cooperative membership under global market competition.
  \item \textbf{Porto Alegre Participatory Budgeting (1989--2017 Suspension):} Combined regional assemblies and recallable delegates with an objective \textbf{Quality of Life Need Matrix} to redirect municipal investment to poor districts, before being \textbf{suspended in 2017} following municipal coalition shifts, fiscal austerity, and $\sim$2,400 unexecuted citizen project requests.
  \item \textbf{Rojava / DAANES (2012--2026 Damascus Integration Accord):} Implemented nested communes, co-chair gender parity, and autonomous women's councils (\emph{Kongra Star}),\footnote{Hahnel, Bohmer, and Chowdhury, ``Reproductive Labour,'' Ch.~10 in Hahnel, \emph{Democratic Economic Planning}, p.~200.} before military and geopolitical pressure led to the \textbf{30 January 2026 SDF--Damascus integration agreement} merging DAANES civil and security institutions into the Syrian state.
  \item \textbf{Ostrom's Common-Pool Resource (CPR) Institutions (1990):} Elinor Ostrom's study of enduring self-governed commons (e.g., Valencia's \emph{Tribunal de las Aguas}) established eight design principles: defined boundaries, proportional equivalence, collective-choice arrangements, accountable monitoring, graduated sanctions, rapid local conflict resolution, right to organise, and \textbf{nested enterprises}.\footnote{Elinor Ostrom, \emph{Governing the Commons} (Cambridge: Cambridge University Press, 1990), Ch.~3, p.~90.}
\end{itemize}

\subsection*{Synthesis: Comparative Institutional-to-Protocol Matrix}
Table \ref{tab:hist-matrix} synthesizes the eight core institutional mechanisms analyzed across Chapters 2.2--2.5, mapping their historical strengths and limitations directly to concrete capabilities in the Distributed Participatory Economics Protocol (DPEP).

\begin{table}[ht!]
\centering
\caption{Comparative Matrix: Institutional Mechanisms Mapped to DPEP Protocol Capabilities}
\label{tab:hist-matrix}
\vspace{3pt}
\scriptsize
\setlength{\tabcolsep}{4.5pt}
\renewcommand{\arraystretch}{1.20}
\begin{tabularx}{\textwidth}{L{1.12in} L{1.02in} L{1.02in} Y}
\toprule
\textbf{Mechanism \& Precedents} & \textbf{Strengths \& Advantages} & \textbf{Historical Limitations} & \textbf{Concrete DPEP Protocol Capability} \\
\midrule
\textbf{1. General Assembly}\newline {\tiny\color{SlateDark}Spain 1936; Athens \emph{Ekklesia}; Occupy; Mondrag\'on} & Builds face-to-face solidarity; surfaces qualitative context. & Strains as attendance grows (working heuristic: $\sim$50--150 speakers; Section~2.5); meeting fatigue. & Scoped strictly to work teams (10--30) and neighborhood councils (50--200) with async pre-reading queues. \\
\addlinespace[4pt]
\textbf{2. Revocable Delegation}\newline {\tiny\color{SlateDark}Paris Commune 1871; Spain CNT; Porto Alegre PB} & Enables multi-tier federation scaling while tethering delegates to base. & Delegates ossify if recall is procedurally slow or difficult. & \textbf{Scoped CapCert Delegate Certificates:} Issued by council $k$-of-$n$ multi-signature; revocable in a single council vote. \\
\addlinespace[4pt]
\textbf{3. Sortition + Rotation}\newline {\tiny\color{SlateDark}Classical Athens (\emph{Boule}, \emph{Logistai}, \emph{Dikasteria})} & Prevents factional elections, charismatic capture, and coordinator elites. & Unsuitable for rare specialist tasks unless drawn from qualified pool. & \textbf{Sortition (beacon open):} Intended to select rotating IFB oversight, audit panels, and dispute juries (Section~\ref{sec:sortition}). \\
\addlinespace[4pt]
\textbf{4. Mandatory Audit}\newline {\tiny\color{SlateDark}Athens \emph{Euthynai} (\emph{Logistai}); Mondrag\'on \emph{Comisi\'on de Vigilancia}} & Systematically detects bias or unauthorised overrides. & Requires tamper-evident records and independent auditors. Mondrag\'on's elected Social Council is a workers' voice body, not the audit committee. & \textbf{Deterministic WASM + Two-Root Audit Logs:} Any node with the leaf archive can recompute IFB price updates; mobile light clients verify headers only. \\
\addlinespace[4pt]
\textbf{5. Endowment Equalization}\newline {\tiny\color{SlateDark}Aragon Depots (1936); Mondrag\'on Funds} & Prevents collectives with superior machinery from extracting rents. & Failed in 1936 Barcelona where rich collectives hoarded revenues. & \textbf{Stand-Alone Capital Asset Ledger:} Automatically debits each $WC_j$'s $SC_j$ for capital/resource user rights ($P_k, P_r$). \\
\addlinespace[4pt]
\textbf{6. Proportional \& Caucus}\newline {\tiny\color{SlateDark}Rojava Co-Chairs \& Women's Councils} & Protects disproportionately affected groups from majority tyranny. & Can deadlock if scope of veto is unbounded. & \textbf{Impact-Weighted \& Dual-Key Governance Rules} (PEPE Ch.~2 p.~24; DEP Ch.~10 p.~200). \\
\addlinespace[4pt]
\textbf{7. Iterative T\^atonnement}\newline {\tiny\color{SlateDark}IETF/Git; Lange-Arrow-Hurwicz; Parecon IFB} & Scales quantitative clearing to millions asynchronously. & Pure prices strip out human context unless paired with qualitative text. & \textbf{Two-Root BLAKE3 Aggregation + Content-Addressed Qualitative \texttt{Lists}}. \\
\addlinespace[4pt]
\textbf{8. Graduated Sanctions}\newline {\tiny\color{SlateDark}Ostrom CPRs (Valencia \emph{Tribunal de las Aguas})} & Intended to keep routine disputes local (Ostrom 1990); no sourced ``95\%'' rate is claimed here. & Requires clear invariants and fast local mediation. & \textbf{Tiered Variance Alerts $\to$ UFB Mediation $\to$ Industry/Environmental Court}. \\
\bottomrule
\end{tabularx}
\end{table}

% ==============================================================================
% CHAPTER 3: SOURCE-GROUNDED PARTICIPATORY ECONOMICS REQUIREMENTS
% ==============================================================================
\chapter{Source-Grounded Requirements of Participatory Economics}

\chapterepigraph{Until now, however, descriptions of alternatives have been unconvincing. Here Michael Albert and Robin Hahnel support the libertarian socialist tradition by presenting a rigorous, well-defined model of how producers and consumers could democratically plan their interconnected activities.}{Michael Albert and Robin Hahnel, \emph{The Political Economy of Participatory Economics} (1991), Front Flap / JSTOR Abstract\protect\footnotemark}
\footnotetext{Albert and Hahnel, \emph{Political Economy of Participatory Economics}, book overview (\url{https://www.jstor.org/stable/j.ctv131bvr2}).}

\section{Source Provenance and Classification Methodology}

Before formulating mathematical equations or distributed protocols, we establish strict traceability to the foundational texts. Every requirement in Table \ref{tab:req-matrix} is classified into one of four provenance tiers:
\begin{itemize}[leftmargin=1.6em,itemsep=2pt]
  \item \textbf{\texttt{[ORIGINAL-1991]}:} Explicitly specified in Michael Albert and Robin Hahnel's two 1991 volumes, \emph{The Political Economy of Participatory Economics} (PEPE) and \emph{Looking Forward} (LF).\footnote{Albert and Hahnel, \emph{Political Economy of Participatory Economics} (1991), print pagination verified via library TOC scan (\texttt{research/sources/pepe\_dandelon.pdf}, pp.~1--132) and full text analyzed via ZNetwork (\texttt{intro.htm}, \texttt{1.htm}--\texttt{6.htm}); Albert and Hahnel, \emph{Looking Forward} (1991), full text analyzed via ZNetwork (\texttt{lfIntroduction.htm}, \texttt{lf1\_1.htm}--\texttt{lf10\_1.htm}).}
  \item \textbf{\texttt{[INTERPRETATION-LATER]}:} Institutional elaborations, formalizations, and accounting procedures introduced in Michael Albert's \emph{Parecon: Life After Capitalism} (PLAC, 2003), Robin Hahnel's \emph{Democratic Economic Planning} (DEP, 2021) and \emph{A Participatory Economy} (2022), Anders Sandstr\"om's \emph{Anarchist Accounting} (2020/2023), Mitchell Szczepanczyk's simulation work (2023), and the Participatory Economy Project's User Story Mapping (2023).\footnote{Hahnel, \emph{Democratic Economic Planning} (2021); Robin Hahnel, \emph{A Participatory Economy} (Chico, CA: AK Press, 2022); Sandstr\"om, \emph{Anarchist Accounting} (2020); Szczepanczyk, ``Pseudocode and Algorithms'' (2023); Participatory Economy Project, ``Participatory Planning: User Story Mapping'' (2023), \url{https://participatoryeconomy.org/project/participatory-planning-user-story-mapping/}.}
  \item \textbf{\texttt{[TECH-EXTENSION]}:} Technological extensions required to execute the model over a modern, decentralised, post-quantum, privacy-preserving computational network.
  \item \textbf{\texttt{[UNRESOLVED]}:} Open theoretical or institutional questions requiring simulation or empirical human research (examined in full in Chapter 10).
\end{itemize}

\section{Normative Foundations and Endogenous Preference Theorems}

In Chapter 1 of \emph{The Political Economy of Participatory Economics} (pp.~8--11), Albert and Hahnel evaluate economic institutions against five valuative criteria: \textbf{Equity}, \textbf{Self-Management}, \textbf{Solidarity}, \textbf{Variety}, and \textbf{Efficiency}.\footnote{Albert and Hahnel, \emph{Political Economy of Participatory Economics}, Ch.~1, pp.~8--9, 11.} For distributive equity, they distinguish four mutually exclusive maxims (PEPE Ch.~1, p.~8; Ch.~3, p.~44; DEP Ch.~1, pp.~16--22):
\begin{enumerate}[leftmargin=1.8em,itemsep=2pt]
  \item \textbf{Maxim 1 (Capitalism):} Payment according to the value of the contribution of one's personal labour \emph{and} productive property owned (rejected due to inheritance---the ``Rockefeller grandson problem''---luck, and compounding wealth accumulation).\footnote{Hahnel, \emph{Democratic Economic Planning}, Ch.~1, pp.~16--18; Albert and Hahnel, \emph{Political Economy of Participatory Economics}, Ch.~3, pp.~45--46.}
  \item \textbf{Maxim 2 (Market Socialism / Coordinatorism):} Payment according to the value of one's personal labour contribution alone (rejected because differences in marginal productivity stem overwhelmingly from the ``genetic lottery,'' luck, job assignment, and publicly funded training rather than personal sacrifice---illustrated by their ``Mozart vs.\ Salieri'' and ``Dave Butz'' examples).\footnote{Albert and Hahnel, \emph{Political Economy of Participatory Economics}, Ch.~3, pp.~46--47.}
  \item \textbf{Maxim 3 (Participatory Economics):} Payment according to \textbf{effort and personal sacrifice} expended in socially valued labour (longer hours, higher intensity, or more onerous/dangerous conditions prior to full BJC equalization).
  \item \textbf{Maxim 4 (Humane Solidarity):} Payment according to \textbf{need}---incorporated into Parecon via universal free public goods (healthcare, education) and council-approved special-need allowances (PEPE Ch.~3, p.~47).
\end{enumerate}

Crucially, drawing on the formal results of \emph{Quiet Revolution in Welfare Economics} (1990), consulted here through PEPE's summary,\footnote{Robin Hahnel and Michael Albert, \emph{Quiet Revolution in Welfare Economics} (Princeton, NJ: Princeton University Press, 1990), Chs.~4, 6, 7; summarised in \emph{Political Economy of Participatory Economics}, Ch.~1, pp.~9--11.} Albert and Hahnel argue that because human preferences are \textbf{endogenous} (people rationally mold their future characteristics $\Delta C$ in anticipation of future terms of supply), any institution that systematically mischarges for activities relative to true social opportunity cost causes: (1) a \textbf{snowballing divergence from Pareto optimality} over time, (2) systematic \textbf{warping of human development trajectories}, and (3) ex post invisibility of the bias to participants. Markets systematically underprice negative externalities and overprice public goods (warping preferences toward anti-social individualism), while central planning systematically under-supplies self-managed work roles (warping workers toward apathy and coordinator submission).\footnote{Albert and Hahnel, \emph{Political Economy of Participatory Economics}, Ch.~1, pp.~12--17; Ch.~5, pp.~98--105.}

\section{Formal Traceability Matrix of Requirements}

{\scriptsize\setlength{\tabcolsep}{3.5pt}
\begin{longtable}{L{0.78in} L{1.95in} L{0.82in} L{0.95in}}
\caption{Master Requirements Matrix Derived from Primary and Secondary Literature}
\label{tab:req-matrix} \\
\toprule
\textbf{Req ID} & \textbf{Requirement Specification} & \textbf{Classification} & \textbf{Primary Citations} \\
\midrule
\endfirsthead
\multicolumn{4}{l}{\scriptsize\itshape Continued from previous page} \\
\toprule
\textbf{Req ID} & \textbf{Requirement Specification} & \textbf{Classification} & \textbf{Primary Citations} \\
\midrule
\endhead
\midrule
\multicolumn{4}{r}{\scriptsize\itshape Continued on next page} \\
\endfoot
\bottomrule
\endlastfoot
\textbf{REQ-NORM-01} & \textbf{Productive Commons \& User-Right Fees:} Abolish private ownership of productive property; all natural resources and capital stocks belong to society. Worker councils acquire temporary user rights during annual planning and pay opportunity-cost charges ($P_r, P_k, P_l$). & \texttt{[ORIGINAL-1991]} + \texttt{[INTERP-LATER]} & PEPE Ch.~1 (pp.~17--20), Ch.~5 (p.~85); DEP Ch.~6 (pp.~98--102), Ch.~12 (p.~244); Sandstr\"om (2020). \\
\addlinespace[3pt]
\textbf{REQ-NORM-02} & \textbf{Proportional Self-Management:} Decision-making input proportional to the degree affected by an economic choice. & \texttt{[ORIGINAL-1991]} & PEPE Ch.~1 (p.~8), Ch.~2 (p.~24); LF Ch.~1; PLAC Ch.~5; DEP Ch.~1 (p.~14). \\
\addlinespace[3pt]
\textbf{REQ-NORM-03} & \textbf{Effort \& Sacrifice Remuneration:} Consumption budgets proportional to effort/sacrifice (Maxim 3) tempered by universal free goods and need exceptions (Maxim 4), with council average effort ratings capped by $SB_j/SC_j$. & \texttt{[ORIGINAL-1991]} + \texttt{[INTERP-LATER]} & PEPE Ch.~3 (pp.~44--55); DEP Ch.~6 (pp.~108--117); Sandstr\"om (2023). \\
\addlinespace[3pt]
\textbf{REQ-NORM-04} & \textbf{Balanced Job Complexes (BJCs):} Tasks partitioned into subsets $\{PIPA\}_i$ ordered by empowerment/desirability and combined into balanced complexes within and across workplaces via Job Complex Committees. & \texttt{[ORIGINAL-1991]} & PEPE Ch.~2 (pp.~25--30), Ch.~5 (pp.~82--83); LF Chs.~1--2. \\
\addlinespace[3pt]
\textbf{REQ-ACT-01} & \textbf{Nested Consumer Councils \& Federations:} Individuals/households $\to$ Neighbourhood $CC_h \to$ Ward $\to$ City/County $\to$ State/Region $\to$ National Federation, requesting public goods matched to spatial scale. & \texttt{[ORIGINAL-1991]} & PEPE Ch.~3 (pp.~40--42), Ch.~5 (pp.~81--82); LF Chs.~3, 4, 8. \\
\addlinespace[3pt]
\textbf{REQ-ACT-02} & \textbf{Worker Councils \& Dual Federations:} Democratic workplace councils federated both by \emph{industry} (scale economies, R\&D, $SB/SC$ benchmarking) and by \emph{region} (local production externalities). & \texttt{[ORIGINAL-1991]} & PEPE Ch.~2 (pp.~23--25), Ch.~4 (p.~62), Ch.~5 (p.~88); LF Chs.~1, 2, 7. \\
\addlinespace[3pt]
\textbf{REQ-ACT-03} & \textbf{Six Functional Facilitation Boards:} Iteration (IFB), Consumption (CFB), Production (PFB), Employment (EFB), Household (HFB), and Updating (UFB) facilitation boards performing purely informational/mediating tasks. & \texttt{[ORIGINAL-1991]} & LF Ch.~6 (\texttt{lf6\_1.htm}), Ch.~9 (\texttt{lf9\_2.htm}); PEPE Ch.~4 (p.~62). \\
\addlinespace[3pt]
\textbf{REQ-ACT-04} & \textbf{Communities of Affected Parties (CAPs) \& Environmental Courts:} Airshed/watershed collectives participating in the Pollution Demand-Revealing Mechanism (PDRM) below hard biophysical caps. & \texttt{[INTERP-LATER]} & DEP Ch.~7 (pp.~138--149), Ch.~14 (pp.~261--267); PEP (2023). \\
\addlinespace[3pt]
\textbf{REQ-ACT-05} & \textbf{Stand-Alone Capital Asset Entities \& Society Account:} Double-entry social accounting entities tracking capital production costs, user-right fees, CAP flows, and GDP division between consumption and investment. & \texttt{[INTERP-LATER]} & Sandstr\"om (2020, Chs.~3--5; 2023). \\
\addlinespace[3pt]
\textbf{REQ-DATA-01} & \textbf{Seven Priced Economic Categories:} Private goods, public goods, intermediate inputs, natural resources, labour categories, capital stocks, and regulated pollutants. & \texttt{[ORIGINAL-1991]} + \texttt{[INTERP-LATER]} & PEPE Ch.~5 (pp.~78--88); DEP Ch.~7, Ch.~9; Szczepanczyk (2023). \\
\addlinespace[3pt]
\textbf{REQ-DATA-02} & \textbf{Coarse Planning Taxonomy vs.\ Granular SKUs:} Plan over coarse resource-equivalent product classes; extrapolate specific sizes/colors statistically and resolve at intra-year fulfilment. & \texttt{[ORIGINAL-1991]} + \texttt{[INTERP-LATER]} & PEPE Ch.~4 (p.~64); LF Ch.~6 (\texttt{lf6\_2.htm}); DEP Ch.~8 (pp.~163--167). \\
\addlinespace[3pt]
\textbf{REQ-DATA-03} & \textbf{Public Books of Blueprints \& Qualitative \texttt{Lists}:} Publish all technical input-output recipes $[A,K,R,L]$ and attach qualitative descriptive \texttt{Lists} to quantitative council vectors. & \texttt{[ORIGINAL-1991]} & PEPE Ch.~4 (pp.~61--62), Ch.~5 (p.~84), Ch.~6; PLAC Chs.~11, 17. \\
\addlinespace[3pt]
\textbf{REQ-DATA-04} & \textbf{Protected Private Consumption Sphere:} Personal consumption \emph{within} warranted budget $B_u$ cannot be vetoed (PEPE). Cryptographic invisibility of the bundle is a DPEP extension, not Albert and Hahnel's 1991 design (Section~\ref{sec:consumption-privacy}). & \texttt{[ORIGINAL-1991]} (non-veto) + \texttt{[TECH-EXT]} (invisibility) & PEPE Ch.~3 (p.~41), Ch.~5 (p.~80): non-veto of in-budget requests. LF Ch.~3 subjects consumption to collective criticism. PLAC Chs.~12, 18: anonymous submission against ``busybody'' intrusion, not cryptographic aggregation. \\
\addlinespace[3pt]
\textbf{REQ-PLAN-01} & \textbf{Multi-Horizon \& Collective-First Sequencing:} Long-run (education, environment) and aggregate investment plans precede annual planning; federation collective goods are settled before final private consumption bundles. & \texttt{[ORIGINAL-1991]} + \texttt{[INTERP-LATER]} & LF Ch.~3 (\texttt{lf3\_2.htm}), Ch.~8 (\texttt{lf8\_1.htm}), Ch.~9 (\texttt{lf9\_1.htm}); DEP Chs.~11--15. \\
\addlinespace[3pt]
\textbf{REQ-PLAN-02} & \textbf{Slack Planning \& Hierarchical Intra-Year Clearing:} Include 3--5\% planned excess supply buffer ($SL^E$); cancel routine consumer swaps inside neighbourhood/ward/city federations; trigger UFB replanning only when net shocks exceed slack. & \texttt{[ORIGINAL-1991]} & PEPE Ch.~4 (pp.~66--67), Ch.~6 (p.~113, 122); DEP Ch.~8 (pp.~167--170). \\
\end{longtable}}

% ==============================================================================
% CHAPTER 4: FORMAL ECONOMIC MODEL AND COMPUTATIONAL SPECIFICATION
% ==============================================================================
\chapter{Formal Economic Model, Accounting Identities, and Algorithms}

\chapterepigraph{It is impossible to show that a feasible nonmarket system at least approaches the productivity of the market unless (1) a rather well-developed theoretical model of the nonmarket system is available, and (2) it is demonstrated that a sufficiently productive approximation of the ideal\dots\ system described in the theoretical model is practically possible.}{Allen Buchanan, \emph{Ethics, Efficiency, and the Market} (1985), p.~29, quoted in Albert and Hahnel, \emph{The Political Economy of Participatory Economics} (1991), Ch.~4, p.~58, and Ch.~5, p.~73\protect\footnotemark}
\footnotetext{Allen Buchanan, \emph{Ethics, Efficiency, and the Market} (Totowa, NJ: Rowman and Littlefield, 1985), p.~29; cited in Albert and Hahnel, \emph{Political Economy of Participatory Economics}, Ch.~4, p.~58 n.~3, and Ch.~5, p.~73.}

\section{Ontology of Economic Actors and Priced Objects}

Let the participatory economy $\mathcal{E}$ over annual planning period $T$ consist of the following finite sets of actors and institutional entities:\footnote{Formalised from Albert and Hahnel, \emph{Political Economy of Participatory Economics}, Ch.~4 (pp.~62--68) and Ch.~5 (pp.~77--95); Hahnel, \emph{Democratic Economic Planning}, Ch.~7 (pp.~130--156); Sandstr\"om, \emph{Anarchist Accounting} (2020), Ch.~3; and Participatory Economy Project, ``User Story Mapping'' (2023).}
\begin{itemize}[leftmargin=1.6em,itemsep=2pt]
  \item \textbf{Individuals ($\mathcal{I} = \{1, \dots, N_{\text{ind}}\}$):} Natural persons participating as consumers, workers, and citizens. Each individual $u \in \mathcal{I}$ belongs to a primary neighbourhood consumer council $h(u) \in \mathcal{H}$ and, if working, a primary worker council $j(u) \in \mathcal{J}$ (plus optional secondary council $j'(u)$ for cross-workplace Balanced Job Complex rotation).\footnote{Albert and Hahnel, \emph{Political Economy of Participatory Economics}, Ch.~2, pp.~27--30; Ch.~5, p.~83.}
  \item \textbf{Neighbourhood Consumer Councils ($\mathcal{H} = \{CC_1, \dots, CC_H\}$) \& Nested Consumer Federations ($\mathcal{F}^C = \bigcup_{\ell=1}^L \mathcal{F}^C_\ell$):} A rooted geographic hierarchy of depth $L=5$ (Level 1 = Ward, Level 2 = City/County, Level 3 = State, Level 4 = Region, Level 5 = National Federation).\footnote{Albert and Hahnel, \emph{Political Economy of Participatory Economics}, Ch.~3, pp.~40--42; Ch.~5, pp.~81--82.}
  \item \textbf{Worker Councils ($\mathcal{J} = \{WC_1, \dots, WC_J\}$) \& Dual Worker Federations ($\mathcal{F}^W = \mathcal{F}^{W,\text{ind}} \cup \mathcal{F}^{W,\text{reg}}$):} Democratic workplace collectives federated both by industry ($IF_k \in \mathcal{F}^{W,\text{ind}}$) and by geographic region ($\mathcal{F}^{W,\text{reg}}$).\footnote{Albert and Hahnel, \emph{Political Economy of Participatory Economics}, Ch.~2, pp.~23--25; Ch.~5, pp.~88, 96--97.}
  \item \textbf{Communities of Affected Parties ($\mathcal{C}^{\text{CAP}} = \{CAP_1, \dots, CAP_Z\}$):} Ecological collectives of residents $\text{pop}(CAP_z) \subseteq \mathcal{I}$ impacted by pollutant $z \in \mathcal{Z}$ in a specific airshed or watershed.\footnote{Hahnel, \emph{Democratic Economic Planning}, Ch.~7 (``The Pollution Demand Revealing Mechanism''), pp.~139--148.}
  \item \textbf{Six Facilitation Boards ($\mathcal{B} = \{\text{IFB}, \text{CFB}, \text{PFB}, \text{EFB}, \text{HFB}, \text{UFB}\}$), Stand-Alone Capital Asset Entities ($\mathcal{A}^{\text{cap}}$), and the Society Account ($\mathcal{A}^{\text{soc}}$).}\footnote{Albert and Hahnel, \emph{Looking Forward}, Ch.~6 (\texttt{lf6\_1.htm}); Sandstr\"om, ``An Accounting System for a Libertarian Socialist Economy'' (2023).}
\end{itemize}

\begin{definition}[Master Catalogue of Seven Priced Economic Categories]
Following Albert and Hahnel (1991a, Chs.~5--6), Hahnel (2021, Chs.~7--9), and Szczepanczyk (2023, Sec.~4), the economy tracks supply, demand, and indicative prices across seven disjoint categories indexed by the master catalogue $\Omega$:
\begin{equation}
  \Omega = \mathcal{G}^{\text{priv}} \sqcup \mathcal{G}^{\text{pub}} \sqcup \mathcal{G}^{\text{int}} \sqcup \mathcal{R} \sqcup \mathcal{L} \sqcup \mathcal{K} \sqcup \mathcal{Z}
\end{equation}
with complete indicative price vector at iteration $t$ given by:
\begin{equation}
  P^{(t)} = \left(p^{\text{priv}}, \, p^{\text{pub}}, \, p^{\text{int}}, \, p^{\text{res}}, \, p^{\text{lab}}, \, p^{\text{cap}}, \, p^{\text{pol}}\right) \in \mathbb{R}_+^{|\Omega|}
\end{equation}
\end{definition}

\section{Primal-Dual Reference Model and Sandstr\"om's Society Account}

In Chapter 5 of \emph{The Political Economy of Participatory Economics} (pp.~78--79), Albert and Hahnel introduce a primal-dual linear programming reference model to clarify the meaning of shadow prices before presenting their decentralised iterative procedure.\footnote{Albert and Hahnel, \emph{Political Economy of Participatory Economics}, Ch.~5, pp.~78--79.} Let $x$ be the vector of production activity levels, $A$ be the produced intermediate input matrix, $K$ be the capital stock coefficient matrix (with available stock $k^*$), $R$ be the non-producible natural resource input matrix (with available ecological supply $r^*$), $L$ be the direct labour input matrix (with available labour $l^*$), and $v$ be the row vector of relative social values of net outputs. The economy-wide primal and dual reference problems are:
\begin{align}
  \textbf{Primal:} \quad \max_{x \ge 0} \quad & v(I - A)x \quad \text{subject to} \quad Kx \le k^*, \;\; Rx \le r^*, \;\; Lx \le l^* \label{eq:primal} \\
  \textbf{Dual:} \quad \min_{P_k, P_r, P_l \ge 0} \quad & P_k k^* + P_r r^* + P_l l^* \quad \text{subject to} \quad P_k K + P_r R + P_l L \ge v(I - A) \label{eq:dual}
\end{align}
As Albert and Hahnel explicitly note (PEPE Ch.~5, p.~79), the shadow prices $P_l$ solved from the static dual (\ref{eq:dual}) reflect only the marginal physical productivities of different labour categories, \emph{not} the differential disutilities or desirabilities of work from the workers' point of view; whereas the indicative prices generated by Parecon's iterative council procedure incorporate worker well-being and effort disutility directly.\footnote{Albert and Hahnel, \emph{Political Economy of Participatory Economics}, Ch.~5, p.~79.}

\subsection{Closing the Macro-Accounting Loop (Sandstr\"om 2020/2023)}
Anders Sandstr\"om's \emph{Anarchist Accounting} (2020) and 2023 technical specification formalise how opportunity-cost user fees paid by Worker Councils reconcile with effort-based consumer incomes inside the \textbf{Society Account} ($\mathcal{A}^{\text{soc}}$):\footnote{Sandstr\"om, \emph{Anarchist Accounting} (2020), Ch.~3; Sandstr\"om, ``An Accounting System for a Libertarian Socialist Economy'' (2023), Diagram 1 (Arrows 1--14).}

\begin{enumerate}[leftmargin=1.8em,itemsep=3pt]
  \item \textbf{Social Value of Net Final Output (GDP) at Indicative Prices $P$:}
  \begin{equation}
    Y^{\text{GDP}}(P) = p^{\text{priv}} \cdot D^{\text{priv}} + p^{\text{pub}} \cdot D^{\text{pub}} + p^{\text{cap,new}} \cdot D^{\text{inv}}
  \end{equation}
  In the Society Account, credits consist of all user-right fees charged to Worker Councils for access to society's labour categories ($p^{\text{lab}} \cdot L$), natural resources ($p^{\text{res}} \cdot R$), and productive capital assets held in Stand-Alone Capital Asset Entities ($p^{\text{cap}} \cdot K$), plus Consumer Council depreciation/discount charges on public consumer assets and net WC operating surpluses.
  \item \textbf{Aggregate Consumption Pool, Effort-Weighted Hours ($H^{\text{eff}}$), and Num\'eraire ($w_{\text{base}} \equiv 1.0$):}
  Let $\sigma_{\text{inv}} \in (0, 1)$ be the share of GDP allocated to aggregate investment (DEP Ch.~11),\footnote{Hahnel and Kerkhoff, ``Aggregate Investment Planning,'' Ch.~11 in Hahnel, \emph{Democratic Economic Planning}, pp.~219--237.} $Y^{\text{allow}}$ be the ratified budget for non-working \emph{cash} allowances (Maxim 4 transfers $A_u^{\text{social}}$), $Y^{\text{CAP}}(P) = \sum_{z \in \mathcal{Z}} p_z^{\text{pol}} E_z$ be total CAP pollution compensation (funded by WC emission charges $\sum_j p^{\text{pol}} \cdot e_j$), and $\Delta S^{\text{agg}} = \sum_{u \in \mathcal{I}} \Delta S_u$ be net aggregate household borrowing. Universal public \emph{services} already appear as $p^{\text{pub}}\cdot D^{\text{pub}}$ inside $Y^{\text{GDP}}$; they must \textbf{not} also be subtracted as $Y^{\text{allow}}$, or the consumption pool is double-counted (Table~\ref{tab:macro-sheet}). Effort-weighted hours $H^{\text{eff}} = \sum_{u \in \mathcal{I}^{\text{work}}} h_u \epsilon_u$ convert the worker pool into a unit wage. The identity $w_{\text{base}}\equiv 1$ is a \textbf{num\'eraire normalisation after clearing}, not an independent price:
  \begin{align}
    Y^{\text{work}}(P) &= (1 - \sigma_{\text{inv}}) Y^{\text{GDP}}(P) - Y^{\text{allow}} - Y^{\text{CAP}}(P) - \Delta S^{\text{agg}} \nonumber \\
    w_{\text{base}}(P) &= \frac{Y^{\text{work}}(P)}{\sum_{u \in \mathcal{I}^{\text{work}}} h_u \epsilon_u} \quad\text{(then rescale all prices so } w_{\text{base}}=1\text{)}
  \end{align}
  \item \textbf{Capping Both Mean Workplace Effort Rating and Total Credited Hours (DEP Ch.~6, p.~117):}
  To prevent either collective ``grade inflation'' or phantom-hours inflation (\texttt{INV-4B}), each worker council $WC_j$'s mean effort multiplier and total credited effort-hours are bounded by its \emph{accredited} labour for the plan year---the initial endowment plus ratified net hiring, $\sum_\ell (l^0_{j,\ell} + \Delta l_{j,\ell})$, the same quantity charged in Eq.~(\ref{eq:sc})---and by its social benefit-to-cost ratio:\footnote{Hahnel, \emph{Democratic Economic Planning}, Ch.~6 (``Capping average effort ratings''), p.~117. Using the endowment $l^0_j$ alone would leave newly hired workers uncredited. The Phase~1 default for $\bar{\epsilon}_{\max}$ is a political choice set out in Section~\ref{sec:effort-default}. Vintage-equipment plants are equalised via lower $p^{\text{cap}} \cdot k_j^0$ rental charges. INV-4B does not help if the labour roster itself is falsified---that remains a physical \emph{euthynai} problem.}
  \begin{align}
    \frac{\sum_{u \in WC_j} h_u \epsilon_u}{\sum_{u \in WC_j} h_u} &\le \min\left(\bar{\epsilon}_{\max}, \; \frac{SB_j(P)}{SC_j(P)}\right) \nonumber \\
    \sum_{u \in WC_j} h_u \epsilon_u &\le \left(\sum_{\ell \in \mathcal{L}} \big(l^0_{j,\ell} + \Delta l_{j,\ell}\big)\right) \min\left(\bar{\epsilon}_{\max}, \; \frac{SB_j(P)}{SC_j(P)}\right) \label{eq:effort-cap}
  \end{align}
  Here $l_{j,\ell}$ is measured in accredited hours. Because $\Delta l_j$ is itself part of the proposal, the cap is checked against the ratified roster at plan ratification and against the proposed roster in each round.
  \item \textbf{Individual Warranted Annual Consumption Budget ($B_u$):}
  Combining Sandstr\"om's four income flows (Sandstr\"om 2023, Arrows 1--4), citizen $u$'s warranted budget is:
  \begin{align}
    B_u(P) &= \underbrace{w_{\text{base}}(P) h_u \epsilon_u}_{\text{1: Work Effort}} + \underbrace{A_u^{\text{social}}}_{\text{2: Benefits}} + \underbrace{\sum_{z: u \in CAP_z} \frac{p_z^{\text{pol}} E_z}{|\text{pop}(CAP_z)|}}_{\text{3: CAP Compensation}} + \underbrace{\Delta S_u}_{\text{4: Loan/Need}} \nonumber \\
    \sum_{u \in \mathcal{I}} B_u(P) &= Y^{\text{work}}(P) + Y^{\text{allow}} + Y^{\text{CAP}}(P) + \Delta S^{\text{agg}} = (1 - \sigma_{\text{inv}}) Y^{\text{GDP}}(P) \label{eq:warranted-budget}
  \end{align}
These last two displays are \textbf{identities} once $Y^{\text{allow}}$ is defined as cash transfers only and $A_u^{\text{social}}$ exhausts $Y^{\text{allow}}$. They hold at every price vector, not only at equilibrium. Relations such as $D^{\text{priv}}=S^{\text{priv}}$ are equilibrium conditions, not identities.
\end{enumerate}

\paragraph{Simultaneity of budgets and demand.}
$Y^{\text{GDP}}$ is valued at the \emph{demanded} quantities $D$, so budgets depend on demand and demand depends on budgets: $B_u = B_u(P, D)$. Within an iteration DPEP breaks the loop with a one-round lag, $B_u^{(t)} = B_u(P^{(t)}, D^{(t-1)})$, so councils always optimise against a published budget. The lag is consistent only at a fixed point, where $B_u^{(t)} = B_u(P^{(t)}, D^{(t)})$; whether the lagged process reaches that fixed point is part of open question P2 and is not implied by any convergence result for fixed budgets. The investment share and the investment bundle are not independent choices: at a cleared plan $\sum_u B_u = p^{\text{priv}}D^{\text{priv}} + p^{\text{pub}}D^{\text{pub}}$, so the identities force $\sigma_{\text{inv}}\, Y^{\text{GDP}} = p^{\text{cap,new}} \cdot D^{\text{inv}}$. Either the federations ratify $\sigma_{\text{inv}}$ and $D^{\text{inv}}$ adjusts in value, or they ratify a physical investment bundle and $\sigma_{\text{inv}}$ is implied. DPEP takes the first option (DEP Ch.~11 treats aggregate investment as a prior collective decision), so $D^{\text{inv}}$ is scaled each round to match.

\begin{table}[htbp]
\centering
\caption{Macro balance sheet (identities vs.\ equilibrium conditions)}
\label{tab:macro-sheet}
\scriptsize
\begin{tabularx}{\textwidth}{L{1.55in} Y L{1.15in}}
\toprule
\textbf{Relation} & \textbf{Content} & \textbf{Status} \\
\midrule
$Y^{\text{GDP}} = p^{\text{priv}}D^{\text{priv}} + p^{\text{pub}}D^{\text{pub}} + p^{\text{cap,new}}D^{\text{inv}}$ & Final uses at indicative prices. & Identity (definition). \\
$Y^{\text{work}} = (1-\sigma_{\text{inv}})Y^{\text{GDP}} - Y^{\text{allow}} - Y^{\text{CAP}} - \Delta S^{\text{agg}}$ & Worker remuneration pool. $Y^{\text{allow}}$ is cash $A_u^{\text{social}}$ only. & Identity, given that split. \\
$\sum_u B_u = Y^{\text{work}} + Y^{\text{allow}} + Y^{\text{CAP}} + \Delta S^{\text{agg}}$ & Warranted budgets exhaust the consumption pool. & Identity at every $P$. \\
$p^{\text{lab}}\cdot L + p^{\text{res}}\cdot R + p^{\text{cap}}\cdot K + p^{\text{pol}}\cdot E = $ WC user-right charges & Society Account credits. & Accounting identity. \\
$D_g \le S_g$, $D_r \le r^*$, $D_z \le \min(E_z,E_z^{\max})$ & One-sided feasibility. & Equilibrium / halt condition, not an identity. \\
$p^{\text{lab}} \ge w_{\text{base}}$ & Labour price floor (INV-4B). & Policy constraint; can prevent labour-market clearing in surplus. \\
\bottomrule
\end{tabularx}
\end{table}

\section{Formal Optimisation Problems of Councils and CAPs}

\subsection{Consumer Councils and Nested Federations (PEPE Ch.~5, pp.~79--82, 88)}
Each consumer council $CC_h \in \mathcal{H}$ seeks to maximise its members' total economic well-being $W_h(U_h, \Delta C_h)$---incorporating both immediate preference fulfilment $U_h$ and endogenous human characteristic development $\Delta C_h$---subject to the equitable budget constraint that its total social cost $SC(CC_h, P)$ does not exceed its members' warranted budget $B(CC_h, P) = \sum_{u \in \text{pop}(CC_h)} B_u(P)$:\footnote{Albert and Hahnel, \emph{Political Economy of Participatory Economics}, Ch.~5, pp.~77--80, 88.}
\begin{align}
  \max_{y_h^{\text{priv}}, \, y_h^{\text{pub,loc}} \ge 0} \quad & W_h\Big(U_h(y_h^{\text{priv}}, y^{\text{pub}}), \; \Delta C_h(y_h^{\text{priv}}, y^{\text{pub}})\Big) \\
  \text{subject to} \quad & SC(CC_h, P) \triangleq p^{\text{priv}} \cdot y_h^{\text{priv}} + p^{\text{pub}} \cdot y_h^{\text{pub,loc}} \nonumber \\
  &\qquad\qquad\quad + |\text{pop}(CC_h)| \sum_{f \succ CC_h} \frac{p^{\text{pub}} \cdot y_f^{\text{pub}}}{|\text{pop}(f)|} \le B(CC_h, P) \nonumber
\end{align}
where $f \succ CC_h$ denotes every higher-level ward, city, state, or national consumer federation containing $CC_h$. Albert and Hahnel state (PEPE Ch.~3, p.~41; Ch.~5, p.~80) that for any individual or council whose request satisfies $SC \le B$, other councils may advise but \textbf{cannot veto} the request; veto power applies only to exception requests exceeding warranted allowances ($SC > B$). That is a normative design rule, not a theorem.

\subsection{Communities of Affected Parties (CAPs) (DEP Ch.~7, pp.~139--148)}
For each regulated pollutant $z \in \mathcal{Z}$ with hard biophysical sustainability ceiling $E_z^{\max}$ (established by Long-Run Environmental Planning; DEP Ch.~14),\footnote{Hahnel, \emph{Democratic Economic Planning}, Ch.~7, pp.~139--148, Appendix pp.~159--162, and Ch.~14, pp.~261--267.} given indicative damage price $p_z^{\text{pol}}$, $CAP_z$ chooses the total permissive emission supply $E_z \in [0, E_z^{\max}]$ to maximise member well-being:
\begin{equation}
  \max_{0 \le E_z \le E_z^{\max}} \quad W_{CAP_z}\Big(p_z^{\text{pol}} E_z, \; -D_{CAP_z}(E_z)\Big)
\end{equation}
At an interior optimum below $E_z^{\max}$, $CAP_z$ sets $E_z$ such that the sum of marginal victim damages across $u \in \text{pop}(CAP_z)$ equals $p_z^{\text{pol}}$ (the Lindahl-Pigou efficiency condition).

\subsection{Worker Councils (PEPE Ch.~5, pp.~85--88; DEP Ch.~7, pp.~152--156)}
Each worker council $WC_j \in \mathcal{J}$ begins the planning cycle with initial labour endowment $l_j^0$ ($\sum_{j=1}^J l_j^0 = l^*$) and initial capital stock endowment $k_j^0$ ($\sum_{j=1}^J k_j^0 = k^*$), for which it is automatically charged $p^{\text{lab}} \cdot l_j^0 + p^{\text{cap}} \cdot k_j^0$.\footnote{Albert and Hahnel, \emph{Political Economy of Participatory Economics}, Ch.~5, pp.~85--87.} In iteration $t$, $WC_j$ proposes output vector $\bar{a}_j = (\bar{a}_j^{\text{priv}}, \bar{a}_j^{\text{pub}}, \bar{a}_j^{\text{int}}, \bar{a}_j^{\text{inv}})$, intermediate inputs $\underline{a}_j$, natural resources $r_j$, net labour additions/releases $\Delta l_j$, net capital additions/releases $\Delta k_j$, and pollutant emissions $e_j$ from its feasible technology set $\mathcal{T}_j$. Defining Social Benefit ($SB_j$) and Social Cost ($SC_j$):
\begin{align}
  SB_j(P) &= p^{\text{priv}} \cdot \bar{a}_j^{\text{priv}} + p^{\text{pub}} \cdot \bar{a}_j^{\text{pub}} + p^{\text{int}} \cdot \bar{a}_j^{\text{int}} + p^{\text{cap,new}} \cdot \bar{a}_j^{\text{inv}} \label{eq:sb} \\
  SC_j(P) &= p^{\text{int}} \cdot \underline{a}_j + p^{\text{res}} \cdot r_j + p^{\text{lab}} \cdot (l_j^0 + \Delta l_j) + p^{\text{cap}} \cdot (k_j^0 + \Delta k_j) + p^{\text{pol}} \cdot e_j \label{eq:sc}
\end{align}
Equation (\ref{eq:sc}) correctly charges $l_j^0 + \Delta l_j$ (endowment plus net hiring). The floor $p^{\text{lab}} \ge w_{\text{base}}$ in INV-4B, however, prevents labour user-rights from falling when labour is in surplus. If prices are to be opportunity costs, a surplus of a labour category should lower $p^{\text{lab}}_\ell$ below $w_{\text{base}}$; the floor converts that surplus into \textbf{open underemployment} rather than a market-clearing wage cut. An underemployment rebate (paying idle accredited hours from the Society Account) can restore household budgets, but it does not restore the opportunity-cost reading of $p^{\text{lab}}$, and it is funded by reducing $Y^{\text{work}}$ or $\sigma_{\text{inv}}$. This trade-off is unresolved.

$WC_j$ maximises its workers' fulfilment and development $W_j(U_j, \Delta C_j)$ subject to the \textbf{Socially Non-Abusive Production Constraint} (PEPE Ch.~5, pp.~87--88):
\begin{equation}
  \max_{\pi_j \in \mathcal{T}_j} \;\; W_j\Big(U_j(\pi_j), \; \Delta C_j(\pi_j)\Big) \quad \text{subject to} \quad SB_j(P) - SC_j(P) \ge 0 \label{eq:wc-opt}
\end{equation}
that is, $SB_j(P)/SC_j(P) \ge 1$.

\section{Critical Algorithmic Audit of \texttt{pequod-cljs} and Convergence Rules}

In Chapter 9 of \emph{Democratic Economic Planning} (2021, pp.~173--194) and Szczepanczyk (2023),\footnote{Hahnel, Weisdorf, and Szczepanczyk, ``Computer Simulations of Participatory Planning,'' Ch.~9 in Hahnel, \emph{Democratic Economic Planning}, pp.~173--194; Szczepanczyk, ``Pseudocode and Algorithms'' (2023), pp.~43--54; audited against upstream \texttt{msszczep/pequod-cljs} at commit \texttt{8bd46b5c233f0e48c767ad6d5069b9437a67c38b} (\texttt{src/clj/pequod\_cljs/csvgen.clj}, 836 lines; trimmed copy in \texttt{research/sources/pequod\_cljs\_csvgen\_apache2.clj}, 799 lines).} the authors simulated 40 experiments (\texttt{dep1ex61..100.csv}, generated via \texttt{bin/bigger\_gen.sh}) with $30,000$ Worker Councils, $30,000$ Consumer Councils, and $500$ categories (100 private, 100 public, 100 intermediate, 100 natural resources, 100 labour categories), alongside 50 smaller 3,000-council runs (\texttt{dep1ex01..50.csv}).

\subsection{The Hahnel-Weisdorf-Szczepanczyk Price Rule and CSV Audit}
In \texttt{csvgen.clj} (upstream lines 391--392 and 463--469), let $S_i^{(t)}$ be total supply, $D_i^{(t)}$ be total demand, and $\text{surplus}_i^{(t)} = S_i^{(t)} - D_i^{(t)}$. Defining the two-sided relative imbalance $\rho_i^{(t)}$:
\begin{equation}
  \rho_i^{(t)} = \frac{2 \left| S_i^{(t)} - D_i^{(t)} \right|}{S_i^{(t)} + D_i^{(t)}}, \qquad \tilde{\delta}_i^{(t)} = 1.05 - 0.5^{\rho_i^{(t)}}
\end{equation}
the damped percentage price delta $\delta_i^{(t)}$ and multiplicative price update are:\footnote{Szczepanczyk, ``Pseudocode and Algorithms'' (2023), Eq.~9, p.~48; upstream \texttt{csvgen.clj} lines 391--392 (\texttt{get-deltas}) and 463--469.}
\begin{align}
  \delta_i^{(t)} &= \max\left(0.001, \; \tilde{\delta}_i^{(t)} \cdot \min\left(0.25, \rho_i^{(t-1)}\right)\right) \nonumber \\
  P_i^{(t+1)} &= P_i^{(t)} \left(1 + \text{sign}\big(D_i^{(t)} - S_i^{(t)}\big)\delta_i^{(t)}\right) \label{eq:pequod-rule}
\end{align}
The upstream code computes $\max(0.001, \min(\tilde\delta, \tilde\delta \cdot \min(0.25,\rho^{(t-1)})))$. Because $\tilde\delta \in [0.05, 0.80]$ is positive and $\min(0.25,\cdot) \le 0.25$, the inner $\min$ always returns its second argument, so the simplified form above is identical. (A commented-out variant in the same file, $\min(\tilde\delta, \rho^{(t-1)})$, is the one in which the outer $\min$ would bind.)
Our programmatic audit of the committed CSV files\footnote{Performed by \path{research/analysis/verify_r2.py} in the project repository.} reports the subsets the book actually cites. Across \texttt{dep1ex61--100.csv} ($n=40$), Year 1 reaches the 5\% threshold (\texttt{:green}) in \textbf{11.85 rounds} on average and the 3\% threshold (\texttt{:blue}) in \textbf{19.20 rounds}; Year 2 reaches 5\% in \textbf{6.58 rounds}. Averaging all 99 files other than \texttt{dep1ex100.csv} yields 11.505 Year-1 5\% rounds and must not be quoted as 11.85. File \texttt{dep1ex100.csv} is a committed large-run log; there is no matching \texttt{dep1ex100.clj}. The 50 \texttt{dep2ex*.csv} files have no corresponding \texttt{.clj} generators in the clone; a code comment in \texttt{csvgen.clj} associates a second experiment family with toothaches, but committed \texttt{dep2ex} CSVs do not show the large between-section exponent jump that would mark a toothache run. Because no wall-clock timing log exists in the upstream repository, any per-iteration wall-clock figure ($\sim 2.5\text{ s}$) is an unverified estimate. Hahnel (2021) Ch.~9 reports the Cobb--Douglas simulations; that chapter was consulted only through secondary notes and does not claim a page-level reading of increasing-returns results.

\subsection{Audit of \texttt{pequod-cljs} and the ``Toothache'' Construction}
Direct inspection of \texttt{wcs.clj}, \texttt{ccs.clj}, \texttt{csvgen.clj}, and all 70 \texttt{dep1ex*.clj} generators establishes four structural properties of \texttt{pequod-cljs}. Fifty files (\texttt{dep1ex01--50}) are stored directly; twenty (\texttt{dep1ex51}, \texttt{dep1ex60--78}) are Git LFS objects of about 160~MB each, fetched from the repository's media endpoint and checked against the SHA-256 in their pointers. Together they hold 750,000 Worker Council records.\footnote{Checked by \path{research/analysis/verify_pequod_datasets.py} in the project repository.}
\begin{enumerate}[leftmargin=1.8em,itemsep=2pt]
  \item \textbf{Unconstrained Closed-Form Cobb-Douglas Optimisation:} Each $WC_j$ maximises $\Pi(x,e) = \lambda a e^c \prod_{m=1}^n x_m^{b_m} - \sum_{m=1}^n p_m x_m - s e^k$ via symbolic Mathematica first-order conditions (\texttt{solution-1..8}, upstream lines 221--401) without enforcing $SB_j / SC_j \ge 1$ or labour stock conservation $\sum_j \Delta l_j = 0$.
  \item \textbf{``Increasing returns'' (\texttt{create-toothaches}) and unboundedness:} In \texttt{wcs.clj} (lines 31--70), each $WC_j$ has $n \in \{3, \dots, 8\}$ inputs with baseline exponents $\sum_{m=1}^n b_m \in [0.75, 0.85]$, effort elasticity $c \in [0.05, 0.10]$, and effort disutility exponent $k \in [3, 4]$ (\texttt{:du}). Along the ray $x_m(t)=t$, $e(t)=t^{1/k}$, revenue scales as $t^{\sum b_m + c/k}$ against linear cost, so $\Pi$ is \textbf{bounded with an interior maximum along that ray iff} $\sum b_m + c/k < 1$. (Joint concavity of $\Pi$ on $(x,e)$ is a stronger, distinct condition and is not claimed.) \texttt{create-toothaches} in \texttt{csvgen.clj} selects 20\% of councils and adds $0.1$ to each $b_m$. Across all 750,000 records the baseline index $\sum b_m + c/k$ lies in $[0.7696, 0.8781]$; with the toothache increment it would lie in $[1.0696, 1.6652] > 1$ for every record, and hence for every selected council, so the FOC is a saddle along the scale ray and the formula predicts \textbf{downward-sloping supply}. That prediction is \textbf{not observed} in committed CSVs: \texttt{-main} currently sets \texttt{toothaches? true}, yet the largest Year-1/Year-2 \texttt{exponent-sum} field change found in \texttt{dep1ex*.csv} and \texttt{dep2ex*.csv} is empty (those rows carry no numeric exponent), and the largest first-green-round change between year sections is 0.03 on the files inspected.
  \item \textbf{Exogenous Consumer Budgets:} In \texttt{consume} (upstream lines 188--204), each Consumer Council's budget \texttt{(:income cc)} is fixed at $5,000$ across all iterations.
  \item \textbf{Two-Sided Resource Tolerance Without Pollutants:} Natural resources are given a fixed supply of $1,000$ (upstream line 168) and priced within the two-sided $\tau = 5\%$ tolerance (permitting up to $\sim 5.1\%$ overshoot), without CAP pollutant pricing or strict one-sided ceiling enforcement.
\end{enumerate}

\subsection{Why General Convergence Cannot Be Guaranteed by Heuristic T\^atonnement}
Simply adding heuristic proximal regularisation and trust regions to Eq.~(\ref{eq:pequod-rule}) cannot guarantee convergence under Leontief complementarities and non-convexities, on both theoretical and empirical grounds:
\begin{itemize}[leftmargin=1.6em,itemsep=2pt]
  \item \textbf{Scarf's (1960) Leontief Instability \& Trust-Region Stalling:} By the Sonnenschein--Mantel--Debreu theorem and Herbert Scarf's (1960) counterexample,\footnote{Herbert Scarf, ``Some Examples of Global Instability of the Competitive Equilibrium,'' \emph{International Economic Review} 1, no. 3 (1960): 157--172; Hugo Sonnenschein, ``Do Walras' Identity and Continuity Characterize the Class of Community Excess Demand Functions?'' \emph{Journal of Economic Theory} 6, no. 4 (1973): 345--354.} price t\^atonnement can cycle indefinitely when income effects violate gross substitutability. Running Eq.~(\ref{eq:pequod-rule}) on Scarf's 3-good Leontief economy\footnote{Run with \path{research/analysis/verify_economics.py} in the project repository.} confirms that over 3,000 rounds, $\max_i \rho_i^{(t)}$ never drops below \textbf{10.54\%} ($\beta \in \{\text{None}, 0.25, 0.10, 0.02\}$), while a decaying trust region ($\beta^{(t)} = 0.25 \cdot 0.99^t$) freezes proposals at a \textbf{56.60\% imbalance}---shrinking trust regions stall, they do not clear the cycle.
  \item \textbf{Indivisible Plants vs.\ One-Good Shapley--Folkman Smoothing:} For a single indivisible facility with $S(p) = 100 \cdot \mathbb{I}_{p \ge 10}$ and demand $D(p) = 70 - p$, no clearing price exists: the theoretical minimum imbalance is $\rho = 50.0\%$, and over 500 rounds Eq.~(\ref{eq:pequod-rule}) never falls below \textbf{50.37\%}. Disaggregating that capacity across $K=20$ plants of size $100/K$ with staggered thresholds $8.5+3i/(K-1)$ yields a \emph{minimum} imbalance of \textbf{0.40\%} across 500 rounds. That is the best round, not a stable gap. The last-100-round maximum at $a=70$ is \textbf{8.4\%}; a sweep $a\in[60,80]$ in steps of $0.5$ has median last-100-round max \textbf{6.0\%}, worst \textbf{9.5\%}, and only 24\% of cases stay $\le 5\%$ in the last 100 rounds.\footnote{Computed by \path{research/analysis/verify_r2.py} in the project repository.} Identical thresholds at 10 give no smoothing ($\min\rho\approx 50\%$). Starr (1969) bounds the Hausdorff gap between the sum of non-convex sets and its convex hull; this one-good staggered-threshold test does not instantiate that theorem for interdependent Leontief markets.\footnote{Ross M.\ Starr, ``Quasi-Equilibria in Markets with Non-Convex Preferences,'' \emph{Econometrica} 37, no. 1 (1969): 25--38.}
\end{itemize}

\subsection{Research-Candidate Solvers (Algorithms 3a--3d)}
\label{sec:alg3}
DPEP separates four candidate coordination algorithms to be benchmarked in Phase 1. None is a production commitment. Table~\ref{tab:method-families} places them among the method families under consideration. All four share one definition of success.

\begin{definition}[Feasible, cleared, and socially acceptable plans]\label{def:success}
Let $u^{\text{out}} = \sum_g p_g \max(0, S_g - D_g) / \sum_g p_g S_g$ be the value-weighted \emph{unused-output share}, $u^{\text{need}} = \sum_g p_g \max(0, D_g - S_g) / \sum_g p_g D_g$ the \emph{unmet-demand share}, $u^{\text{lab}}$ and $u^{\text{cap}}$ the shares of accredited labour hours and capital services left idle, and $\Delta W$ the relative welfare or need-satisfaction gap to a reference oracle (a central LP or MILP on the same instance, Table~\ref{tab:method-families}) where one exists.
\begin{enumerate}[leftmargin=1.6em,itemsep=1pt]
  \item A plan is \textbf{feasible} if $D_g \le S_g$ for every good and every hard ecological or resource cap holds ($D_r \le r^*$, $D_z \le E_z^{\max}$). The one-sided rule belongs here: caps must never be exceeded, and slack against a cap is not waste.
  \item A plan is \textbf{cleared} if it is feasible and, for produced goods, labour and capital, $\max_i \rho_i \le \tau$, $u^{\text{out}} \le \bar u^{\text{out}}$, $u^{\text{lab}} \le \bar u^{\text{lab}}$ and $u^{\text{cap}} \le \bar u^{\text{cap}}$, net of the planned slack buffer $SL^E$.
  \item A plan is \textbf{socially acceptable} if it is cleared, $u^{\text{need}} \le \bar u^{\text{need}}$ for priority (Maxim~4) categories, and $\Delta W \le \bar{\Delta}$ where an oracle exists. These thresholds are ratified politically, not derived.
\end{enumerate}
Proposed defaults for Phase~1 (for author ratification): $\tau = 5\%$, $\bar u^{\text{out}} = \bar u^{\text{lab}} = \bar u^{\text{cap}} = 5\%$ beyond planned slack, $\bar u^{\text{need}} = 2\%$, $\bar\Delta = 5\%$ median and 15\% at the 90th percentile.
\end{definition}

The distinction matters. In the single-plant case of Section~4.4.3, Eq.~(\ref{eq:pequod-rule}) reaches one-sided feasibility at round~3 with $S=100$ and $D=58.48$: 42\% of output is unused, and the case has no clearing price. Under the definition above the run is reported as \emph{feasible but not cleared} in every one of 500 rounds. With $K=20$ staggered plants the first feasible round (round~3) still leaves 8.4\% unused, the first cleared round is round~5, and 57\% of rounds are cleared. Scarf's cycle is never even feasible: some good is over-demanded in every one of 3,000 rounds.\footnote{Computed by \path{research/analysis/verify_r3.py} in the project repository.}

\paragraph{Algorithm 3a --- Proximal damped t\^atonnement (heuristic).}
\begin{enumerate}[leftmargin=1.6em,itemsep=1pt]
  \item \textbf{Input:} $P^{(0)}$, $\mu>0$, thresholds of Definition~\ref{def:success}, $T_{\max}$, caps $r^*,E_z^{\max}$.
  \item For $t=1,\dots,T_{\max}$: each council solves $\max_{\pi\in\mathcal{T}_c} W_c(\pi;P^{(t)}) - \tfrac{\mu}{2}\|\pi-\pi^{(t-1)}\|_2^2$ subject to $SB/SC\ge 1$ or $SC\le B$ as appropriate; IFB updates $P$ by Eq.~(\ref{eq:pequod-rule}) and rescales so $w_{\text{base}}=1$.
  \item \textbf{Halt with success} when the plan is \emph{cleared}. If $T_{\max}$ is reached, publish the best feasible plan with its $u^{\text{out}}$, $u^{\text{lab}}$, $u^{\text{cap}}$ and $u^{\text{need}}$, labelled \emph{feasible, not cleared}, or report infeasibility. No rate is claimed.
\end{enumerate}

\paragraph{Algorithm 3b --- Proximal Jacobian ADMM (Deng, Lai, Peng \& Yin 2017, Algorithm~4; fixed-budget convex case only).}
Deng et al.\ treat $\min \sum_{i=1}^N f_i(x_i)$ subject to $\sum_i A_i x_i = c$ with each $f_i$ closed proper convex. DPEP's coupling constraints are inequalities $\sum_c A_c x_c \le b$ (demand within supply and caps), so a non-negative slack block $s \ge 0$ with $f_s \equiv 0$ is added and the constraint becomes $\sum_c A_c x_c + s = b$; the number of blocks is $N = N_{\text{active}} + 1$. With $f_c = -W_c + \iota_{\mathcal{T}_c}$:
\begin{enumerate}[leftmargin=1.6em,itemsep=1pt]
  \item \textbf{Input:} closed convex $\mathcal{T}_c$, \emph{fixed} $B_u$, penalty $\rho>0$, dual damping $\gamma\in(0,2)$, proximal matrices $P_i \succ \rho\big(\tfrac{N}{2-\gamma}-1\big)A_i^\top A_i$ (Deng et al.\ eq.~2.11; for $P_i=\tau_i I$ this is $\tau_i > \rho(\tfrac{N}{2-\gamma}-1)\|A_i\|^2$, and for the prox-linear choice $P_i = \tau_i I - \rho A_i^\top A_i$ it is $\tau_i > \tfrac{\rho N}{2-\gamma}\|A_i\|^2$).
  \item In parallel for every block $i$ (councils and slack), with $r^k = \sum_j A_j x_j^k - b$:
  \[
  x_i^{k+1} = \arg\min_{x_i} \; f_i(x_i) + \tfrac{\rho}{2}\Big\| A_i x_i + \textstyle\sum_{j\ne i} A_j x_j^k - b - \lambda^k/\rho \Big\|^2 + \tfrac12 \|x_i - x_i^k\|_{P_i}^2 .
  \]
  \item Dual step: $\lambda^{k+1} = \lambda^k - \gamma\rho\,(\sum_i A_i x_i^{k+1} - b)$. With the slack block the dual needs no projection; indicative prices are $-\lambda$.
  \item \textbf{Halt} when the plan is cleared (Definition~\ref{def:success}) and $\|u^k - u^{k+1}\|$ is below tolerance; otherwise fail at $T_{\max}$.
\end{enumerate}
Deng et al.\ prove convergence and an $o(1/k)$ rate on the successive-iterate quantity $\|u^k-u^{k+1}\|_{G'}^2$ (primal and dual together), not a rate to optimal welfare. The theorem does not cover endogenous budgets, $\varepsilon$-caps that depend on prices, renormalising duals so that $w_{\text{base}}\equiv1$ mid-run, or non-convex $\mathcal{T}_c$ (convergence results for ADMM on non-convex problems, such as Hong, Luo and Razaviyayn 2016, need structural conditions that DPEP's lumpy technologies do not meet). Writing the objective as $\sum_c W_c$ assumes cardinal, interpersonally comparable well-being. Because $N$ includes every active council, $\tau_i$ grows with $N_{\text{active}}$; whether that makes convergence too slow at $N=10^6$ is untested.

A simpler-looking iteration, $x_c^{t+1} = \arg\min -W_c(x) + \langle\lambda^t, A_c x\rangle + \tfrac12\|x - x_c^t\|_{P_c}^2$ with a projected dual ascent, is easily mistaken for ADMM. It is an \textbf{Arrow--Hurwicz-type primal--dual proximal iteration}: it omits the augmented residual term, and the Deng et al.\ theorem does not apply to it. In a two-council test with linear welfare and a shared capacity of~1, it oscillates indefinitely (residual $\approx 1.0$, price swinging between 0 and 2 over 20,000 rounds), while Algorithm~4 with a slack block converges to residual~0 and price~1; both converge when welfare is strictly concave.\footnote{Computed by \path{research/analysis/reference_tests_r3.py} and \path{research/analysis/verify_r3.py} in the project repository.}

\paragraph{Algorithm 3d --- Damped budget fixed point around a fixed-budget inner solve (candidate for endogenous income).}
Endogenous income makes full DPEP an equilibrium problem: budgets $B_u(P,D)$ (Eq.~\ref{eq:warranted-budget}) depend on the prices and quantities they help determine, and the effort cap couples incomes to $SB/SC$. The natural formulation is a mixed complementarity problem (MCP) in $(x, P, B)$. Phase~1 needs a solver it can actually run, so the proposed candidate is an outer loop:
\begin{enumerate}[leftmargin=1.6em,itemsep=1pt]
  \item Hold budgets $B^{k}$ fixed and solve the inner fixed-budget problem with Algorithm~3b (convex instances) or 3c (lumpy instances), giving $(x^k, P^k)$.
  \item Recompute warranted budgets $\hat B^{k} = B(P^k, D^k)$ from the Society Account identities and update $B^{k+1} = (1-\alpha)B^k + \alpha \hat B^k$ with damping $\alpha \in (0,1]$.
  \item Stop when $\|B^{k+1}-B^k\|/\|B^k\| \le 10^{-3}$ and the inner plan is cleared; report cycling or divergence otherwise.
\end{enumerate}
\emph{Assumptions and limits.} Each inner solve inherits 3b's or 3c's assumptions. No convergence theorem is claimed for the outer loop; a fixed point need not be unique, and the loop can cycle when income effects are strong (the same mechanism that defeats t\^atonnement on Scarf's example). As a cross-check on small instances, the MCP can be passed to a PATH-type complementarity solver; PATH is proprietary and its free licence is limited to small problems, so licensing terms must be checked before relying on it, and no open-source MCP solver of comparable robustness is assumed. Lemke's method applies only to the linear complementarity case and is used as a diagnostic.

\emph{Interaction with Algorithm 3c.} The 3c objective below is written at the current budgets $B_u$. Under endogenous income a lumpy plant decision changes incomes, and so the demand it is meant to meet; 3c is therefore run as the inner solve of 3d, and binary decisions taken at one budget vector are re-checked after each budget update. Phase~1 tests 3c first with fixed budgets and then inside 3d.

\paragraph{Algorithm 3c --- Two-tier federation MILP + continuous ADMM (undefined until the following are fixed).}
\begin{enumerate}[leftmargin=1.6em,itemsep=1pt]
  \item \textbf{What councils reveal:} each $WC_j$ posts a discrete menu of plant/shift options $z_{j,k}\in\{0,1\}$ with publicly listed input--output rows (REQ-DATA-03), not its full private technology set $\mathcal{T}_j$. Continuous operating bands around each chosen $z$ remain local.
  \item \textbf{Objective (provisional):} lexicographic (1) one-sided ecological and resource feasibility, (2) minimise unmet consumer demand at current $B_u$, (3) minimise unused output and idle labour and capital (Definition~\ref{def:success}), (4) minimise a public lumpy-cost index. This objective is a federation \emph{proposal}, not a social-welfare theorem. Under endogenous income 3c runs inside Algorithm~3d.
  \item \textbf{Decision rights:} binary $z$ are ratified by the relevant industry and regional federations plus affected CAPs; a central \texttt{HiGHS} MILP is an aid, not an authority (Kuch 2024).
  \item \textbf{Infeasibility:} if the MILP is infeasible, return the IIS (irreducible infeasible subsystem), raise the relevant cap or investment debate, and \textbf{do not} silently relax $E_z^{\max}$. Halt and publish the IIS after $T_{\max}$ unsuccessful repairs.
\end{enumerate}

\begin{table}[htbp]
\centering
\caption{Ten method families by problem class (research map, not a ranking)}
\label{tab:method-families}
\scriptsize
\begin{tabularx}{\textwidth}{L{1.35in} L{1.25in} Y}
\toprule
\textbf{Family} & \textbf{Problem class} & \textbf{DPEP status} \\
\midrule
Unconstrained Cobb--Douglas FOC & Smooth convex, fixed income & \texttt{pequod-cljs} baseline; rejected for production. \\
Damped t\^atonnement (3a) & Heuristic, any $\mathcal{T}_c$ & Candidate; fails Scarf/single-plant tests. \\
Proximal Jacobian ADMM (3b) & Convex, \emph{fixed} budgets & Candidate; Deng et al.\ Algorithm~4 with slack block; successive-iterate rate only. \\
Arrow--Hurwicz primal--dual & Convex, fixed budgets & Resembles 3b but lacks the augmented term; oscillates on linear welfare; not used. \\
Fisher / Eisenberg--Gale & Linear utilities, fixed budgets & Fits fixed $B_u$; not endogenous income. \\
Negishi welfare weights & Convex equilibrium & Requires comparable $W_c$; open normatively. \\
Variational inequality / MCP & Endogenous income + prices & Leading formulation. Phase-1 candidate: Algorithm~3d (damped outer loop); PATH-type solver as small-instance cross-check, subject to licence. \\
Federation MILP (3c) & Discrete lumpy $z\in\{0,1\}$ & Candidate; needs ratification and IIS exit. \\
Scarf / Lemke path-following & Linear complementarity & Diagnostic for cycles; not a planning procedure. \\
Central LP (Cockshott-style) & Known $A,K,R,L$ & Rejected as workplace authority; usable as Phase-1 oracle. \\
\bottomrule
\end{tabularx}
\end{table}

\section{Protocol State Machine and Canonical CBOR Schemas}

\begin{figure}[htbp]
\centering
\begin{tikzpicture}[
  node distance=0.42in and 0.45in,
  state/.style={rectangle, rounded corners=3pt, draw=NavyPrimary, fill=SoftBg, thick, minimum width=1.65in, minimum height=0.42in, align=center, font=\scriptsize\bfseries},
  loopstate/.style={rectangle, rounded corners=3pt, draw=CrimsonAccent, fill= CrimsonAccent!8!white, thick, minimum width=1.65in, minimum height=0.42in, align=center, font=\scriptsize\bfseries},
  arrow/.style={-{Stealth[length=2.2mm]}, thick, color=SlateDark}
]
  \node[state] (s0) {S0: Pre-Plan Parameters\\[-1pt]\tiny ($l^*, r^*, E^{\max}, \sigma_{\text{inv}}, P^{(0)}$)};
  \node[state, right=of s0] (s1) {S1: Collective Goods Phase\\[-1pt]\tiny (Federation Votes $y_f^{\text{pub}}$)};
  \node[loopstate, below=of s0] (s2) {S2: Round Open ($t=1..R$)\\[-1pt]\tiny (Local Council Solvers $\pi_c^{(t)}$)};
  \node[loopstate, right=of s2] (s3) {S3: Two-Root Aggregation\\[-1pt]\tiny (BLAKE3 $H_{\text{prop}}, H_{\text{tot}}$)};
  \node[loopstate, below=of s3] (s4) {S4: WASM IFB Price Update\\[-1pt]\tiny (Check $\max_i \rho_i^{(t)} \le \tau$ \& Caps)};
  \node[state, left=of s4] (s5) {S5: Late-Round Synthesis\\[-1pt]\tiny (Optional 5-Plan Vote, PEPE p.71)};
  \node[state, below=of s5] (s6) {S6: Plan Ratification\\[-1pt]\tiny (Signed Commitments \& BJC)};
  \node[state, right=of s6] (s7) {S7: Intra-Year Clearing\\[-1pt]\tiny (Ward Slack Cancellation)};
  \node[state, below=of s6] (s9) {S9: Year-End \emph{Euthynai} Audit\\[-1pt]\tiny (Society \& CAP Reconciliation)};
  \node[state, right=of s9] (s8) {S8: Scoped Replanning\\[-1pt]\tiny (UFB Sub-Graph Adjustment)};

  \draw[arrow] (s0) -- (s1);
  \draw[arrow] (s1.south) -- ++(0,-0.15in) -| (s2.north);
  \draw[arrow] (s2) -- (s3);
  \draw[arrow] (s3) -- (s4);
  \draw[arrow, CrimsonAccent] (s4.east) -- ++(0.22in,0) |- node[pos=0.25, right, font=\tiny, align=left]{$\rho_i > \tau$\\$t \leftarrow t+1$} (s3.east);
  \draw[arrow, CrimsonAccent] (s4.north west) -- (s2.south east);
  \draw[arrow] (s4) -- node[above, font=\tiny]{$t \ge T_{\text{switch}}$} (s5);
  \draw[arrow] (s5) -- (s6);
  \draw[arrow] (s4.south) -- ++(0,-0.14in) -| node[pos=0.25, above, font=\tiny]{$\rho_i \le \tau$} (s6.north);
  \draw[arrow] (s6) -- (s7);
  \draw[arrow] (s7) -- node[right, font=\tiny]{Shock $> SL^E$} (s8);
  \draw[arrow] (s8.west) -- ++(-0.15in,0) |- (s7.south);
  \draw[arrow] (s7.west) -- (s9.north east);
\end{tikzpicture}
\caption{Complete 10-State DPEP Annual Planning and Intra-Year Operational State Machine}
\label{fig:state-machine}
\end{figure}

To prevent IEEE-754 floating-point non-determinism across CPU architectures (x86\_64, ARMv8 smartphones, and WebAssembly), all quantities and prices in DPEP are encoded as deterministic 64-bit fixed-point integers (\texttt{FixedDecimal}, default scale $10^{-6}$ micro-units):

\noindent\textbf{Canonical wire encoding.}\label{sec:fixed-point-bounds} Listings~\ref{lst:cbor-idl} and~\ref{lst:cbor-idl-2} are an IDL sketch (proto3-like field names for readability). It is \textbf{not} sent as protobuf. All signed payloads use \textbf{RFC 8949 \S 4.2 Canonical CBOR} (CDE) with a \textbf{fixed global scale of $10^{-6}$} (\texttt{int64} micro-units). Multiplication uses 128-bit intermediate products, round-half-to-even, saturating on overflow (the round fails closed rather than wrapping). Division by zero is a fault. The map $0.5^{\rho_i}$ is a degree-6 minimax polynomial on $\rho\in[0,2]$ with published coefficients; outside that interval the update is clipped. Protocol checks verify the rounded inequalities $SB/SC\ge 1$ and $SC\le B$ after this arithmetic. They do not prove that a council maximised a truthful $W_c$, and they do not exist as KKT certificates for MILPs.

\begin{lstlisting}[language=Java,float=tbp,captionpos=t,caption={IDL sketch of DPEP item vectors and consumer-council proposals. Signed encoding is RFC 8949 Canonical CBOR, not protobuf.},label=lst:cbor-idl]
// IDL sketch only --- not the wire format
package dpep.v1;

message SparseItemVector {
  repeated uint32 item_ids = 1;          // Strictly ascending catalog indices in Omega
  repeated int64 micro_quantities = 2;   // Strictly positive quantities (> 0) in 10^-6 micro-units
}

message QualitativeAnnotation {          // Implements PEPE Ch. 6 'Lists' with GDPR redactability
  uint32 item_id = 1;
  string summary_code = 2;
  string narrative_text = 3;
  bytes redactable_pointer_cid = 4;      // Mutable pointer supporting tombstone erasure
}

message ConsumerCouncilProposal {
  string domain_tag = 1;                 // Fixed: "DPEP-WD-v0-CC-PROPOSAL"
  uint32 plan_year = 2;
  uint32 round_number = 3;
  bytes council_id = 4;
  uint32 accredited_member_count = 5;
  SparseItemVector private_demand = 6;   // Coarse-category totals (HRU x s_g) from K=3 Prio3 aggregators
  SparseItemVector local_public_demand = 7;
  int64 total_social_cost_micro = 8;
  int64 warranted_budget_micro = 9;
  bool requests_exception = 10;
  repeated QualitativeAnnotation qualitative_list = 11;
  bytes vdaf_transcript_root = 12;       // Commits Prio3 aggregate-share commitments
  bytes signer_bitmap = 13;
  repeated bytes council_hybrid_signatures = 14; // k-of-n Ed25519+ML-DSA-65 multi-signature
}
\end{lstlisting}

\begin{lstlisting}[language=Java,float=tbp,captionpos=t,caption={IDL sketch of worker-council proposals and IFB round checkpoints (continuation of Listing~\ref{lst:cbor-idl}).},label=lst:cbor-idl-2]
message WorkerCouncilProposal {
  string domain_tag = 1;                 // Fixed: "DPEP-WD-v0-WC-PROPOSAL"
  uint32 plan_year = 2;
  uint32 round_number = 3;
  bytes council_id = 4;
  uint32 industry_federation_id = 5;
  SparseItemVector proposed_outputs = 6;
  SparseItemVector intermediate_inputs = 7;
  SparseItemVector natural_resources = 8;
  SparseItemVector total_labor_requested = 9;
  SparseItemVector total_capital_requested = 10;
  SparseItemVector pollutant_emissions = 11;
  int64 effective_effort_hours_micro = 12; // INV-4B: sum(h_u*eps_u) <= (l0+dl)*min(eps_max,SB/SC)
  int64 social_benefit_micro = 13;
  int64 social_cost_micro = 14;
  int64 sb_sc_ratio_micro = 15;
  bool requests_appeal = 16;               // Subject to soft-budget cap at IF_k
  repeated QualitativeAnnotation qualitative_list = 17;
  bytes signer_bitmap = 18;
  repeated bytes council_hybrid_signatures = 19;
}

message RoundAggregationAndPrices {
  string domain_tag = 1;                   // Fixed: "DPEP-WD-v0-ROUND-CHECKPOINT"
  uint32 plan_year = 2;
  uint32 round_number = 3;
  repeated int64 indicative_prices_micro = 4; // Normalised so w_base_micro == 1_000_000
  repeated int64 total_supply_micro = 5;
  repeated int64 total_demand_micro = 6;
  repeated int64 planned_slack_micro = 7;
  repeated int64 one_sided_shortage_micro = 8; // max(0, D_i - S_i)
  int64 base_hourly_remuneration_micro = 9;    // Fixed numeraire: 1_000_000
  bool strict_one_sided_feasible = 10;
  bytes proposals_merkle_root = 11;            // 32-B BLAKE3 root over sorted council proposals (512-B proof)
  bytes category_totals_merkle_root = 12;      // 32-B BLAKE3 root over category sums
  bytes council_participation_bitmap = 13;     // 2 bits/council: 01=Active, 10=Partitioned, 00=Excluded
  bytes deterministic_wasm_code_hash = 14;
  bytes validator_quorum_bitmap = 15;
  repeated bytes validator_hybrid_signatures = 16;
}
\end{lstlisting}

% ==============================================================================
% CHAPTER 5: DISTRIBUTED SYSTEMS ARCHITECTURE
% ==============================================================================
\chapter{Distributed Systems Architecture, Consistency, and Low-Power Participation}

\chapterepigraph{Planning is computerized. All councils can access data banks of all facilitation boards and all other councils. The only thing demanding about flexible information access is that the computer system not only automatically accept but also sum and compare proposals to ascertain the status of any good.}{Michael Albert and Robin Hahnel, \emph{Looking Forward} (1991), Ch.~6, ``Preparing First Proposals''\protect\footnotemark}
\footnotetext{Albert and Hahnel, \emph{Looking Forward}, Ch.~6 (\texttt{lf6\_2.htm}).}

\section{Why Participatory Economics Is Not a High-TPS Blockchain Problem}

When we examine the computational structure formalised in Chapter 4, three architectural truths emerge:
\begin{enumerate}[leftmargin=1.8em,itemsep=2pt]
  \item \textbf{Annual Planning Is an Epoch-Synchronised Bulk-Synchronous Parallel (BSP) Process:} Millions of participants and tens of thousands of councils draft proposals locally, aggregate them up nested geographic and industrial federation trees, and compute updated indicative prices $P^{(t+1)}$ once per planning round (10--20 rounds per year). Global state finality is required \textbf{only at iteration boundaries} and at plan ratification.
  \item \textbf{Intra-Year Execution Is Hierarchically Scoped:} As Albert and Hahnel (1991a, Ch.~4, pp.~66--67) specify, routine intra-year consumer preference shifts cancel out at the neighbourhood, ward, or city federation level within planned slack buffers ($SL^E \approx 3\text{--}5\%$) without altering aggregate production plans.
  \item \textbf{Participation Must Work Offline and on Resource-Constrained Devices:} Workers and consumers in rural collectives or intermittently connected networks must be able to deliberate, optimize local proposals, and sign ballots offline.\footnote{Martin Kleppmann, Adam Wiggins, Peter van Hardenberg, and Mark McGranaghan, ``Local-First Software: You Own Your Data, in spite of the Cloud,'' \emph{Proceedings of the 2019 ACM SIGPLAN International Symposium on New Ideas, New Paradigms, and Reflections on Programming and Software (Onward! '19)}, pp.~154--178, \url{https://doi.org/10.1145/3359591.3359737}.}
\end{enumerate}

\section{The Four-Layer Hybrid Architecture (Architecture D)}

\begin{figure}[ht!]
\centering
\begin{tikzpicture}[
  node distance=0.34in,
  layer/.style={rectangle, rounded corners=4pt, draw=NavyPrimary, line width=0.9pt, text width=4.52in, align=left, inner xsep=9pt, inner ysep=8.5pt, font=\scriptsize},
  arrow/.style={{Stealth[length=2.4mm]}-{Stealth[length=2.4mm]}, line width=1.2pt, color=CrimsonAccent}
]
  \node[layer, fill=SoftBg] (l4) {
    {\footnotesize\textbf{\textcolor{CrimsonAccent}{LAYER 4: HUMAN DELIBERATION, ACCESSIBILITY \& LOCAL SOLVER EDGE}}}\par\vspace{3.5pt}
    $\bullet$ Offline-First Mobile/Desktop PWA \& Native Clients (\texttt{Role-LP}) + Embedded \texttt{Clarabel.rs} / \texttt{HiGHS} WASM Solver\par\vspace{2pt}
    $\bullet$ ``What-If Collective Consumption Explorer, Qualitative \texttt{List} Browser \& Screen-Reader/Voice Interface\par\vspace{1pt}
  };
  \node[layer, fill=CodeBg, below=of l4] (l3) {
    {\footnotesize\textbf{\textcolor{NavyPrimary}{LAYER 3: LOCAL-FIRST COUNCIL \& HOUSEHOLD STATE LAYER (ASYNC / OFFLINE)}}}\par\vspace{3.5pt}
    $\bullet$ Encrypted SQLite (SQLCipher) Vector Ledger + Loro/Automerge Delta-CRDT Council Workspaces (\texttt{Role-CC})\par\vspace{2pt}
    $\bullet$ \texttt{Prio3} VDAF shares to $K=3$ aggregators + Revocable Capability Certificates (CapCert)\par\vspace{1pt}
  };
  \node[layer, fill=SoftBg, below=of l3] (l2) {
    {\footnotesize\textbf{\textcolor{TealAccent}{LAYER 2: FEDERATED AGGREGATION \& CONTENT-ADDRESSED DATA PLANE (P2P / QUIC)}}}\par\vspace{3.5pt}
    $\bullet$ \texttt{iroh} / \texttt{quinn} (QUIC) P2P Transport with Delay-Tolerant Store-and-Forward Gateways (\texttt{Role-GW})\par\vspace{2pt}
    $\bullet$ \textbf{Two-Root BLAKE3 Aggregation} (\texttt{Role-PCW}) + Verified Streaming of Qualitative \texttt{List} Blobs (\texttt{Role-SR})\par\vspace{1pt}
  };
  \node[layer, fill=NavyPrimary!8!white, below=of l2] (l1) {
    {\footnotesize\textbf{\textcolor{SlateDark}{LAYER 1: SELECTIVE EPOCH-BFT CHECKPOINT \& DETERMINISTIC IFB ENGINE}}}\par\vspace{3.5pt}
    $\bullet$ Rotating Sortition-Selected BFT Committee (\texttt{Role-BFT}, CometBFT / Simplex) Committing 1 Checkpoint / Round\par\vspace{2pt}
    $\bullet$ Sandboxed Deterministic Fixed-Point \texttt{Wasmtime} IFB Price Engine + Sandstr\"om Society Account \& CAP Roots\par\vspace{1pt}
  };

  \draw[arrow] ($(l4.south)+(0,-1.5pt)$) -- ($(l3.north)+(0,1.5pt)$);
  \draw[arrow] ($(l3.south)+(0,-1.5pt)$) -- ($(l2.north)+(0,1.5pt)$);
  \draw[arrow] ($(l2.south)+(0,-1.5pt)$) -- ($(l1.north)+(0,1.5pt)$);
\end{tikzpicture}
\caption{Four-Layer Architecture of the Distributed Participatory Economics Protocol (Architecture D)}
\label{fig:arch-d}
\end{figure}

Figure~\ref{fig:round-seq} traces one planning round message by message across the layers, and Figure~\ref{fig:topology} (Section~\ref{sec:roles}) shows where each node role sits.

\begin{figure}[htbp]
\centering
\begin{tikzpicture}[x=1in,y=1in]
  \foreach \x/\t in {0.47/{Person\\[-1pt]\normalfont(Role-LP, phone)},
                     1.46/{Council\\[-1pt]\normalfont(Role-CC)},
                     2.45/{Prio3 aggregators\\[-1pt]\normalfont(A, B, C)},
                     3.44/{Federation servers\\[-1pt]\normalfont(Role-PCW, Role-SR)},
                     4.43/{BFT committee\\[-1pt]\normalfont(Role-BFT, $n{=}128$)}}{
    \node[dhead, text width=0.8in, minimum height=0.34in] at (\x,0) {\t};
    \draw[dlife] (\x,-0.2) -- (\x,-5.0);
  }
  % 1
  \draw[darr] (4.43,-0.45) -- node[dlab, pos=0.5]{1\; checkpoint $t{-}1$ with prices $P^{(t)}$} (0.47,-0.45);
  % 2
  \node[dbox, text width=0.8in] at (0.47,-0.82) {2\; local solver drafts the request; the person edits it};
  % 3
  \draw[darr, TealAccent] (0.47,-1.22) -- node[dlab, pos=0.5]{3\; report in 3 shares\\(milestone rounds only)} (2.45,-1.22);
  % 4
  \draw[darr, TealAccent] (2.45,-1.62) -- node[dlab, above=1pt, pos=0.5, fill=none]{4\; neighbourhood total ($\ge 30$)} (1.46,-1.62);
  % 5
  \node[dbox, text width=0.8in] at (1.46,-2.02) {5\; council deliberates; $k$ of $n$ members sign $\pi_c^{(t)}$};
  % 6
  \draw[darr] (1.46,-2.42) -- node[dlab, pos=0.5]{6\; $\pi_c^{(t)}$ to primary server + $\ge 2$ relays} (3.44,-2.42);
  % 7
  \draw[darr] (3.44,-2.74) -- node[dlab, pos=0.5]{7\; signed \texttt{SubmissionReceipt}} (1.46,-2.74);
  % 8
  \node[dbox, text width=0.8in] at (3.44,-3.14) {8\; Two-Root reduction, conservation check, WASM price update};
  % 9
  \draw[darr] (3.44,-3.56) -- node[dlab, above=1pt, pos=0.5, fill=none]{9\; roots, bitmap, $P^{(t+1)}$} (4.43,-3.56);
  % 10
  \draw[darr, CrimsonAccent, dashed] (1.46,-3.96) -- node[dlab, pos=0.42, text=CrimsonAccent]{10\; if omitted: receipt within $T_{\text{challenge}}$} (4.43,-3.96);
  % 11
  \node[dbox, text width=0.8in] at (4.43,-4.34) {11\; $\ge 86$ of 128 validators sign checkpoint $t$};
  % 12
  \draw[darr] (4.43,-4.74) -- node[dlab, pos=0.5]{12\; checkpoint $t$ with $P^{(t+1)}$ (coordinator attestation or sync committee)} (0.47,-4.74);
  \node[font=\tiny\itshape, text=SlateDark] at (2.45,-5.14) [text width=4.7in, align=flush center] {Next round: $t \leftarrow t+1$ until the plan clears. Steps 3--4 run only at milestone rounds; otherwise councils revise at council level.};
\end{tikzpicture}
\caption{One planning round as a message sequence. Teal steps carry private household data and run only at milestone rounds; the dashed crimson step is the censorship challenge (Section~\ref{sec:two-root}).}
\label{fig:round-seq}
\end{figure}

\subsection{Two-Root Authenticated Aggregation, Completeness Bitmaps, and Physical Boundaries}
\label{sec:two-root}
Naively embedding an $80\text{ KB}$ dense 10,000-category vector inside every internal node of a depth-16 Sparse Merkle-Sum Tree would cause a single mobile inclusion path to bloat to $16 \times 80\text{ KB} = \mathbf{1.28\text{ MB}}$---and still could not prevent a dishonest aggregator from fabricating off-path sibling subtrees unless auditors recomputed the entire tree from the leaves up. Moreover, a Merkle inclusion proof for council $c$ proves \emph{membership}, not \emph{completeness} (whether dissident council $c'$ was omitted) or \emph{physical truth} (whether reported output physically exists).

A uniform 5.2~kB proposal size is too small: 5,200~B less one 3,373~B ML-DSA-65 signature leaves 1,827~B, about 150 sparse items at 12~B each (4-byte id, 8-byte quantity), while a neighbourhood's summed demand over 3,000 private categories is about 36~kB, and every extra multi-signature signer adds a further 3,373--3,437~B. Table~\ref{tab:proposal-sizes} reworks sizes by council type. Council threshold signatures are not specified: until a standard post-quantum threshold scheme exists (open question P12), a council signs with a $k$-of-$n$ multi-signature, i.e.\ $k$ separate hybrid signatures plus a signer bitmap.

\begin{table}[htbp]
\centering
\caption{Council proposal sizes and round archive at $N=10^6$ (raw canonical CBOR; analytical)}
\label{tab:proposal-sizes}
\scriptsize
\setlength{\tabcolsep}{3pt}
\begin{tabularx}{\textwidth}{L{1.55in} R{0.55in} R{0.55in} Y}
\toprule
\textbf{Council type} & \textbf{1 signer} & \textbf{5 signers} & \textbf{Basis} \\
\midrule
Worker council, 150 items & 5.6 kB & 19.4 kB & 400~B header + $150\times12$~B + hybrid signatures (3,437~B each) \\
Consumer council, 200 coarse categories & 6.2 kB & 20.0 kB & Proposed coarse household baseline (Section~\ref{sec:privacy-baseline}) \\
Consumer council, 3,000 categories & 39.8 kB & 53.6 kB & Dense neighbourhood sum \\
\midrule
\textbf{Round archive, 20k WC + 20k CC} & & & \\
Coarse, single key & \multicolumn{2}{r}{238 MB / 40k sig.\ checks} & 2.4~GB raw per 10-round year \\
Coarse, 5 signers & \multicolumn{2}{r}{787 MB / 200k sig.\ checks} & 7.9~GB raw per year \\
Dense, single key & \multicolumn{2}{r}{910 MB / 40k sig.\ checks} & 9.1~GB raw per year \\
Dense, 5 signers & \multicolumn{2}{r}{1.46 GB / 200k sig.\ checks} & 14.6~GB raw per year \\
\bottomrule
\end{tabularx}
\par\smallskip{\scriptsize\raggedright Header size (400~B) is an estimate; compression ratios for signature-dominated payloads are poor because signatures are incompressible. Source: \path{research/analysis/verify_r3.py}, section E.\par}
\end{table}

Even at these sizes, a round's proposals are small enough that full replication is cheaper than heavy internal-vector proofs. DPEP therefore uses a \textbf{Two-Root Authenticated Aggregation Protocol} instead of internal-vector Merkle-sum proofs:
\begin{align}
  H_{\text{proposals}}^{(t)} &= \text{BLAKE3-Root}\Big(\big\{\big(\text{id}_c, \, \text{BLAKE3}(\text{CBOR}(\pi_c^{(t)}))\big)\big\}_{c=1}^M\Big) \\
  H_{\text{totals}}^{(t)} &= \text{BLAKE3-Root}\Big(\big\{\big(i, \, S_i^{(t)}, \, D_i^{(t)}, \, P_i^{(t+1)}\big)\big\}_{i \in \Omega}\Big)
\end{align}
\begin{itemize}[leftmargin=1.6em,itemsep=2pt]
  \item \textbf{512-Byte Mobile Inclusion Proofs:} Any mobile client (\texttt{Role-LP}) verifies that its council's exact proposal was included in \texttt{proposals\_merkle\_root} using a standard 16-level hash path of $16 \times 32\text{ B} = \mathbf{512\text{ bytes}}$.
  \item \textbf{Full Leaf Data Availability (DA) \& Conservation Verification:} Auditing servers (\texttt{Role-SR}, \texttt{Role-PCW}, \texttt{Role-BFT}) download the round's leaf archive (0.24--1.46~GB raw; Table~\ref{tab:proposal-sizes}), verify the council signatures against the accredited roster, and check that the exact fixed-point sum equals \texttt{category\_totals\_root} before signing the checkpoint. That is 40,000--200,000 ML-DSA-65 verifications per round. Per-signature verification time on server hardware has not been benchmarked here, so whether every server can do this inside the round window is a Kill Gate~2 measurement, not an assumption; the fallback is to shard verification across servers and have each attest to its shard.
  \item \textbf{Distinguishing Censorship from Network Partitions (\path{council_participation_bitmap}):} Each checkpoint commits a 2-bit-per-council bitmap (\texttt{01 = Fresh Round-$t$ Proposal}, \texttt{10 = Timeout/Partition Carry-Forward}, \texttt{00 = Excluded}). Councils submit proposals to both their primary aggregator and $\ge 2$ independent \texttt{Role-SR} relays, receiving timestamped signed \texttt{SubmissionReceipts}. Any council falsely marked \texttt{10} or \texttt{00} broadcasts its receipt during the pre-finality challenge window $T_{\text{challenge}}$ to force inclusion and fault the censoring aggregator.
  \item \textbf{Explicit Cryptographic vs.\ Physical Boundary:} Cryptography verifies digital authenticity, inclusion, and arithmetic conservation; it \textbf{cannot} verify that physical grain was harvested or that emissions sensors were not bypassed, which remain governed by federated physical \emph{euthynai} spot-audits and counterparty delivery reconciliation.
\end{itemize}

Figure~\ref{fig:two-root} shows the two trees, the phone's inclusion path, the participation bitmap and the challenge that turns an omission into a public fault.

\begin{figure}[htbp]
\centering
\begin{tikzpicture}[x=1in,y=1in,
  hn/.style={circle, draw=NavyPrimary, fill=SoftBg, thick, minimum size=0.11in, inner sep=0pt},
  hp/.style={hn, draw=CrimsonAccent, fill=CrimsonAccent!25!white},
  hs/.style={hn, draw=GoldAccent!80!black, fill=GoldAccent!35!white},
  tl/.style={thin, NavyPrimary}]
  % checkpoint on top
  \node[dhead, text width=4.6in] (cp) at (2.45,3.95) {Checkpoint $t$ (Layer 1): $H_{\text{proposals}}$, $H_{\text{totals}}$, participation bitmap, $\ge 86$ of 128 hybrid signatures};
  % ---- left tree: proposals
  \foreach \ox/\name in {0/L, 2.5/R}{
    \node[hn] (\name r) at (\ox+1.2,3.25) {};
    \node[hn] (\name a) at (\ox+0.66,2.75) {};  \node[hn] (\name b) at (\ox+1.74,2.75) {};
    \node[hn] (\name c) at (\ox+0.39,2.25) {};  \node[hn] (\name d) at (\ox+0.93,2.25) {};
    \node[hn] (\name e) at (\ox+1.47,2.25) {};  \node[hn] (\name f) at (\ox+2.01,2.25) {};
    \foreach \i in {0,...,7}{ \node[hn] (\name l\i) at (\ox+0.25+\i*0.27,1.75) {}; }
    \draw[tl] (\name r)--(\name a) (\name r)--(\name b) (\name a)--(\name c) (\name a)--(\name d) (\name b)--(\name e) (\name b)--(\name f)
              (\name c)--(\name l0) (\name c)--(\name l1) (\name d)--(\name l2) (\name d)--(\name l3)
              (\name e)--(\name l4) (\name e)--(\name l5) (\name f)--(\name l6) (\name f)--(\name l7);
  }
  % highlight proof path in left tree: leaf 2 -> d -> a -> r ; siblings l3, c, b
  \node[hp] at (Ll2) {}; \node[hp] at (Ld) {}; \node[hp] at (La) {}; \node[hp] at (Lr) {};
  \node[hs] at (Ll3) {}; \node[hs] at (Lc) {}; \node[hs] at (Lb) {};
  \draw[very thick, CrimsonAccent] (Ll2)--(Ld)--(La)--(Lr);
  \draw[darr] (Lr) -- (Lr |- cp.south);
  \draw[darr] (Rr) -- (Rr |- cp.south);
  \node[font=\tiny\bfseries, anchor=east] at (Lr.west) {$H_{\text{proposals}}$\,};
  \node[font=\tiny\bfseries, anchor=west] at (Rr.east) {\,$H_{\text{totals}}$};
  \node[dbox, text width=2.2in] at (1.2,1.27) {Leaf per council: $(\mathrm{id}_c,\ \mathrm{BLAKE3}(\mathrm{CBOR}(\pi_c^{(t)})))$\\
    \textcolor{CrimsonAccent}{\textbf{path}} + \textcolor{GoldAccent!70!black}{\textbf{siblings}}: a phone checks its council's\\leaf with 16 levels $\times$ 32 B $=$ \textbf{512 B}};
  \node[dbox, text width=2.2in] at (3.7,1.27) {Leaf per category: $(i,\ S_i^{(t)},\ D_i^{(t)},\ P_i^{(t+1)})$\\
    auditors holding the full leaf archive (0.24--1.46 GB)\\recompute every sum: \textbf{conservation check}};
  % ---- bitmap and censorship
  \node[font=\tiny\bfseries, anchor=west] at (0.02,0.78) {Participation bitmap (2 bits per council):};
  \foreach \i/\v/\c in {0/01/TealAccent, 1/01/TealAccent, 2/10/GoldAccent!80!black, 3/01/TealAccent, 4/00/CrimsonAccent, 5/01/TealAccent, 6/01/TealAccent, 7/10/GoldAccent!80!black}{
    \node[rectangle, draw=\c, thick, fill=white, minimum width=0.22in, minimum height=0.15in, font=\tiny\ttfamily, inner sep=0pt] at (2.2+\i*0.25,0.78) {\v};
  }
  \node[font=\tiny, anchor=west, align=left] at (0.02,0.55) {\texttt{01} fresh proposal \quad \texttt{10} carried forward (timeout or partition) \quad \texttt{00} excluded};
  \node[dwarn, text width=1.35in] (x1) at (0.72,0.2) {Council wrongly marked \texttt{00} or \texttt{10}};
  \node[dbox, text width=1.35in] (x2) at (2.45,0.2) {shows its signed \texttt{SubmissionReceipt} within $T_{\text{challenge}}$};
  \node[dinst, text width=1.35in] (x3) at (4.18,0.2) {forced inclusion; the censoring server is publicly faulted};
  \draw[darr] (x1)--(x2); \draw[darr] (x2)--(x3);
\end{tikzpicture}
\caption{Two-Root authenticated aggregation. A phone verifies its council's leaf with a 512-byte path; auditors with the leaf archive recompute every category sum; the participation bitmap and signed receipts separate censorship from partition.}
\label{fig:two-root}
\end{figure}

\section{Granular Consistency Mapping and Node Role Taxonomy}
\label{sec:roles}

\begin{table}[htbp]
\centering
\caption{Consistency Models and Partition Tradeoffs Across Economic Subsystems}
\label{tab:consistency}
\scriptsize
\setlength{\tabcolsep}{3.5pt}
\begin{tabularx}{\textwidth}{L{0.95in} L{1.00in} C{0.62in} L{1.00in} Y}
\toprule
\textbf{Economic Subsystem} & \textbf{Consistency Model} & \textbf{Latency Target} & \textbf{CAP Partition Behaviour} & \textbf{Architectural Justification} \\
\midrule
\textbf{1. Proposal Drafting (\texttt{S1--S2})} & Strong Eventual Consistency (Delta-CRDT + SQLite)\footnotemark & $<50\text{ ms}$ local & \textbf{AP (Offline-First):} Drafts continue on local LAN/device. & Unsubmitted drafts affect no external council; offline usability is paramount. \\
\addlinespace[3pt]
\textbf{2. Council Ratification Vote} & Council-Scoped Quorum ($k$-of-$n$ multi-signature) & Seconds to hours & \textbf{Locally CP on Council LAN:} Independent of WAN. & Only members of $WC_j$ or $CC_h$ hold authority over their council's proposal. \\
\addlinespace[3pt]
\textbf{3. Federated Round Aggregation (\texttt{S3})} & Deterministic Two-Root Reduction & $1\text{--}10\text{ min}$ per round & \textbf{Bounded-Wait CP + Fallback:} Carries forward $\pi_c^{(t-1)}$ if WAN severed past grace window. & Prevents a single disconnected hamlet from halting national planning while preserving verifiable inclusion. \\
\addlinespace[3pt]
\textbf{4. IFB Price Update \& Final Plan (\texttt{S4, S6})} & Global Total Order / Linearizable Epoch BFT\footnotemark & $2\text{--}15\text{ s}$ per block & \textbf{CP (Safety over Liveness):} Requires $\ge 2f+1$ of $3f+1$ federation validators. & Every council must optimise round $t+1$ against the \emph{exact same} indicative price vector $P^{(t+1)}$. \\
\addlinespace[3pt]
\textbf{5. Intra-Year Fulfilment (\texttt{S7})} & Bilateral Signed Receipts + Federated Rollup & $<500\text{ ms}$ (offline-capable) & \textbf{AP with Bounded Credit:} Daily depot credit cap reconciled at Ward rollup. & Point-of-delivery pickups at neighbourhood depots must never block on global consensus. \\
\bottomrule
\end{tabularx}
\end{table}
\addtocounter{footnote}{-1}
\footnotetext{Marc Shapiro, Nuno Pregui\c{c}a, Carlos Baquero, and Marek Zawirski, ``Conflict-Free Replicated Data Types,'' \emph{Symposium on Self-Stabilizing Systems (SSS 2011)}, LNCS 6976, pp.~386--400.}
\stepcounter{footnote}
\footnotetext{Ethan Buchman, Jae Kwon, and Zarko Milosevic, ``The Latest Gossip on BFT Consensus,'' \emph{arXiv:1807.04938} (2018); Maofan Yin et al., ``HotStuff: BFT Consensus with Linearity and Responsiveness,'' \emph{ACM PODC '19}, pp.~347--356.}

Figure~\ref{fig:topology} places the six node roles of Table~\ref{tab:node-roles} in the federation structure.

\begin{figure}[htbp]
\centering
\begin{tikzpicture}[x=1in,y=1in, tree/.style={thick, NavyPrimary}]
  \node[dhead, text width=3.3in] (bft) at (2.45,4.2) {Epoch BFT committee: 128 validators drawn by sortition (Role-BFT)};
  \node[dbox, text width=1.75in] (cf) at (1.2,3.5) {\textbf{Consumer federation}\\region and city};
  \node[dbox, text width=1.45in] (if) at (3.85,3.5) {\textbf{Industry federation}};
  \node[dbox, text width=1.75in] (ws) at (1.2,2.8) {\textbf{Ward and city servers}\\Role-PCW, Role-SR (community servers)};
  \node[dbox, text width=1.45in] (is) at (3.85,2.8) {\textbf{Industry servers}\\Role-PCW, Role-SR};
  \node[dbox, text width=0.62in] (c1) at (0.4,2.0) {Neighbourhood council\\Role-CC (Pi,~5~W)};
  \node[dbox, text width=0.62in] (c2) at (1.2,2.0) {Neighbourhood council\\Role-CC};
  \node[dbox, text width=0.62in, dashed] (c3) at (2.0,2.0) {Offline hamlet council\\Role-CC};
  \node[dbox, text width=0.62in] (w1) at (3.45,2.0) {Worker council\\Role-CC};
  \node[dbox, text width=0.62in] (w2) at (4.25,2.0) {Worker council\\Role-CC};
  \foreach \n/\x in {p1/0.4, p2/1.2, p3/2.0, q1/3.45, q2/4.25}{
    \node[dcrypto, text width=0.62in] (\n) at (\x,1.22) {members' phones\\Role-LP};
  }
  \draw[tree] (bft.south -| cf.north) -- (cf.north);  \draw[tree] (bft.south -| if.north) -- (if.north);
  \draw[tree] (cf)--(ws); \draw[tree] (if)--(is);
  \draw[tree] (ws.south)--(c1.north); \draw[tree] (ws.south)--(c2.north);
  \draw[tree, dashed] (ws.south)-- node[dlab, pos=0.55]{Role-GW\\data mule} (c3.north);
  \draw[tree] (is.south)--(w1.north); \draw[tree] (is.south)--(w2.north);
  \foreach \c/\p in {c1/p1, c2/p2, c3/p3, w1/q1, w2/q2}{ \draw[tree] (\c)--(\p); }
  \draw[SlateDark, densely dotted, thick] (p2.south east) to[bend right=14] node[dlab, pos=0.5]{the same person: one consumer and one worker council} (q1.south west);
  % aggregators outside the tree
  \node[dcrypto, text width=1.2in] (A) at (0.75,0.12) {\textbf{Aggregator A}\\consumer-side operator};
  \node[dcrypto, text width=1.2in] (B) at (2.45,0.12) {\textbf{Aggregator B}\\producer-side operator};
  \node[dcrypto, text width=1.2in] (C) at (4.15,0.12) {\textbf{Aggregator C}\\external operator};
  \foreach \t in {A,B,C}{ \draw[darr, TealAccent, dashed] (p1.south) -- (\t.north); }
  \node[font=\tiny\itshape, text=TealAccent, align=center] at (2.45,-0.22) {Each member's phone sends its \texttt{Prio3} report shares to three aggregators that sit \emph{outside} the federation tree (non-nested).};
\end{tikzpicture}
\caption{Federation topology and node placement (illustrative). Councils and servers nest by geography and industry; the three \texttt{Prio3} aggregators sit outside that nesting, and the BFT committee is drawn by sortition across federations.}
\label{fig:topology}
\end{figure}

\begin{table}[htbp]
\centering
\caption{Six-Role Hardware Realism Matrix (Supporting Legacy Smartphones and Raspberry Pis)}
\label{tab:node-roles}
\scriptsize
\setlength{\tabcolsep}{3.5pt}
\begin{tabularx}{\textwidth}{L{0.98in} L{1.35in} C{0.75in} C{0.75in} Y}
\toprule
\textbf{Node Role} & \textbf{Core Responsibilities} & \textbf{Entry Smartphone (3GB RAM ARM)?} & \textbf{Raspberry Pi 4/5 (4--8GB RAM)?} & \textbf{Community Server (8+ Cores)?} \\
\midrule
\textbf{1. Lightweight Participant (\texttt{Role-LP})} & Local household budget QP solver, \texttt{Prio3} share generation, CRDT sync, Two-Root \& BFT header verifier, ballot signer. & \textbf{YES (Primary)} (RAM unverified) & \textbf{YES} & \textbf{YES} \\
\addlinespace[3pt]
\textbf{2. Council Coordinator (\texttt{Role-CC})} & Hosts $WC_j$ or $CC_h$ local Wi-Fi CRDT relay, runs workplace LP/QP solver, assembles the council $k$-of-$n$ multi-signature. & \textbf{YES (Failover Hotspot)} & \textbf{YES (Ideal \$50 Appliance at 5W)} & \textbf{YES} \\
\addlinespace[3pt]
\textbf{3. Storage Replica (\texttt{Role-SR})} & Pins regional content-addressed BLAKE3 blobs for qualitative \texttt{Lists} and audit logs. & \textbf{NO} (Flash wear \& cell data limits) & \textbf{YES} (with USB SSD) & \textbf{YES (Primary)} \\
\addlinespace[3pt]
\textbf{4. Planning Compute Worker (\texttt{Role-PCW})} & Computes Two-Root reductions, deterministic WASM price updates, and input-output analytics. & \textbf{NO} & \textbf{PARTIAL} (Up to $5,000$ councils) & \textbf{YES (Primary IFB Node)} \\
\addlinespace[3pt]
\textbf{5. Epoch Consensus Validator (\texttt{Role-BFT})} & Participates in sortition-elected BFT committee finalising round checkpoints. & \textbf{NO} (Requires 24/7 fixed routing) & \textbf{YES} (For PoC / local federations) & \textbf{YES (Primary)} \\
\addlinespace[3pt]
\textbf{6. Delay-Tolerant Gateway (\texttt{Role-GW})} & Bridges offline mesh/LoRa/USB sneakernet bundles to federation relays. & \textbf{YES (As ``Data Mule'')} & \textbf{YES (Mesh Box)} & \textbf{YES} \\
\bottomrule
\end{tabularx}
\end{table}

% ==============================================================================
% CHAPTER 6: TECHNOLOGY AND ARCHITECTURE COMPARISON
% ==============================================================================
\chapter{Weighted Architecture and Implementation Stack Comparison}

\chapterepigraph{Devising and choosing from among these and other possibilities is a practical issue in implementing any actual participatory economy. Assuming the procedures chosen do not violate principles essential to participatory planning, considerations include (1) the extent to which iteration workers could bias outcomes, (2) the extent of reductions in the number of iterations required to reach a plan\dots}{Michael Albert and Robin Hahnel, \emph{The Political Economy of Participatory Economics} (1991), Ch.~4, p.~68\protect\footnotemark}
\footnotetext{Albert and Hahnel, \emph{Political Economy of Participatory Economics}, Ch.~4, p.~68.}

\section{Weighted Comparison of the Four Architectural Paradigms}

We compare the four candidate architectures across all 14 required criteria: \textbf{Architecture A} (Conventional Federated SQL/REST Services), \textbf{Architecture B} (Pure P2P Gossip + CRDTs without finality), \textbf{Architecture C} (Full Per-Transaction BFT Blockchain), and \textbf{Architecture D} (Hybrid Local-First + Two-Root Aggregation + Selective Epoch-BFT Checkpoints). Scores are expert judgements of the \emph{current} design, not measured performance. Rows 4, 5, 6, 9 and 12 are scored against the measured \texttt{Prio3} costs, the reworked proposal sizes (Table~\ref{tab:proposal-sizes}) and the carry-forward constraint in Section~\ref{sec:carry-forward}. A leads D on five criteria (throughput, bandwidth, hardware, implementation simplicity, energy) and D equals or beats A on the other nine.

\begin{table}[htbp]
\centering
\caption{Weighted 14-criterion comparison (1--5 scale): A federated SQL/REST, B pure P2P gossip, C per-transaction blockchain, D hybrid}
\label{tab:arch-comp}
\scriptsize
\setlength{\tabcolsep}{2.8pt}
\begin{tabularx}{\textwidth}{L{0.98in} C{0.24in} C{0.24in} C{0.24in} C{0.24in} C{0.26in} Y}
\toprule
\textbf{Criterion} & \textbf{Wt.} & \textbf{A} & \textbf{B} & \textbf{C} & \textbf{D} & \textbf{Rationale for D (judgement, not measurement)} \\
\midrule
1. Decentralisation & 9\% & 2.0 & 5.0 & 3.0 & \textbf{4.5} & Local autonomy plus rotating BFT; A centralises database administration, C heavy validators. \\
2. Trust Assumptions & 10\% & 2.0 & 3.0 & 4.5 & \textbf{4.0} & One honest aggregator for privacy, \emph{all} for correct totals (VDAF-22 \S 9); correctness rests on institutions, not cryptography. \\
3. Censorship Resistance & 8\% & 2.0 & 4.5 & 4.0 & \textbf{4.5} & Signed receipts let a censored council prove omission. \\
4. Throughput \& Latency & 8\% & 5.0 & 4.0 & 2.0 & \textbf{4.0} & $10^6$ report verifications cost each aggregator 4--61 desktop core-minutes; A has no such cost. \\
5. Network Bandwidth & 7\% & 4.5 & 2.5 & 1.5 & \textbf{3.5} & No $\mathcal{O}(N^2)$ gossip, but 290~kB quorum signatures and 0.24--1.46~GB round archives. \\
6. Hardware Requirements & 8\% & 4.5 & 3.5 & 1.5 & \textbf{4.0} & Raspberry Pi and Android are targets; no device benchmark yet. \\
7. Offline Capability & 9\% & 1.5 & 5.0 & 1.0 & \textbf{5.0} & SQLite + CRDTs + DTN store-and-forward. \\
8. Byzantine Fault Tolerance & 8\% & 1.0 & 2.5 & 5.0 & \textbf{4.0} & Detects leaf tampering given data availability; does not bind aggregator shares. \\
9. Privacy & 9\% & 2.5 & 2.0 & 1.0 & \textbf{3.5} & Coarse \texttt{Prio3}; no report carry-forward; credentials unspecified; differencing remains. \\
10. Post-Quantum Readiness & 6\% & 4.0 & 3.5 & 2.0 & \textbf{4.0} & FIPS 203--205 hybrids; BLAKE3 128-bit; not a uniform Level-3 claim. \\
11. Implementation Simplicity & 6\% & 5.0 & 3.0 & 2.5 & \textbf{3.5} & Four layers can be rolled out in phases, but far more code than A. \\
12. Operational Cost & 4\% & 4.0 & 4.5 & 2.0 & \textbf{4.0} & $\approx 298$~MB/yr checkpoints, 2.4--14.6~GB/yr raw archive; audits dominate (note). \\
13. Energy Efficiency & 4\% & 4.5 & 4.0 & 2.5 & \textbf{4.0} & No proof-of-work; no sourced energy ratio claimed. \\
14. Maintainability & 4\% & 4.0 & 3.5 & 2.0 & \textbf{4.0} & Economic rules as upgradable, ratified WASM modules. \\
\midrule
\textbf{Weighted score} & \textbf{100\%} & \textbf{3.10} & \textbf{3.59} & \textbf{2.54} & \textbf{4.07} & D leads under these judgements (Section~\ref{sec:arch-sensitivity}). \\
\bottomrule
\end{tabularx}
\par\smallskip
{\scriptsize\raggedright \emph{Note on row 12} (\path{research/analysis/parameters.py}, \path{verify_r3.py}; SI units, $1\text{ kB}=1{,}000\text{ B}$): at $N=10^6$, $G=10{,}000$, $n=128$, a Layer-1 block is 5 \texttt{FixedDecimal} arrays ($400{,}000\text{ B}$), $\lceil 2n/3\rceil = 86$ signatures ($290{,}078\text{ B}$), a 10{,}000~B participation bitmap and two 32~B roots: $700{,}142\text{ B} \approx 700\text{ kB}$. A 10-round plan has $10+365+50=425$ blocks/year, $\approx 298\text{ MB/year}$; at 20 rounds (435 blocks), $\approx 305\text{ MB/year}$. Layer-2 archive sizes come from Table~\ref{tab:proposal-sizes}; index, receipt, qualitative-\texttt{List} and replica overhead is unspecified. Hosting 100 federation servers at \$30--\$60/month is \$3,000--\$6,000/month and excludes audits, facilitation labour and research. All figures are unbenchmarked.\par}
\end{table}

\FloatBarrier
\subsection{Sensitivity Analysis of Architecture Rankings}
\label{sec:arch-sensitivity}
All figures below are computed from the cells of Table~\ref{tab:arch-comp} by a script that parses the table from the book source.\footnote{The script is \path{research/analysis/verify_r3.py} in the project repository.}
\begin{itemize}[leftmargin=1.6em,itemsep=1pt]
  \item \textbf{Baseline weights:} D~4.07, B~3.59, A~3.10, C~2.54. \textbf{Equal weights} ($1/14$ each): D~4.04, B~3.61, A~3.32, C~2.46.
  \item \textbf{Weight on simplicity, throughput and cost.} Put a share $w$ of the weight on criteria 4, 11 and 12 equally and spread the rest in proportion to the baseline. A overtakes D at $w \approx 0.63$; at $w=0.70$ the scores are A~4.08, D~3.92, B~3.75, C~2.30. The ranking of A against D therefore depends on how much a community values operational simplicity, and that is a legitimate reason to use A for Phase~1H and the earliest pilots.
  \item \textbf{Harsher D cells} (trust 3.5, bandwidth 2.5, hardware 3.0, privacy 3.0): D~3.82, still ahead of B~3.59 and A~3.10 at baseline weights.
  \item \textbf{Score uncertainty.} Perturbing every cell of every architecture independently by $U(-1,1)$ points (clipped to $[1,5]$; 20,000 draws) leaves D first in 99.2\% of draws and B first in 0.8\%; D's median lead over the runner-up is 0.50 points and its 5th-percentile lead 0.15.
\end{itemize}
D's lead at baseline weights is robust to plausible errors in individual scores, but not to a weighting that puts simplicity first.

\section{Two-Stage Language Strategy: Python-First Research to Rust Protocol}

Rather than locking the project prematurely into systems-level Rust before the non-convex economic coordination algorithms are validated, we adopt a \textbf{two-stage implementation strategy}:
\begin{enumerate}[leftmargin=1.8em,itemsep=2pt]
  \item \textbf{Phase 1 (Economic Coordination Research) --- Python Scientific Stack (\texttt{NumPy}/\texttt{SciPy}, \texttt{highspy} MILP/LP, \texttt{Clarabel} Python bindings):} Minimises iteration time and maximises reproducibility for testing candidate ADMM, MILP federation decomposition, and Scarf/Leontief stress tests.
  \item \textbf{Phase 2 \& 3 (Validated Distributed Protocol \& Mobile Engine) --- Rust + WebAssembly (\texttt{dpep-core}, \texttt{dpep-node}):} Translates the Phase-1-validated clearing rules into memory-safe Rust (\texttt{Clarabel.rs},\footnote{Paul J.\ Goulart and Yuwen Chen, ``Clarabel: An Interior-Point Solver for Conic Programs with Quadratic Objectives,'' \emph{Mathematical Programming Computation} (2026), \url{https://doi.org/10.1007/s12532-026-00320-7}; software at \url{https://github.com/oxfordcontrol/Clarabel.rs}.} FIPS 203/204/205 implementations with independent assurance (\texttt{libcrux}, formally verified with hax and F$^\ast$; \texttt{aws-lc-rs}, FIPS 140-3 validated; the unaudited \texttt{RustCrypto} crates only as a WASM fallback),\footnote{The \texttt{RustCrypto} \texttt{ml-kem} documentation states that the implementation ``has never been independently audited'' (\url{https://docs.rs/ml-kem}); \texttt{libcrux-ml-kem} is verified with hax and F$^\ast$ (\url{https://docs.rs/libcrux-ml-kem}) and its generated C is vendored in Mozilla NSS; AWS-LC FIPS~3.0 is a FIPS 140-3 validated module including ML-KEM (\url{https://aws.amazon.com/blogs/security/aws-lc-fips-3-0-first-cryptographic-library-to-include-ml-kem-in-fips-140-3-validation}). Checked November 2026.} \texttt{quinn} QUIC with \texttt{rustls} and its hybrid \texttt{X25519MLKEM768} key exchange, \texttt{Wasmtime}, and Mozilla UniFFI). Local council solvers use IEEE-754 floating-point heuristics to \emph{draft} proposals. Protocol-level checks ($SB_j/SC_j \ge 1$, $SC_h \le B_h$, IFB price updates) are specified in deterministic 64/128-bit fixed-point arithmetic (Section~\ref{sec:fixed-point-bounds}). Those checks verify constraint satisfaction after rounding, not a truthful objective or MILP optimality.
\end{enumerate}

\begin{table}[htbp]
\centering
\caption{Programming Language Comparison Across Core Protocol Requirements (1--5 Scale)}
\label{tab:lang-comp}
\scriptsize
\setlength{\tabcolsep}{3.0pt}
\begin{tabularx}{\textwidth}{L{1.25in} C{0.38in} C{0.38in} C{0.52in} C{0.45in} C{0.40in} Y}
\toprule
\textbf{Technical Dimension} & \textbf{Rust} & \textbf{Go} & \textbf{TypeScript} & \textbf{Kotlin} & \textbf{Swift} & \textbf{WASM (from Rust)} \\
\midrule
Memory Safety \& Zero GC Pauses & \textbf{5.0} & 3.5 & 2.5 & 3.0 & 4.0 & \textbf{5.0} (Sandboxed linear memory) \\
NIST PQC (\texttt{libcrux}, \texttt{aws-lc-rs}, \texttt{RustCrypto}) & \textbf{5.0} & 4.0 & 2.0 & 2.5 & 2.5 & \textbf{5.0} (Direct compilation) \\
Local-First \& P2P (\texttt{iroh}, CRDT, SQLite) & \textbf{5.0} & 4.5 & 4.0 & 3.0 & 3.0 & \textbf{4.5} (Browser worker target) \\
Operations Research (\texttt{Clarabel.rs}, \texttt{HiGHS}) & \textbf{4.5} & 2.5 & 2.0 & 2.5 & 2.5 & \textbf{5.0} (Runs on every phone/PC) \\
Cross-Platform Mobile Reuse (UniFFI) & \textbf{5.0} & 3.0 & 4.5 & 3.0 & 1.5 & \textbf{5.0} (Universal execution) \\
Accessible UI \& Screen-Reader Tooling & 2.5 & 1.5 & \textbf{5.0} & \textbf{4.5} & \textbf{4.5} & Paired with TypeScript PWA UI \\
\bottomrule
\end{tabularx}
\end{table}

\section{Component Maturity}
\label{sec:maturity}

A component can be mature at one level and not at another, so Table~\ref{tab:maturity} rates three levels separately: the \textbf{algorithm or protocol}, the \textbf{library} this book names, and the \textbf{DPEP combination} (the way DPEP uses it). The scale runs from M5 (\emph{ubiquitous}: final standard, several independent implementations, billions of users, ten or more years of public analysis) through M4 (\emph{production-proven} by several major operators, with audits or formal verification), M3 (\emph{early production}: a new standard, an Internet-Draft, or few operators), M2 (\emph{pre-production}: unaudited, pre-1.0 or beta), M1 (\emph{research}) to M0 (\emph{unspecified}). Ratings are judgements as of November 2026, not measurements. Evidence, production users, the researchers behind each component and sources are in the project's maturity dossier and register.\footnote{\path{reports/STACK_MATURITY.md} and \path{research/stack_maturity_register.csv} in the project repository.}

\begin{table}[htbp]
\centering
\caption{Component maturity at three levels (M0--M5; judgement as of November 2026)}
\label{tab:maturity}
\scriptsize
\setlength{\tabcolsep}{3pt}
\begin{tabularx}{\textwidth}{L{1.30in} C{0.42in} C{0.62in} C{0.46in} Y}
\toprule
\textbf{Component} & \textbf{Algorithm} & \textbf{Library} & \textbf{DPEP use} & \textbf{Main gap} \\
\midrule
X25519 / Ed25519 & M5 & M4--M5 & M4 & None material. \\
Hybrid X25519 + ML-KEM-768 & M4 & M4 & M2 & Default in major browsers, OpenSSH~10 and Apple OS~26 TLS; the gap is DPEP's own use. \\
Hybrid Ed25519 + ML-DSA-65 & M3 & M2--M3 & M2 & Composite format is a draft; 3.4~kB per signature. \\
SLH-DSA-SHA2-192s & M3 & M2--M3 & M3 & Little production use; acceptable for rare root signing. \\
BLAKE3 & M3 & M4 & M3 & Not standardised; SHA3-256 registered as alternative. \\
\texttt{Prio3} VDAF / DAP & M3 & M3 & M1 & Budget circuit unspecified; library tracks draft~18. \\
Post-quantum HPKE for shares & M2 & M2 & M1 & No standard suite yet. \\
Two-Root aggregation & M5 & M4 & M1 & Merkle logs mature (Certificate Transparency); DPEP protocol original. \\
One-report-per-person credentials & M4 / M1\textsuperscript{a} & --- & M1 & Interim classical (Section~\ref{sec:pq-urgency}); $B_u$ binding open. \\
Sortition beacon & M4 / M1\textsuperscript{a} & --- & M1 & Interim threshold beacon; key generation and custody unspecified. \\
Capability certificates (CapCert) & M2 & M2 & M0 & Format not chosen (UCAN 1.0 release candidate, Biscuit). \\
QUIC: \texttt{quinn} + \texttt{rustls} & M5 & M4 & M3 & Hybrid key exchange must be enabled explicitly. \\
\texttt{iroh} & M3 & M3 & M2 & 1.0 released June 2026; relays must be self-hosted. \\
SQLite / SQLCipher; CRDTs & M5 / M4 & M3--M5 & M3 & CRDT schema evolution and metadata growth. \\
CometBFT & M4--M5 & M4 / M2\textsuperscript{b} & M2 & Post-quantum keys are pure ML-DSA, not hybrid. \\
Fixed-point WASM on \texttt{Wasmtime} & M4 & M4 & M2 & Pin a long-term-support release; fix NaN and fuel settings. \\
\texttt{HiGHS} / \texttt{Clarabel} & M4 & M4 / M2\textsuperscript{c} & M2 & No phone benchmark. \\
Mozilla UniFFI & M3--M4 & M3--M4 & M2 & Interface still changing. \\
\bottomrule
\end{tabularx}
\par\smallskip{\scriptsize\raggedright \textsuperscript{a}\,Classical / post-quantum. \textsuperscript{b}\,Ed25519 keys / ML-DSA-65 keys (v0.40). \textsuperscript{c}\,Native / WASM on a phone.\par}
\end{table}

Nine DPEP combinations have no deployment precedent, so they inherit no maturity: the \texttt{Prio3} budget circuit, Two-Root aggregation, the hybrid signature format, CapCerts, credentials, the beacon, the light-client policy (Table~\ref{tab:light-clients}), non-nested aggregator governance, and ratified WASM economic rules. Two-Root aggregation, WASM rules and server-side signature verification are tested at Kill Gate~2. The budget circuit, hybrid format, CapCerts, credentials, light-client policy and aggregator governance are tested at Kill Gate~3, whose external audit covers them. The beacon has an interim construction and stays an open question. Post-quantum deployment is moving quickly, so the hybrid-signature, \texttt{Prio3}, share-encryption, credential, beacon, \texttt{iroh} and CometBFT rows are re-checked before each gate.

% ==============================================================================
% CHAPTER 7: POST-QUANTUM CRYPTOGRAPHY, SECURITY, AND PRIVACY
% ==============================================================================
\chapter{Post-Quantum Cryptography, Democratic Membership, and Privacy}

\chapterepigraph{Decisions about what individuals wish to consume will be subject to collective criticism by fellow council members, though with specific guarantees for preserving individual freedoms and privacy.}{Michael Albert and Robin Hahnel, \emph{Looking Forward} (1991), Ch.~3, ``Consumption Norms''\protect\footnotemark}
\footnotetext{Albert and Hahnel, \emph{Looking Forward}, Ch.~3 (\texttt{lf3\_2.htm}); Michael Albert, \emph{Parecon: Life After Capitalism} (Verso, 2003), Ch.~12 (``Consuming'') and Ch.~18 (``Privacy / Frenzy: A Busybody Economy?''), where Albert specifies that individual consumption proposals within warranted effort-rated budgets are submitted anonymously to prevent ``busybody'' intrusion; see also Jan Philipp Dapprich, ``Democratic Theory and Economic Planning: A Critique of Participatory Economics,'' \emph{Erasmus Journal for Philosophy and Economics} 17, no. 2 (Winter 2024), \url{https://ejpe.org/journal/article/view/867}.}

\section{What the Source Tradition Requires of Consumption Privacy}
\label{sec:consumption-privacy}
The epigraph is the 1991 design: consumption is \emph{subject to collective criticism}, with unspecified ``guarantees'' for freedom and privacy. PEPE Ch.~3 (p.~41) and Ch.~5 (p.~80) add a non-veto rule for in-budget requests. Albert (2003, Chs.~12, 18) argues for anonymous submission against busybody intrusion. None of those texts specifies cryptographic aggregation. DPEP's \texttt{Prio3} invisibility is a \textbf{technical extension}: neighbours cannot see or comment on the anonymous bundle, so the 1991 criticism channel is lost. That is a genuine normative cost, not a proven improvement. REQ-DATA-04 therefore splits: non-veto is \texttt{[ORIGINAL-1991]}; invisibility is \texttt{[TECH-EXT]}.

\section{Post-Quantum Cryptography (NIST FIPS 203, 204, 205)}

To protect multi-decade household records against harvest-now-decrypt-later adversaries, DPEP uses cryptographic agility (\texttt{CryptoSuiteId}) and \textbf{hybrid} classical + NIST PQC primitives. This is \textbf{not} a uniform NIST Level-3 claim, and it is not a ``zero lattice'' claim: ML-KEM and ML-DSA are lattice schemes. BLAKE3's specification targets 128-bit security for all goals. VDAF-22 makes no post-quantum claim. Table~\ref{tab:pq-components} states the status of each component.\footnote{NIST, \emph{FIPS 203 (ML-KEM)}, \emph{FIPS 204 (ML-DSA)}, and \emph{FIPS 205 (SLH-DSA)} (13 August 2024). Byte sizes in Table~\ref{tab:pqc-bench} are specification constants. The latency columns are placeholders carried from an unsourced estimate; in particular $\sim 80\ \mu\text{s}$ for a hybrid X25519+ML-KEM-768 operation looks optimistic for an entry-level phone. They must be replaced by device benchmarks (Kill Gate~3) before use.}

Two formats are still in standardisation. \textbf{Hybrid signatures} follow the composite ML-DSA Internet-Draft of Ounsworth et al.\ (2026), which defines an Ed25519 combination but is not yet an RFC; until it is, DPEP's hybrid wire format is DPEP-defined under \texttt{CryptoSuiteId}. Plain ML-DSA in X.509 certificates is standardised (Massimo et al.\ 2025). \textbf{Share encryption}: DAP encrypts report shares with HPKE, whose standard suites are classical, so the hybrid share encryption in Table~\ref{tab:pqc-bench} needs a post-quantum HPKE suite (ML-KEM or the X-Wing combiner), which is still in IETF drafting.

\begin{table}[htbp]
\centering
\caption{Post-quantum primitives: specification byte sizes and \emph{placeholder} ARMv8 latencies (unbenchmarked)}
\label{tab:pqc-bench}
\scriptsize
\setlength{\tabcolsep}{2.8pt}
\begin{tabularx}{\textwidth}{L{1.12in} C{0.52in} R{0.44in} R{0.48in} R{0.48in} R{0.48in} Y}
\toprule
\textbf{Primitive \& Standard} & \textbf{Security} & \textbf{Pub Key} & \textbf{Cipher/Sig} & \textbf{ARM Sign} & \textbf{ARM Verify} & \textbf{Assigned Role in DPEP} \\
\midrule
\textbf{X25519 + ML-KEM-768} (FIPS 203 Hybrid KEM) & NIST L3 + Classical & $1,216\text{ B}$ & $1,120\text{ B}$ & $\sim 80\ \mu\text{s}$\textsuperscript{p} & $\sim 85\ \mu\text{s}$\textsuperscript{p} & Point-to-point QUIC channels \& \texttt{Prio3} aggregator share encryption. \\
\addlinespace[3pt]
\textbf{Ed25519 + ML-DSA-65} (FIPS 204 Hybrid Sig) & NIST L3 + Classical & $1,984\text{ B}$ & $3,373\text{ B}$ & $\sim 230\ \mu\text{s}$\textsuperscript{p} & $\sim 175\ \mu\text{s}$\textsuperscript{p} & \textbf{Workhorse Signature} for council proposals, CapCerts \& receipts. \\
\addlinespace[3pt]
\textbf{SLH-DSA-SHA2-192s} (FIPS 205 Hash-Based) & \textbf{NIST L3} (Hash Only) & $48\text{ B}$ & $16,224\text{ B}$ & unverified & unverified & \textbf{Level-3 Root Anchor} for Federation accreditation \& WASM rule signing. Signing is slow by design; device timings not benchmarked. \\
\addlinespace[3pt]
\textbf{BLAKE3} (BLAKE3 spec) / SHA3-256 (FIPS 202) & BLAKE3: 128-bit (all goals). SHA3-256: 128-bit collision. & N/A & $32\text{ B}$ & $<1\ \mu\text{s/kB}$\textsuperscript{p} & $<1\ \mu\text{s/kB}$\textsuperscript{p} & Two-Root Merkle trees \& content CIDs. BLAKE3 targets 128-bit security for preimage as well as collision resistance; quantum (Grover/BHT) reductions are not claimed. \\
\bottomrule
\end{tabularx}
\par\smallskip{\scriptsize\raggedright \textsuperscript{p}\,Placeholder, not measured. Council signatures are $k$-of-$n$ multi-signatures of the hybrid scheme (3,437~B per signer); no post-quantum threshold signature is assumed.\par}
\end{table}

\begin{table}[htbp]
\centering
\caption{Component-by-component post-quantum status}
\label{tab:pq-components}
\scriptsize
\begin{tabularx}{\textwidth}{L{1.35in} L{1.15in} Y}
\toprule
\textbf{Component} & \textbf{Primitive} & \textbf{Claim type} \\
\midrule
Share encryption & Hybrid X25519 + ML-KEM-768 & Standards-based (FIPS 203 + classical). Lattice. NIST L3 for ML-KEM-768. \\
Council / validator signatures & Hybrid Ed25519 + ML-DSA-65 & Standards-based (FIPS 204 + classical). Lattice. \\
Root / WASM signatures & SLH-DSA-SHA2-192s & Standards-based (FIPS 205). Hash-based NIST L3. \\
Transcript hashes / Merkle & BLAKE3 (SHA3-256 registered as alternative) & Spec: 128-bit for all goals. Not NIST Level 3. Collision classically $\approx 128$ bits. Not standardised by NIST or the IETF, so SHA3-256 (FIPS 202) is registered under \texttt{CryptoSuiteId} for checkpoint roots. \\
FLP soundness / joint randomness & Finite-field FLP + XofTurboShake128 & Inferred from VDAF-22; no PQ theorem. Joint randomness is a known attack surface if the XOF is broken. \\
$B_u$ credentials / sortition beacon & Interim classical (Section~\ref{sec:pq-urgency}) & Not post-quantum. Accepted as interim because forgery and bias need a live CRQC; post-quantum replacement is open research. \\
\texttt{Prio3} as a whole & Mix of the above & VDAF-22 Informational draft; no PQ claim. \\
\bottomrule
\end{tabularx}
\end{table}

\subsection{Which Properties Need Post-Quantum Protection Now, and Crypto Agility}
\label{sec:pq-urgency}

Post-quantum protection is urgent for one property and not for the others. \textbf{Confidentiality} must be post-quantum now: ciphertext recorded today can be decrypted once a CRQC exists, so household reports, share encryption and transport use hybrid ML-KEM from the first deployment and are never downgraded. \textbf{Authenticity, Sybil resistance and randomness} need post-quantum protection \emph{before a CRQC exists}, not before deployment: a CRQC can forge future signatures and bias future draws, but it cannot alter a checkpoint whose history is already anchored under a hash-based signature, or redo a committee draw already made.

\begin{designrule}[Interim Classical Constructions Under Crypto Agility (adopted)]
Where a post-quantum construction is standardised and practical, DPEP uses it in hybrid. Where none exists (one-report-per-person credentials, the sortition beacon) or the consensus engine does not yet support the hybrid (Layer-1 validator keys), DPEP uses the most mature available construction, on four conditions:
\begin{enumerate}[leftmargin=1.6em,itemsep=1pt]
  \item \textbf{Suite inside the signature.} Every signed or encrypted object carries its \texttt{CryptoSuiteId} inside the signed bytes, next to its domain-separation tag, so a suite cannot be stripped or swapped.
  \item \textbf{No downgrade below a ratified floor.} The minimum suite for each object type is fixed in the checkpoint and changed only by the dual-federation ratification used for WASM rules (Section~\ref{sec:sortition}); objects below the floor are rejected, never negotiated.
  \item \textbf{Published migration triggers.} A suite is replaced when the first of these occurs: the replacement is a final standard (composite ML-DSA RFC; FIPS~206; a standardised post-quantum credential or VRF); the software in use supports it; or a credible public warning of CRQC progress. Each migration has a stated window with dual signing during the transition.
  \item \textbf{History survives the switch.} Checkpoint history signed under an outgoing suite is re-anchored under SLH-DSA before that suite is withdrawn; per-epoch credentials and the beacon's threshold key are re-issued, not converted.
\end{enumerate}
\end{designrule}

\begin{table}[htbp]
\centering
\caption{Post-quantum urgency, current construction and upgrade path by function}
\label{tab:pq-urgency}
\scriptsize
\setlength{\tabcolsep}{3pt}
\begin{tabularx}{\textwidth}{L{1.05in} L{0.62in} L{1.45in} Y}
\toprule
\textbf{Function} & \textbf{PQ needed} & \textbf{Construction now} & \textbf{Upgrade path and trigger} \\
\midrule
Transport, report-share encryption & Now & Hybrid X25519 + ML-KEM-768 & Already post-quantum; add HQC as a backup KEM once standardised (final expected 2027). \\
Council signatures, receipts, CapCerts & Before CRQC & Hybrid Ed25519 + ML-DSA-65 (DPEP-defined format) & Composite ML-DSA RFC; FN-DSA once FIPS~206 is final (smaller signatures). \\
Layer-1 validator keys & Before CRQC & Strongest suite the engine supports at deployment & Hybrid once a composite standard and engine support exist. Sizes in this book assume the hybrid, the worst case. \\
Root anchors, WASM rules, re-anchoring & Before CRQC (long-lived) & SLH-DSA-SHA2-192s & None needed (hash-based). \\
One-report-per-person credentials & Before CRQC (forgery) & Classical anonymous credential issued jointly by $k$-of-$n$ councils, with a per-epoch nullifier & Lattice credentials once they fit the phone budget of Kill Gate~3. \\
Sortition beacon & Before CRQC (bias) & Threshold beacon run by the federations (threshold BLS) & A post-quantum VRF or VDF beacon once one is standardised. \\
Binding $B_u$ to the credential & --- & Unspecified (Section~\ref{sec:bu-cred}) & Open research. \\
\bottomrule
\end{tabularx}
\end{table}

Figure~\ref{fig:agility} draws the same table as migration lanes.

\begin{figure}[htbp]
\centering
\begin{tikzpicture}[x=1in,y=1in]
  \def\lw{0.98in}\def\nw{1.12in}\def\tw{0.92in}\def\gw{1.28in}
  \node[font=\tiny\bfseries] at (0.5,0.02) {Function};
  \node[font=\tiny\bfseries] at (1.66,0.02) {Construction now};
  \node[font=\tiny\bfseries] at (2.83,0.02) {Trigger (first of)};
  \node[font=\tiny\bfseries] at (4.15,0.02) {Target};
  \foreach \y/\f/\n/\t/\g/\st in {
     -0.32/{Transport and share encryption}/{Hybrid X25519 + ML-KEM-768 (post-quantum now)}/{HQC standard (expected 2027)}/{add HQC as backup KEM}/dcrypto,
     -0.80/{Council signatures, receipts, CapCerts}/{Hybrid Ed25519 + ML\mbox{-}DSA\mbox{-}65, DPEP format}/{composite RFC; FIPS~206}/{composite format; FN-DSA option}/dbox,
     -1.28/{Layer-1 validator keys}/{engine's suite (Ed25519 or ML-DSA-65 keys)}/{composite standard and engine support}/{hybrid validator keys by key rotation}/dbox,
     -1.76/{One-report-per-person credentials}/{classical, threshold-issued, per epoch}/{post-quantum credential fits a 150~kB round}/{lattice credential, re-issued}/dbox,
     -2.24/{Sortition beacon}/{threshold BLS run by the federations}/{post-quantum VRF or VDF standard}/{post-quantum beacon, key re-shared}/dbox,
     -2.72/{Root anchors and history}/{SLH\mbox{-}DSA\mbox{-}SHA2\mbox{-}192s\\(hash-based)}/{none}/{unchanged}/dcrypto}{
    \node[font=\tiny\bfseries, text width=\lw, align=flush left] (f) at (0.5,\y) {\f};
    \node[\st, text width=\nw, minimum height=0.32in] (n) at (1.66,\y) {\n};
    \node[dwarn, dashed, text width=\tw, minimum height=0.32in] (t) at (2.83,\y) {\t};
    \node[\st, text width=\gw, minimum height=0.32in] (g) at (4.15,\y) {\g};
    \draw[darr] (n)--(t); \draw[darr] (t)--(g);
  }
  \node[font=\tiny\itshape, text=CrimsonAccent] at (2.45,-3.08) {Any lane switches early on a credible public warning of progress towards a CRQC.};
  \node[dhead, text width=1.0in] (m1) at (0.56,-3.5) {1\; ratify new suite floor};
  \node[dhead, text width=1.0in] (m2) at (1.85,-3.5) {2\; dual-sign window (old + new)};
  \node[dhead, text width=1.0in] (m3) at (3.14,-3.5) {3\; re-anchor history under SLH-DSA};
  \node[dhead, text width=1.0in] (m4) at (4.33,-3.5) {4\; withdraw old suite};
  \draw[darr] (m1)--(m2); \draw[darr] (m2)--(m3); \draw[darr] (m3)--(m4);
  \node[font=\tiny, text width=4.7in, align=flush center] at (2.45,-3.88) {Every migration follows these four steps. \texttt{CryptoSuiteId} sits inside the signed bytes; objects below the floor are rejected.};
\end{tikzpicture}
\caption{Crypto-agility lanes. Each function moves from its current construction to its target when the first trigger occurs, through the same four-step migration.}
\label{fig:agility}
\end{figure}


\paragraph{Credentials.}
Threshold issuance, in which no single council can mint or link a credential, follows Coconut (Sonnino et al.\ 2019), a classical pairing-based construction. Lattice alternatives exist as research (Bootle et al.\ 2023), but reported sizes, under 80~kB per credential in one recent construction,\footnote{Sven Argo et al., ``Practical Post-Quantum Signatures for Privacy,'' \emph{ACM CCS 2024}, \url{https://doi.org/10.1145/3658644.3670297} (IACR ePrint 2024/131).} do not fit beside a 55~kB report in a 150~kB round. \textbf{Selection criterion for the interim scheme:} its unlinkability should not rest on an assumption a CRQC breaks (statistically zero-knowledge presentation proofs, statistical blinding). This must be verified for the chosen scheme; it is not a property of every classical credential. RFC~9474, for example, attributes the blindness of RSA blind signatures to a separate proof without stating its assumptions, and states that the protocol is not post-quantum ready. If no scheme meets the criterion, a future CRQC could link past credential showings to identities. That reveals who submitted, not what was submitted, since report contents stay secret-shared under ML-KEM.

\paragraph{Beacon.}
The interim beacon uses threshold-BLS randomness of the kind produced by the open-source \texttt{drand} software,\footnote{\url{https://drand.love} (checked November 2026).} with the threshold key shared among the federations rather than delegated to an outside network. It has no last-revealer problem, and a CRQC could bias only draws made after it exists. Its distributed key generation and share custody are not yet specified.

\paragraph{Layer 1.}
CometBFT remains the reference engine (provisional). Its v0.40 release accepts ML-DSA-65 consensus keys through key rotation, as a pure rather than hybrid scheme, and the Cosmos documentation estimates about 331~kB of signatures per block for 100 validators, consistent with the 290--296~kB for 86 hybrid signatures in Table~\ref{tab:light-clients}.\footnote{Cosmos SDK documentation, ``Post-quantum keys'' (requires Cosmos SDK 0.55 and CometBFT 0.40), \url{https://docs.cosmos.network/sdk/latest/keys/post-quantum-keys}; release dates not confirmed (checked November 2026).} The validator suite therefore moves to ML-DSA by key rotation and to the hybrid when a composite standard and engine support exist, without a change of engine. A Rust engine with a pluggable signing interface (Malachite) is the fallback if a hybrid consensus signature becomes mandatory earlier; it is alpha software without an external audit.

\section{Democratic Membership, Sybil Resistance, and Civil Liberties}
\label{sec:membership}

\begin{designrule}[Four-Way Separation of Identity, Eligibility, Accountability, and Legitimacy]
DPEP strictly separates: (1) \textbf{Cryptographic Identity} (control of a keypair), (2) \textbf{Institutional Eligibility} (accredited membership in a council), (3) \textbf{Economic Accountability} (correspondence of reported claims to physical reality), and (4) \textbf{Democratic Legitimacy} (self-management proportional to degree affected).
\end{designrule}

Figure~\ref{fig:identity} shows how the separation is carried into membership, delegation and private reporting.

\begin{figure}[htbp]
\centering
\begin{tikzpicture}[x=1in,y=1in]
  \foreach \x/\t in {0.6/{\textbf{1 Identity}\\control of a keypair},
                     1.83/{\textbf{2 Eligibility}\\accredited council membership},
                     3.07/{\textbf{3 Accountability}\\claims match physical reality},
                     4.3/{\textbf{4 Legitimacy}\\say in proportion to how one is affected}}{
    \node[dhead, text width=1.08in, minimum height=0.36in, font=\tiny] at (\x,0) {\t};
  }
  \node[font=\tiny\itshape, text=SlateDark] at (2.45,-0.32) {Kept separate by design (Section \ref{sec:membership}): holding a key never by itself grants eligibility, budget or voting weight.};
  % left branch
  \node[font=\tiny\bfseries] at (1.2,-0.62) {Membership and delegation};
  \node[dbox, text width=1.9in] (l1) at (1.2,-0.95) {Person generates a hybrid keypair on the phone};
  \node[dbox, text width=1.9in] (l2) at (1.2,-1.42) {Council accredits the member by a $k$-of-$n$ multi-signature};
  \node[dbox, text width=1.9in] (l3) at (1.2,-1.89) {CapCert: membership certificate (signed under the active \texttt{CryptoSuiteId})};
  \node[dbox, text width=1.9in] (l4) at (1.2,-2.36) {CapCert: time-bounded delegate certificate for a federation, revocable by one council vote};
  \draw[darr] (l1)--(l2); \draw[darr] (l2)--(l3); \draw[darr] (l3)--(l4);
  % right branch
  \node[font=\tiny\bfseries] at (3.7,-0.62) {Private consumption reports};
  \node[dcrypto, text width=1.9in] (r1) at (3.7,-0.95) {$k$ of $n$ councils jointly issue a per-epoch anonymous budget credential (threshold issuance; interim classical)};
  \node[dcrypto, text width=1.9in] (r2) at (3.7,-1.5) {Credential shown to the aggregators at upload with a per-epoch nullifier; no identity revealed};
  \node[dcrypto, text width=1.9in] (r3) at (3.7,-2.0) {\texttt{Prio3} report bound to the credential through its nonce and \texttt{ctx}};
  \node[dwarn, text width=1.9in] (r4) at (3.7,-2.42) {Open: binding the attested budget $B_u$ inside the report circuit};
  \draw[darr] (r1)--(r2); \draw[darr] (r2)--(r3); \draw[darr, CrimsonAccent, dashed] (r3)--(r4);
  \draw[darr, dashed, RuleGray] (l2.east) -- node[dlab, above, pos=0.5]{accreditation} (r1.west);
  \node[dbox, text width=4.72in] at (2.45,-2.9) {\textbf{Goal:} residency, workplace and consumption records cannot be linked to one another, and no member's report or budget is visible to another (design goal; Section \ref{sec:membership}).};
\end{tikzpicture}
\caption{Identity, eligibility and the two credential paths. Membership and delegation are attributed and revocable; private consumption reports use a separate, unlinkable credential whose budget binding is still open.}
\label{fig:identity}
\end{figure}

While phantom digital accounts in $CC_h$ cannot mint adult work income without $WC_j$ effort certificates capped by $SB_j/SC_j$, four deeper Sybil and collusion vectors must be closed explicitly:
\begin{enumerate}[leftmargin=1.8em,itemsep=2pt]
  \item \textbf{Closing Workplace Hours Inflation (\texttt{INV-4B}):} Capping only the \emph{mean} effort rating $\bar{\epsilon}_j$ fails if a $WC_j$ fabricates phantom workers to inflate total hours $\sum h_u$. Equation (\ref{eq:effort-cap}) closes this by capping \emph{total credited effort-hours} $\sum_{u \in WC_j} h_u \epsilon_u \le \sum_\ell (l^0_{j,\ell}+\Delta l_{j,\ell}) \min(\bar{\epsilon}_{\max}, SB_j/SC_j)$ (accredited hours) and flooring labour user-right charges at $p^{\text{lab}}_\ell \ge w_{\text{base}}$.
  \item \textbf{Detecting Collusive ``Carousel Fraud'' Between Fake Producer and Consumer Councils:} If a colluding $WC_j$ and $CC_h$ sign fictitious delivery receipts for non-existent services to mint effort credits, cryptography alone cannot detect the lie. DPEP pairs automated input-output ratio anomaly flags (detecting closed-loop service claims without corresponding physical/intermediate inputs) with mandatory \textbf{Industry Federation physical \emph{euthynai} spot-audits} prior to accrediting capacity.
  \item \textbf{Preventing Vote \& Sortition Dilution:} Sortition tickets and federation referenda are weighted by \textbf{cross-audited council accreditation quotas} rather than raw public keys.
  \item \textbf{Civil-Liberties Protection Against a Coercive Population Register (design goal, not yet constructed):} To prevent housing (\texttt{HFB}) and workplace (\texttt{EFB}) registries from forming a surveillance database or trapping victims of domestic coercion, the design \emph{aims} to let each person turn council attestations into unlinkable per-epoch budget credentials, separating residency from workplace and consumption. No such construction is specified (Section~\ref{sec:bu-cred}), so this protection does not yet exist. Budgets and reports are \textbf{individual}: the index $u$ in Section~\ref{sec:prio3} is one budget-holding person, never a household head. People who share a household and wish to plan jointly each submit their own report against their own budget; no member's report or budget is visible to, or controlled by, another. Unaffiliated migrants and transients receive unconditional Maxim-4 baseline provisioning through community clinics.
\end{enumerate}

\section{Private Household Aggregation via \texttt{Prio3} VDAFs}
\label{sec:prio3}

A discarded design in which each household $u$ generates a standalone zero-knowledge budget proof $\Pi_u$ on a commitment $C_u$ while separately submitting a pairwise-masked vector $\tilde{y}_u$ has three flaws: (1) nothing binds $C_u$ to $\tilde{y}_u$; (2) if even one of many neighbours drops offline, pairwise masks fail to cancel; (3) aggressive FRI/STIR/WHIR parameterisations rely on Reed--Solomon proximity-gap conjectures that Crites and Stewart (IACR ePrint 2025/2046) disproved at up-to-capacity rates.\footnote{Elizabeth Crites and Alistair Stewart, ``On Reed--Solomon Proximity Gaps Conjectures,'' IACR Cryptology ePrint Archive \texttt{2025/2046} (December 2025), \url{https://eprint.iacr.org/2025/2046.pdf}.} Those constructions are not part of Architecture D.

\paragraph{Catalogues.}
$\Omega$ with $G = |\Omega| = 10{,}000$ is the full priced catalogue, and $\mathcal{G}^{\text{priv}}$ with $G_{\text{priv}} = 3{,}000$ the private-consumption categories at national scale. Household reports need not use $\mathcal{G}^{\text{priv}}$ directly. Section~\ref{sec:privacy-baseline} proposes a coarser \emph{reporting} catalogue $\mathcal{G}^{\text{rep}}$ of 100--250 categories, consistent with REQ-DATA-02.

\paragraph{VDAF construction and what it checks.}
Each budget holder $u$ (one person; Section~7.3) encodes its request as a vector $q_u \in \{0,\dots,2^{b}-1\}^{|\mathcal{G}^{\text{rep}}|}$ in Household Reporting Units (Section~\ref{sec:prio3-cost}) and splits it into $K=3$ additive Field128 shares, encrypted with hybrid \texttt{X25519MLKEM768} to three aggregators. Standard \texttt{Prio3SumVec} proves only that every coordinate lies in $[0, 2^b-1]$ (VDAF-22 \S 7.4.3). The budget predicate
\begin{equation}
  \sum_{g \in \mathcal{G}^{\text{rep}}} \pi_g^{(t)}\, q_{u,g} + \text{PubShare}_u \le B_u, \qquad \pi_g^{(t)} = p_g^{(t)} s_g \label{eq:budget-circuit}
\end{equation}
needs a \textbf{custom validity circuit}, which is not specified. Two facts constrain it. First, the prices $\pi^{(t)}$ must enter the circuit as public constants, so each price vector defines a different VDAF instance: a report proves the predicate for the round's prices and for no other. The application context string \texttt{ctx} should bind the report to $(\text{epoch}, t, \text{BLAKE3}(P^{(t)}), \text{aggregator set})$, but \texttt{ctx} enters only the domain separation of the XOF, not the circuit's arithmetic. Second, $B_u$ must reach the circuit either as a public value (which identifies the person) or as a secret-shared input whose correctness is vouched for by a credential; the second option depends on the open credential problem (Section~\ref{sec:bu-cred}). Reference: IETF CFRG \texttt{draft-irtf-cfrg-vdaf-22} (14 August 2026; Informational), full text.\footnote{Henry Corrigan-Gibbs and Dan Boneh, ``Prio: Private, Robust, and Scalable Computation of Aggregate Statistics,'' \emph{USENIX NSDI '17}, pp.~259--282; Dan Boneh et al., ``Zero-Knowledge Proofs on Secret-Shared Data via Fully Linear PCPs,'' \emph{CRYPTO 2019}; Richard Barnes et al., \emph{Verifiable Distributed Aggregation Functions (VDAFs)}, IETF CFRG \texttt{draft-irtf-cfrg-vdaf-22} (14 August 2026), \url{https://datatracker.ietf.org/doc/html/draft-irtf-cfrg-vdaf-22}. Security analysis: Hannah Davis, Christopher Patton, Mike Rosulek, and Phillipp Schoppmann, ``Verifiable Distributed Aggregation Functions,'' \emph{PoPETs} 2023, no.~4, pp.~578--592; cite the full version (IACR ePrint 2023/130), whose 2024 revision corrects the proof of Theorem~1 with a weaker bound and a larger PRG seed.}

\paragraph{Privacy versus correctness (VDAF-22 \S 9).}
These must not be collapsed. Privacy of a measurement requires \textbf{at least one honest aggregator}. Correctness of the aggregate requires \textbf{all aggregators}: ``An Aggregator can, without detection, manipulate the aggregate result by modifying its own aggregate share.'' Any one aggregator can inflate or deflate demand and so move prices.

Figure~\ref{fig:privacy-flow} shows the asymmetry.

\begin{figure}[htbp]
\centering
\begin{tikzpicture}[x=1in,y=1in]
  \node[dbox, text width=0.82in] (p) at (0.47,1.6) {\textbf{Person's report}\\$q_u$: 200 categories\\$\times$ 16 bits\\+ credential and nullifier};
  \node[dcrypto, text width=0.92in] (a) at (2.0,2.5) {\textbf{Aggregator A}\\consumer-side operator};
  \node[dcrypto, text width=0.92in] (b) at (2.0,1.6) {\textbf{Aggregator B}\\producer-side operator};
  \node[dcrypto, text width=0.92in] (c) at (2.0,0.7) {\textbf{Aggregator C}\\external operator};
  \draw[darr, TealAccent] (p.east) -- node[dlab, sloped, above]{share 1} (a.west);
  \draw[darr, TealAccent] (p.east) -- node[dlab, above]{share 2} (b.west);
  \draw[darr, TealAccent] (p.east) -- node[dlab, sloped, below]{share 3} (c.west);
  \node[dbox, text width=0.72in] (m) at (3.33,1.6) {aggregate shares\\+ published\\commitments};
  \draw[darr] (a.east) -- (m.north west);
  \draw[darr] (b.east) -- (m.west);
  \draw[darr] (c.east) -- (m.south west);
  \node[dbox, text width=0.7in] (t) at (4.43,1.6) {\textbf{Neighbourhood total}\\($\ge 30$ reports)};
  \draw[darr] (m) -- (t);
  \node[dinst, text width=0.7in] (chk) at (4.43,0.7) {public check\\$p\cdot\textstyle\sum y \le \sum B_u$};
  \draw[darr] (t) -- (chk);
  \node[dinst, text width=1.5in] (aud) at (3.75,2.62) {\textbf{Independent auditor}: canary batches with known totals, sent to all three};
  \draw[darr, GoldAccent!80!black, dashed] (aud.west) -- (a.east);
  \node[dcrypto, text width=4.72in, align=flush left] at (2.45,0.22) {\textbf{Privacy (cryptographic):} each share alone reveals nothing; the report stays private if \emph{at least one} of A, B, C is honest. Shares are encrypted with hybrid ML-KEM.};
  \node[dwarn, text width=4.72in, align=flush left] at (2.45,-0.2) {\textbf{Correct totals (institutional):} need \emph{all three} honest. Any one aggregator can change its own aggregate share and shift the total without detection by cryptography. Detection only: the public check (gross overspend), commitments (attribution) and canary batches (indiscriminate tampering).};
\end{tikzpicture}
\caption{Household report flow and the privacy/correctness asymmetry. Privacy needs one honest aggregator; correct totals need all three, and the remaining safeguards detect rather than prevent manipulation.}
\label{fig:privacy-flow}
\end{figure}

\paragraph{How much the public budget check limits a cheating aggregator.}
The check $p\cdot\sum_u y_u \le \sum_u B_u$ is cheap and should be run, but it constrains only one number. A malicious aggregator can (i) shift demand between categories while keeping the value total fixed, which is undetectable and moves relative prices; (ii) inflate the total up to the slack between actual spending and $\sum_u B_u$ (if households leave 10\% of budgets unspent, about 10\% of value can be added undetected); and (iii) deflate demand without limit, since the inequality has no lower bound. It is a tripwire for gross overspend, not a correctness guarantee.

\begin{table}[htbp]
\centering
\caption{Aggregator-correctness options (experimental except the first row)}
\label{tab:agg-options}
\scriptsize
\begin{tabularx}{\textwidth}{L{1.30in} L{1.40in} Y}
\toprule
\textbf{Option} & \textbf{Guarantee} & \textbf{Cost / institutional demand} \\
\midrule
Status-quo \texttt{Prio3} ($K=3$), nested CC / Ward / Federation & Privacy if $\ge 1$ honest; correct totals only if all honest. & Cheapest; but the three organs belong to one federation, so their independence is doubtful. \\
Non-nested aggregators ($K=3$) & Same guarantees; independence more plausible. & E.g.\ consumer federation, an industry federation, and an external civil-society or academic operator. \\
Replicated aggregators ($K'>3$) & Privacy if one honest; correctness still needs all shares. & More KEMs and operational load; no correctness gain without robust aggregation. \\
Robust / maliciously secure MPC & Detects or bounds a cheating aggregator. & Super-linear communication; not \texttt{Prio3}; unaudited for DPEP. \\
Public check $p\cdot\sum y \le \sum B_u$ & Detects gross aggregate overspend only (see text). & Cheap; complements every option. \\
Canary batches & Detects indiscriminate tampering with aggregate shares, with probability rising with the canary share. & An independent auditor submits whole batches with known totals. Needs per-batch aggregate shares and batches the aggregators cannot tell apart; does not catch tampering aimed at batches known to be real. \\
Committee-run aggregation & Ties aggregation to the $n=128$ BFT set. & Concentrates privacy-breaking power in the body that finalises prices. \\
Differential privacy on released totals & Bounds one-household differencing; outside VDAF. & Bias against planning accuracy; parameters unspecified. \\
\bottomrule
\end{tabularx}
\end{table}

\paragraph{Recommended correctness model for Phase~3 (proposed, for author ratification).}
Use \textbf{three non-nested aggregators run by organisations with different interests} (for a co-op pilot: the consumer side, the producer side, and an external operator), plus the public budget check and published commitments to each aggregator's aggregate share, so that a dispute can at least be attributed. Nested CC/Ward/Federation aggregators should \textbf{not} be the default. Add \textbf{canary batches} run by an auditor independent of all three aggregators. \textbf{Cryptography in this model guarantees privacy, not correct totals:} correctness rests on institutional independence plus detection, and no mature construction removes that trust while keeping \texttt{Prio3}'s privacy. Robust MPC is deferred until a pilot shows that aggregator manipulation is a real rather than a hypothetical threat. Trusted execution environments are rejected for this role, because they would make the chip vendor the party trusted for correctness. This choice cannot be made final before Phase~3 because it depends on which organisations exist in a pilot.

\subsection{Report Sizes, Units, and Measured Costs}
\label{sec:prio3-cost}

\paragraph{Household Reporting Unit (HRU).}
At the protocol's fixed $10^{-6}$ scale, a $b$-bit coordinate can express at most $(2^b-1)/10^6$ units: 1.05 units at 20 bits, which is useless, and a full 64-bit field makes reports enormous. Household reports therefore use a per-category \textbf{Household Reporting Unit}: the catalogue fixes for each reporting category $g$ an integer scale $s_g$ (micro-units per HRU, for example 100{,}000 for 0.1~kg), and a report carries $q_{u,g} \in \{0,\dots,2^{16}-1\}$. The aggregate in protocol units is
\[
  D_g^{\text{micro}} = s_g \sum_u q_{u,g},
\]
which is exact integer arithmetic; no rounding occurs after the report is made. Rounding occurs once, when a person's intended quantity is rounded to whole HRUs (round-half-to-even): at most half an HRU per category per person, so at most $N s_g/2$ micro-units in aggregate and, for independent errors, a standard deviation of $s_g\sqrt{N/12}$ (for $s_g$ = 0.1 unit and $N=10^6$, about 29 units against a total in the hundreds of thousands). Sixteen bits gives 0--65,535 HRU per category per report; categories whose range does not fit choose a larger $s_g$.

\begin{table}[htbp]
\centering
\caption{\texttt{Prio3SumVec} with 3 aggregators: measured sizes and desktop timings (libprio 0.18.1)}
\label{tab:prio3-bench}
\scriptsize
\setlength{\tabcolsep}{3pt}
\begin{tabular}{rrrrrrrl}
\toprule
\textbf{Length} & \textbf{Bits} & \textbf{Leader share} & \textbf{Client upload} & \textbf{Shard} & \textbf{Verify (leader)} & \textbf{Verify (helper)} & \textbf{Fits 75 kB?} \\
\midrule
100 & 20 & 35.1 kB & 35.3 kB & 0.55 ms & 0.21 ms & 0.23 ms & Yes \\
100 & 32 & 54.9 kB & 55.1 kB & 0.65 ms & 0.21 ms & 0.22 ms & Yes \\
100 & 64 & 107.7 kB & 108.0 kB & 1.13 ms & 0.30 ms & 0.32 ms & No \\
230 & 20 & 78.0 kB & 78.3 kB & 0.81 ms & 0.23 ms & 0.25 ms & No \\
230 & 32 & 123.6 kB & 123.8 kB & 1.28 ms & 0.33 ms & 0.35 ms & No \\
230 & 64 & 243.3 kB & 243.6 kB & 2.52 ms & 0.82 ms & 0.86 ms & No \\
3,000 & 20 & 975.8 kB & 976.0 kB & 10.6 ms & 3.46 ms & 3.64 ms & No \\
3,000 & 32 & 1,556.3 kB & 1,556.5 kB & 16.2 ms & 4.85 ms & 5.21 ms & No \\
3,000 & 64 & 3,100.4 kB & 3,100.7 kB & 33.3 ms & 12.0 ms & 12.7 ms & No \\
\midrule
\multicolumn{8}{l}{\emph{Computed with the same size formula (exact on all nine measured rows); not timed}} \\
150 & 16 & 41.7 kB & 41.9 kB & & & & Yes \\
200 & 16 & 54.9 kB & 55.1 kB & & & & Yes \\
250 & 16 & 68.1 kB & 68.3 kB & & & & Yes \\
10,000 & 64 & 10,292.8 kB & 10,293.0 kB & & & & No \\
\bottomrule
\end{tabular}
\par\smallskip{\scriptsize\raggedright Medians of 30 runs, one core, Apple M4 Pro, macOS 26.6, rustc 1.98.1, \texttt{--release}, \texttt{chunk\_length} chosen by libprio's \texttt{optimal\_chunk\_length}. ``Verify'' is \texttt{verify\_init} (proof-share computation); combining verifier shares and \texttt{verify\_next} take under 5~$\mu$s. Helper shares are 64~B seeds; the public share is 96~B. libprio 0.18.1 implements VDAF draft~18, not draft~22; the encoded sizes follow the same \texttt{Prio3SumVec} layout. These are desktop figures: phones and Raspberry Pi boards will be slower by an unmeasured factor. Code and commands: \path{research/analysis/bench_prio3/}.\par}
\end{table}

\paragraph{What the measurements change.}
A dense 3,000-category report is about 1~MB, thirteen times the 75~kB budget. A 50-coordinate ``sparse'' report cannot be used to send 50 items chosen from 3,000: \texttt{Prio3SumVec} sums fixed coordinates, so a short vector either covers only 50 fixed categories or reveals which categories each coordinate stands for, the same leak that dense reports exist to prevent. No private-sparse construction is specified here; any such design would be experimental. Proof checking is not negligible either: verifying one dense report costs each aggregator, helper as well as leader, about 3.5~ms on a fast desktop core, so $10^6$ reports cost about 61 core-minutes per aggregator per fresh-report round (3.8 minutes on 16 cores), before network transfer and on hardware much faster than a neighbourhood Raspberry Pi. A coarse 200-category, 16-bit report costs about 0.25~ms, roughly 4 core-minutes per aggregator at $10^6$, and a neighbourhood of 200 people verifies in about 50~ms of desktop time. The illustrative ``1~ms per report'' used before measurement was optimistic for dense reports even on a desktop.

\subsection{A Working Privacy Baseline (proposed)}
\label{sec:privacy-baseline}

\begin{table}[htbp]
\centering
\caption{Privacy baselines compared at $N=10^6$ (sizes per fresh report; CPU per aggregator per fresh-report round)}
\label{tab:privacy-options}
\scriptsize
\setlength{\tabcolsep}{3pt}
\begin{tabularx}{\textwidth}{L{1.05in} L{0.62in} L{0.70in} L{1.05in} Y}
\toprule
\textbf{Option} & \textbf{Upload} & \textbf{Agg.\ CPU (desktop)} & \textbf{Hidden from councils} & \textbf{Costs and risks} \\
\midrule
(i) Dense \texttt{Prio3}, 3,000 categories, 20 bits & 976 kB & $\approx 61$ core-min & Every category choice & Exceeds mobile budgets; 20 bits too narrow for many goods. \\
(ii) Coarse \texttt{Prio3}, 100--250 categories, 16-bit HRU & 28--68 kB & $\approx 4$ core-min & Choice among coarse categories; not the product within a category & Loses detail; more reliance on disaggregation to products (P5) and on local intra-year fulfilment; Hayek's objection bites harder. \\
(iii) Council-level trust, no household cryptography & Plain request to council tallier & Negligible & Bundle hidden from neighbours by policy only; tallier sees it & Trusted tallier can leak or alter; fine for small consenting pilots with data-protection safeguards. \\
(iv) Pilot-selectable modes & (ii) or (iii) & -- & As chosen & More code paths to audit; mode choice must itself be democratic and disclosed. \\
\bottomrule
\end{tabularx}
\end{table}

\paragraph{Recommendation (proposed, for author ratification).}
Adopt \textbf{option (ii) with $|\mathcal{G}^{\text{rep}}| = 200$ and 16-bit HRUs} (55~kB upload) as the working baseline for household cryptography, and \textbf{option (iii)} for Phase~1H and the first Phase~3 cooperative pilots. Household-level cryptography is therefore \textbf{not} required for early pilots: with 10--30 consenting cooperatives a disclosed tallier, data minimisation and a data-protection impact assessment give adequate protection at a fraction of the audit cost, and the pilot can test whether members value bundle secrecy enough to justify option (ii). Option (i) is not a baseline at any scale considered here. The coarse baseline ties three open questions together: P5 (how well coarse totals disaggregate to products, now with 200 rather than 3,000 reported categories), P13 (the minimum batch $m_{\min}=30$ now applies per coarse category, which makes small-hamlet anonymity easier), and Kill Gate~3's bandwidth and verification caps (Table~\ref{tab:gates}). The resulting mobile budget per fresh-report round is about 55~kB up plus the light-client download in Section~\ref{sec:sortition}.

\subsection{Prices, Carry-Forward, and Replay}
\label{sec:carry-forward}

\paragraph{What an offline household can prove.}
Eq.~(\ref{eq:budget-circuit}) uses the round's prices. A report made at round $t$ proves the budget predicate at $P^{(t)}$ and nothing about $P^{(t+1)}$; scaling an old report proportionally does not create a new proof. VDAF-22 adds a harder limit: ``Prio3 only permits a report to be aggregated once'' (\S 7.2.3), because re-aggregation can reuse correlated randomness and compromise confidentiality (\S 9.4). Carrying an individual report forward into a later round's batch is therefore not available as specified.

\begin{table}[htbp]
\centering
\caption{Handling households between fresh reports}
\label{tab:carry-options}
\scriptsize
\begin{tabularx}{\textwidth}{L{1.30in} L{1.30in} Y}
\toprule
\textbf{Option} & \textbf{What is verified} & \textbf{Cost} \\
\midrule
(a) Fresh report every round & Budget at every round's prices & 10 uploads and verifications per year; excludes anyone offline in a round. \\
(b) Fresh reports at milestone rounds only; aggregate carried between & Budget at milestone prices; between milestones, only the aggregate check at current prices & Household demand frozen between milestones; councils adjust at council level. \\
(c) Bounded overspend & Aggregate $p^{(t)}\cdot\sum y \le (1+\kappa)\sum B_u$ & Individual overspend between milestones is not detected. \\
(d) Exclude stale reports & Nothing stale & Biases totals against people with poor connectivity. \\
\bottomrule
\end{tabularx}
\end{table}

\paragraph{Proposed rule (for author ratification).}
Combine (b) and (c). Households submit fresh, budget-checked reports at the three milestone rounds of the hybrid schedule ($t=1,4,8$; P3). Between milestones, each neighbourhood's demand is held at its last milestone \emph{aggregate}: aggregators reuse their aggregate shares as constants, with no individual report re-verified or re-aggregated. The public aggregate check is applied at current prices with a tolerance $\kappa$ (proposed 5\%), and price-driven revisions between milestones happen at council level. A person offline at a milestone may submit within a grace window; otherwise they are counted at their council's per-capita average, flagged as imputed. Each released total at a milestone is a fresh aggregate of at least $m_{\min}=30$ reports, which limits but does not remove differencing between milestones (minimum batches, coarser release and DP noise remain the mitigations). This raises the per-person upload to three fresh reports a year: about 0.17~MB at the coarse baseline, against 9.8~MB for ten dense reports.

Figure~\ref{fig:calendar} places the milestone rounds in the planning season and operating year.

\begin{figure}[htbp]
\centering
\begin{tikzpicture}[x=1in,y=1in]
  \node[font=\tiny\bfseries, anchor=west, text width=0.75in, align=flush left] at (0,0) {Planning season\\\normalfont(illustrative 10-round plan)};
  \foreach \t in {1,...,10}{
    \pgfmathsetmacro{\x}{0.95+(\t-1)*0.31}
    \ifnum\t=1 \def\st{dwarn}\else\ifnum\t=4 \def\st{dwarn}\else\ifnum\t=8 \def\st{dwarn}\else\def\st{dcrypto}\fi\fi\fi
    \node[\st, minimum width=0.24in, minimum height=0.22in, font=\tiny\bfseries] (r\t) at (\x,0) {\t};
  }
  \foreach \t/\m in {1/M1, 4/M2, 8/M3}{ \node[font=\tiny\bfseries, text=CrimsonAccent] at ($(r\t)+(0,0.22)$) {\m}; }
  \node[dbox, text width=0.82in] (syn) at (4.42,0) {S5--S6 synthesis\\and ratification};
  \draw[darr] (r10)--(syn);
  \node[font=\tiny, anchor=north west, align=left] at (0.95,-0.2) {round $t$ \quad (one Layer-1 checkpoint each)};
  \node[dwarn, text width=2.25in, align=flush left] at (1.25,-0.62) {\textbf{Milestone round ($t = 1, 4, 8$):} people deliberate; fresh, budget-checked \texttt{Prio3} reports; phones run a full quorum check.};
  \node[dcrypto, text width=2.25in, align=flush left] at (3.68,-0.62) {\textbf{Automated round:} councils revise at council level; household demand held at the last milestone \emph{aggregate}; aggregate check within $\kappa = 5\%$.};
  \node[font=\tiny\bfseries, anchor=west, text width=0.75in] at (0,-1.3) {Operating year};
  \node[dbox, text width=2.2in, minimum height=0.3in] (s7) at (2.05,-1.3) {S7 intra-year clearing: one daily rollup block (365 per year)};
  \node[dinst, text width=0.66in, minimum height=0.3in] (s8) at (3.65,-1.3) {S8 scoped replans ($\approx 50$)};
  \node[dbox, text width=0.55in, minimum height=0.3in] (s9) at (4.45,-1.3) {S9 audit};
  \draw[darr] (s7)--(s8); \draw[darr] (s8)--(s9);
  \node[dhead, text width=4.72in] at (2.45,-1.75) {Layer-1 blocks per year: 10 round checkpoints + 365 daily rollups + 50 scoped replans $=$ 425 (435 at 20 rounds)};
\end{tikzpicture}
\caption{Annual calendar for an illustrative 10-round plan: three human milestone rounds with fresh reports, automated rounds in between, then the operating year and its Layer-1 block count.}
\label{fig:calendar}
\end{figure}

\paragraph{Replay, re-runs and retention.}
Aggregators keep the nonce of every accepted report for the epoch and reject duplicates (VDAF-22 \S 9.2 recommends the nonce for replay protection); \texttt{ctx} binds epoch, round, price-vector hash and aggregator set, so a report cannot be replayed into another round or instance. A round \emph{cannot} be re-run from stored reports with a replacement aggregator: the stored shares are encrypted to the failed aggregator, and re-verifying them would aggregate them a second time. Recovery therefore means failing that milestone and requesting fresh reports, or carrying the previous milestone's aggregate. Aggregators retain encrypted input shares, nonces and aggregate-share commitments only until the challenge window $T_{\text{challenge}}$ closes (proposed 14 days), to settle inclusion disputes, and then delete them; that bounded retention is what ``ephemeral shares'' means in Section~11.2. People keep their own reports locally.

\subsection{Budget Credentials and One-Report-Per-Person}
\label{sec:bu-cred}
The protocol needs each eligible person to submit at most one report per milestone, against their own warranted budget, without a register linking reports to identities. No construction is specified, and Table~\ref{tab:pqc-bench} lists no post-quantum blind or anonymous credential. VDAF-22 notes that clients may be anonymous and that Sybil and rate-limit enforcement is outside the VDAF. Three points narrow the problem. (1) BBS-style anonymous credentials are pairing-based and offer no post-quantum security; lattice-based anonymous credentials are an active research area without a standard, and none is named here. (2) Checking the credential \emph{inside} the FLP is not realistic: the FLP validates low-degree arithmetic over the secret-shared measurement, while verifying a signature or hash chain needs a very large circuit, and the credential itself would have to be secret-shared. The credential should instead be shown to the aggregators at upload, outside the VDAF, with a per-epoch nullifier to block double submission, and bound to the report through its nonce and \texttt{ctx}. (3) What remains open is binding a secret-shared $B_u$ inside Eq.~(\ref{eq:budget-circuit}) to the credential's attested budget. Until that is solved, one-report-per-person rests on council accreditation and rate limits, which is why option (iii) above is the early-pilot default. Section~\ref{sec:pq-urgency} adopts a classical credential with threshold issuance as the interim construction for the one-report-per-person part; the binding of $B_u$ remains open.

\section{Threat Model, Equivocation, Recovery, and Ballots}
\label{sec:threat-model}
Figure~\ref{fig:trust-map} summarises what each part of the design rests on; the list below gives the detail.

\begin{figure}[htbp]
\centering
\begin{tikzpicture}[x=1in,y=1in, node distance=0.06in]
  \node[dcrypto, font=\tiny\bfseries, text width=1.45in, minimum height=0.36in] (h1) at (0.78,0) {Cryptography guarantees\\(under stated assumptions)};
  \node[dinst, font=\tiny\bfseries, text width=1.45in, minimum height=0.36in] (h2) at (2.45,0) {Detected, not prevented};
  \node[dwarn, font=\tiny\bfseries, text width=1.45in, minimum height=0.36in] (h3) at (4.12,0) {Rests on institutions\\and physical audit};
  % column 1
  \node[dcrypto, text width=1.45in, below=of h1] (a1) {Confidentiality of reports and channels (hybrid ML-KEM)};
  \node[dcrypto, text width=1.45in, below=of a1] (a2) {Authenticity of proposals, receipts, checkpoints (hybrid signatures)};
  \node[dcrypto, text width=1.45in, below=of a2] (a3) {Inclusion of a council's proposal (Two-Root path), given the leaf archive};
  \node[dcrypto, text width=1.45in, below=of a3] (a4) {Privacy of a household report if $\ge 1$ aggregator is honest (\texttt{Prio3})};
  \node[dcrypto, text width=1.45in, below=of a4] (a5) {Same prices on every node (deterministic fixed-point WASM)};
  % column 2
  \node[dinst, text width=1.45in, below=of h2] (b1) {Omitted or tampered proposals (receipts + challenge window)};
  \node[dinst, text width=1.45in, below=of b1] (b2) {BFT equivocation (if $\ge 1$ honest validator records both votes)};
  \node[dinst, text width=1.45in, below=of b2] (b3) {Gross aggregate overspend (public budget check)};
  \node[dinst, text width=1.45in, below=of b3] (b4) {Indiscriminate tampering with aggregate shares (canary batches)};
  \node[dinst, text width=1.45in, below=of b4] (b5) {Second report in an epoch (credential nullifier)};
  % column 3
  \node[dwarn, text width=1.45in, below=of h3] (c1) {\textbf{Correct totals}: all three aggregators honest and independent};
  \node[dwarn, text width=1.45in, below=of c1] (c2) {Physical output, emissions and deliveries (\emph{euthynai} spot audits)};
  \node[dwarn, text width=1.45in, below=of c2] (c3) {Labour rosters and carousel fraud};
  \node[dwarn, text width=1.45in, below=of c3] (c4) {Key and appliance custody; who issues credentials};
  \node[dwarn, text width=1.45in, below=of c4] (c5) {Consequences of faulting (attribution, de-accreditation)};
  \node[dbox, text width=4.72in, align=flush left, anchor=north] at ($(a5.south -| h2)+(0,-0.1)$) {\textbf{Not yet specified:} binding $B_u$ inside the report circuit; a post-quantum beacon and credential; secret ballots for consumer-federation votes; key recovery. Interim classical constructions cover the credential and beacon until post-quantum ones exist (Section \ref{sec:pq-urgency}).};
\end{tikzpicture}
\caption{Trust map: what cryptography guarantees, what is only detected, and what rests on institutions and physical audit.}
\label{fig:trust-map}
\end{figure}
\begin{description}[leftmargin=1.4em,itemsep=2pt]
  \item[Holds under stated assumptions.] Hybrid KEMs/signatures as in FIPS 203/204/205; BLAKE3 collision/preimage as in the BLAKE3 spec (128-bit target); Two-Root inclusion if a verifier has the leaf archive; \texttt{Prio3} measurement privacy if $\ge 1$ aggregator is honest.
  \item[Detected, not prevented.] Leaf omission given receipts and a challenge window; BFT equivocation given $\ge 1$ honest validator recording both votes; aggregate overspend if the public $p\cdot\sum y\le\sum B_u$ check is implemented.
  \item[Depends on social institutions.] Physical output, emissions, labour rosters, carousel fraud, aggregator independence, sortition appliance custody, and whether ``faulting'' an aggregator has consequences. There is no stake to slash; \textbf{faulting} means public attribution, dual-federation de-accreditation, and exclusion from the next sortition pool.
  \item[Aggregator compromise.] Replace the aggregator and either request fresh reports for the affected milestone or carry the previous milestone's aggregate. Stored reports cannot be re-verified with a replacement aggregator (Section~\ref{sec:carry-forward}). Old committee keys are a long-range risk: checkpoints must be weakly subjective (users pin a recent checkpoint) or re-signed by the new committee.
  \item[Key compromise.] Individual keys: revoke by council $k$-of-$n$ decision and re-issue. No specified recovery UX or secure-element/mobile-storage profile; community-hosted appliances may make the \emph{host} the real validator unless keys are in a citizen-held token. Until Android exposes post-quantum keys in hardware, keys use \textbf{split hybrid custody}: the classical half of each hybrid key stays in the platform keystore (Android Keystore or StrongBox, Apple Secure Enclave) and the ML-DSA half in software, encrypted under a hardware-held key. Until a CRQC exists a forger needs both halves, so the hybrid inherits the hardware's protection. Apple's Secure Enclave can hold ML-DSA keys directly from OS~26.\footnote{Apple, ``Enhancing your app's privacy and security with quantum-secure workflows,'' CryptoKit documentation, \url{https://developer.apple.com/documentation/cryptokit/enhancing-your-app-s-privacy-and-security-with-quantum-secure-workflows} (checked November 2026). No equivalent Android hardware profile was found.} Forward-secure erasure of committee keys is an \textbf{operational promise} until a ratchet (e.g.\ per-round key derivation with deletion) is specified.
  \item[Replay and domain separation.] Every signed object includes \texttt{epoch\_id}, \texttt{round}, \texttt{role}, and a domain-separation tag of the form \texttt{DPEP-WD-v0-<OBJECT>} (Listing~\ref{lst:cbor-idl}); the version component changes when the specification leaves working-draft status. Unspecified objects are not to be accepted.
  \item[Secret ballot.] Multi-signature bitmaps reveal who voted yes. Binding council votes may stay attributed; consumer-federation ballots that must be secret need a separate coercion-resistant scheme (open).
\end{description}

\section{Constitutional Governance and Sortition}
\label{sec:sortition}

To prevent software maintainers or server-rich councils from becoming a new coordinator class:
\begin{enumerate}[leftmargin=1.8em,itemsep=2pt]
  \item \textbf{Dual-federation ratification of WASM rules, catalogues, and caps:} Changes to \texttt{wasm\_rule\_hash}, product categories $\Omega$, or 5-plan voting grids require a $\ge 2/3$ dual supermajority of the National Consumer Federation and Congress of Industry Federations with a 60-day reproducible-build audit window, and cannot be altered mid-plan.
  \item \textbf{Sortition beacon (unresolved).} A beacon ``anchored'' on outgoing \texttt{SLH-DSA-192s} signatures is grindable if those signatures are randomised (FIPS 205 permits this). Commit-reveal has a last-revealer problem. A verifiable delay function could address grinding but is \textbf{not} specified here. A post-quantum beacon remains an open cryptographic problem, and there is no VRF in this draft. The interim construction is a threshold beacon run by the federations (Section~\ref{sec:pq-urgency}); its key generation and custody are unspecified.
  \item \textbf{Custody:} Sortition draws human delegates and provisions \texttt{Role-BFT} appliances. If keys live on the appliance, the host is the validator. Required: citizen-held tokens, or mark custody unresolved.
  \item \textbf{Committee size $n=128$:} $f=42$, quorum $\lceil 2n/3\rceil=86$. Exact binomial tails\footnote{Computed by \path{research/analysis/verify_r2.py} in the project repository.} for a single annual draw, corrupt share $p$: at $p=15\%$, $1.3\times 10^{-5}\%$ ($n=128$) vs.\ $9.5\times 10^{-3}\%$ ($n=64$); at $p=20\%$, $0.022\%$; at $p=25\%$, $1.82\%$; at $p=30\%$, $21\%$. Larger committees give little protection at a 30\% corrupt pool. An annual draw means a captured committee lasts a year; 1--4 rotations/year reduce but do not eliminate that window.
\end{enumerate}

\paragraph{Light clients.}
At $n=128$ a checkpoint carries 86 hybrid signatures, about 290--296~kB, which would dominate a phone's per-round traffic. Table~\ref{tab:light-clients} sets out what each lighter option trusts.

\begin{table}[htbp]
\centering
\caption{Light-client verification options for \texttt{Role-LP} (per checkpoint)}
\label{tab:light-clients}
\scriptsize
\begin{tabularx}{\textwidth}{L{1.35in} L{0.55in} Y}
\toprule
\textbf{Option} & \textbf{Download} & \textbf{What the phone trusts} \\
\midrule
Fetch and verify all 86 signatures & $\approx 296$ kB & Only the BFT assumption ($\le 42$ faulty validators). \\
Council coordinator verifies and attests & $\approx 3.4$ kB & The council's own \texttt{Role-CC}; a corrupt coordinator can feed a false checkpoint until a member checks. \\
Random sample of 8 signatures & $\approx 27$ kB & Probabilistic: if only the 42 corrupt validators really signed a checkpoint claiming 86, all 8 samples pass with probability $(42/86)^8 \approx 0.3\%$. Not a proof, and it fails if the sample is predictable to the adversary. \\
Full check at milestone rounds only & 296 kB $\times 3$/yr & BFT at milestones; coordinator attestation in between. \\
Rotating sync committee: 24 validators drawn by the beacon, quorum 16 & $\approx 55$ kB & That fewer than 16 of the 24 are faulty. With 42 faulty validators among 128, a draw yields a capturing committee with probability $0.017\%$ (16 members, quorum 11: $\approx 38$ kB, $0.19\%$). A captured committee lasts until the next rotation, and the draw is only as good as the beacon. \\
\bottomrule
\end{tabularx}
\end{table}

\noindent\textbf{Proposed default (for author ratification):} coordinator attestation every round, plus a full quorum check at each milestone round and whenever a member's client is on Wi-Fi, with every client able to opt into full checks. The sync-committee row, the pattern used by Ethereum light clients,\footnote{Ethereum consensus specifications, Altair light-client sync protocol, \url{https://github.com/ethereum/consensus-specs/blob/dev/specs/altair/light-client/sync-protocol.md}. Probabilities are exact hypergeometric tails from \path{research/analysis/verify_sync_committee.py}.} replaces trust in a member's own coordinator with trust in a rotating, randomly drawn committee, and is the preferred upgrade once the interim beacon (Section~\ref{sec:pq-urgency}) is running. The resulting mobile target, written into Kill Gate~3 (Table~\ref{tab:gates}), is $\le 150$~kB per ordinary round (55~kB report at milestones, price deltas, attestation, 512~B inclusion proof) and $\le 450$~kB in a round with a full quorum check.

% ==============================================================================
% CHAPTER 8: AI, MATHEMATICAL OPTIMISATION, AND ECONOMIC FAILURE MODES
% ==============================================================================
\chapter{Accountable AI, Mathematical Optimisation, and Failure Modes}

\chapterepigraph{Finally, we point out to those who fear that computers could become the new dictators, the computer programmes are human products, socially evaluated and improved each year\dots\ In short, neither consumers nor producers need accept computer facilitated alternatives.}{Michael Albert and Robin Hahnel, \emph{Looking Forward} (1991), Ch.~9, ``Revising Data in Subsequent Iterations''\protect\footnotemark}
\footnotetext{Albert and Hahnel, \emph{Looking Forward}, Ch.~9 (\texttt{lf9\_2.htm}).}

\section{Architectural Rule: Deterministic Core, Advisory Edge}

AI must never become an unaccountable central planner or surrogate coordinator class. DPEP enforces a strict boundary:
\begin{itemize}[leftmargin=1.6em,itemsep=2pt]
  \item \textbf{Core Clearing \& Pricing Engine (non-AI by specification):} All price updates, $SB/SC$ ratios, and Society Account settlements are specified as deterministic fixed-point arithmetic. Local IEEE-754 solver output, including any KKT certificate, checks stationarity of \emph{the programme the council chose to solve}. It does not prove a truthful objective, and MILPs have no KKT certificate of global optimality (Section~\ref{sec:fixed-point-bounds}).\footnote{Stephen Boyd, Neal Parikh, Eric Chu, Borja Peleato, and Jonathan Eckstein, ``Distributed Optimization and Statistical Learning via the Alternating Direction Method of Multipliers,'' \emph{Foundations and Trends in Machine Learning} 3, no. 1 (2011): 1--122.}
  \item \textbf{Edge Advisory \& Accessibility Support:} Classical hierarchical time-series models (ETS/ARIMA) forecast SKU size/color distributions; local QP solvers draft initial proposals from last year's bundle (eliminating David Schweickart's ``consumption list burden'' objection);\footnote{David Schweickart, \emph{Against Capitalism} (Cambridge: Cambridge University Press, 1993), pp.~329--334.} and optional local open-weight LLMs are restricted to translation, voice/screen-reader accessibility, and summarising qualitative \texttt{Lists} with mandatory links to raw council text.
\end{itemize}

\begin{table}[htbp]
\centering
\caption{Evaluation of key decision-support capabilities against verifiability and override criteria}
\label{tab:ai-eval}
\scriptsize
\setlength{\tabcolsep}{3.2pt}
\begin{tabularx}{\textwidth}{L{0.92in} L{1.02in} C{0.45in} L{1.05in} Y}
\toprule
\textbf{Capability} & \textbf{Best Method (Is AI Needed?)} & \textbf{Local?} & \textbf{Verifiability \& Challengeability} & \textbf{Override \& Power Concentration Safeguards} \\
\midrule
\textbf{1. SKU Demand Forecasting} & \textbf{No LLM needed:} Hierarchical ETS/ARIMA + demographics. & Yes ($<50\text{ ms}$) & Published MAPE error bounds \& open parameters. & Non-binding baseline; councils edit freely. \\
\addlinespace[3pt]
\textbf{2. Workplace Input Optimisation} & \textbf{No AI needed:} Conic QP / LP (\texttt{Clarabel.rs} / \texttt{HiGHS}). & Yes (WASM time unverified) & KKT certificates, where they exist, witness stationarity of the local convex programme, not a truthful $W_j$ or MILP optimum. & Presents 3 Pareto-frontier options for worker council vote. \\
\addlinespace[3pt]
\textbf{3. Household Proposal Drafting} & \textbf{No AI needed:} Deterministic diff + constrained utility QP. & Yes (on device) & Line-by-line diff against prior year's bundle. & Mandatory review UI highlights top changes; zero data leaves user's device. \\
\addlinespace[3pt]
\textbf{4. Tradeoff Explanation} & \textbf{Deterministic Leontief Cost-Push Attribution} ($(I-A)^{-1}$). & Yes ($<50\text{ ms}$) & Exact algebraic decomposition of $\Delta P_i$. & Avoids LLM hallucination by rendering exact resource/labour/pollution cost shares. \\
\addlinespace[3pt]
\textbf{5. Qualitative \texttt{List} Translation} & \textbf{Yes:} Local Neural Machine Translation + TTS. & Yes (\texttt{Role-CC}) & Original source text \& BLAKE3 CID always linked alongside. & Bilingual delegates can post verified overrides; expands cross-regional inclusion. \\
\addlinespace[3pt]
\textbf{6. Late-Round 5-Plan Synthesis} & \textbf{No AI needed:} Parametric LP over 5 fixed weight vectors. & Yes ($<2\text{ s}$) & Deterministic WASM LP verifiable by any council. & Councils can vote to reject all 5 plans and continue iterating (LF Ch.~9 \texttt{lf9\_2.htm}). \\
\bottomrule
\end{tabularx}
\end{table}

\section{Substantive Engagement with Major Economic \& Democratic Critics}
\label{sec:critics}

This section engages the main theoretical critiques of Participatory Economics across philosophy, economics and mechanism design: the Winter 2024 special issue of the \emph{Erasmus Journal for Philosophy and Economics} (John O'Neill, Hannes Kuch, Jan Philipp Dapprich), read in the open-access issue PDF;\footnote{\emph{Erasmus Journal for Philosophy and Economics} 17, no.~2 (Winter 2024), issue PDF at \url{https://ejpe.org/journal/issue/download/35/67}: John O'Neill, ``Pluralism, Ecology and Planning,'' pp.~27--48, \url{https://doi.org/10.23941/ejpe.v17i2.884}; Hannes Kuch, ``Centralization Tendencies in Participatory Planning,'' pp.~49--72, \url{https://doi.org/10.23941/ejpe.v17i2.885}; Jan Philipp Dapprich, ``Democratic Theory and Economic Planning: A Critique of Participatory Economics,'' pp.~73--93, \url{https://doi.org/10.23941/ejpe.v17i2.867}; John O'Neill, ``Incommensurability, Environment and Planning: A Response to Hahnel's Reply,'' pp.~117--124, \url{https://doi.org/10.23941/ejpe.v17i2.933}.} David Laibman (2014), consulted through the publisher's abstract, with Robin Hahnel's (2014) response;\footnote{David Laibman, ``Horizontalism and Idealism in Socialist Imagination: An Appraisal of the Participatory Economy,'' \emph{Science \& Society} 78, no.~2 (2014): 207--234, \url{https://doi.org/10.1521/siso.2014.78.2.207}; Robin Hahnel, ``Response to David Laibman's `Appraisal of the Participatory Economy','' \emph{Science \& Society} 78, no.~3 (2014): 379--389, \url{https://doi.org/10.1521/siso.2014.78.3.379}. Only the abstract of Laibman's article and the bibliographic record of the response were consulted.} Paul Cockshott and Allin Cottrell (1993);\footnote{W.~Paul Cockshott and Allin Cottrell, \emph{Towards a New Socialism} (Nottingham: Spokesman Books, 1993).} J\'anos Kornai (1980, 1986); Martin Weitzman (1974, 1980); and David Schweickart (1993).

\begin{enumerate}[leftmargin=1.8em,itemsep=3pt]
  \item \textbf{Dapprich (2024) on democratic theory, liberty, and general procedures:} Dapprich asks whether three justifications of democracy (protecting interests, revealing preferences, and liberty in the neo-Roman republican sense) support citizen participation in concrete plans, and argues that they do not: citizens should instead choose \emph{general planning procedures}, an ``algorithmic planning process'' in his words, which protects them from the arbitrary power of co-workers and neighbours as well as of central planners. He also argues that most participatory input in ParEcon actually comes from individual consumption requests, leaving consumer councils with little to vote on, and that without constitutional constraints a worker is exposed to the arbitrary will of the majority on their own council.
  \begin{itemize}[leftmargin=1.4em,itemsep=1pt]
    \item \emph{Where Dapprich is right:} Neighbour veto of in-budget personal consumption is a liberty problem; DPEP's non-veto rule follows PEPE, and cryptographic invisibility goes further (Section~\ref{sec:consumption-privacy}). His point about workplace majorities is also well taken: DPEP needs entrenched individual rights inside councils (safety, exit to another council through the EFB, BJC protections), not only majority procedures.
    \item \emph{Where this monograph disagrees:} Workplace pace, safety and emissions, and city-scale public goods, are collective activities whose terms cannot be fixed in advance by a general procedure without someone writing the procedure's parameters, which re-creates a planning authority one level up. DPEP therefore keeps councils as decision-makers for collective activities and uses general procedures (WASM rules ratified by dual supermajority, Section~\ref{sec:sortition}) as constraints on them. This answers Dapprich only in part.
  \end{itemize}
  \item \textbf{Kuch (2024) on centralisation tendencies:} Kuch argues that a participatory economy needs more central institutions than Hahnel suggests, for strong normative reasons: industry federations act as gatekeepers, appellate bodies and adjustment agencies; national federations decide investment and society-wide public goods; disciplinary bodies may be needed to counter the procedure's false incentives; and conflicts between federations may require a national legislature or something like a ministry of planning. He doubts that council-democratic remedies resolve this. DPEP's responses (Algorithm 3c's federation MILP as an aid rather than an authority, sortition, the Innovation Jury below, dual-federation ratification) are institutional proposals that Kuch's argument puts directly in question; they are tested in P7 and are not shown to work.
  \item \textbf{Cockshott \& Cottrell (1993) on direct LP and scalar labour-time:} Cockshott and Cottrell propose replacing iterative council negotiation with central LP in embodied labour hours. Two objections are raised here (from the bibliographic record and the secondary literature, not a full reading of \emph{Towards a New Socialism}): (a) scalar labour-time omits the independent scarcity of \textbf{non-producible natural resources ($r^*$), ecological sinks ($E_z^{\max}$) and fixed capital ($k^*$)}, and differential workplace disutility; (b) a central LP that needs each workplace's full private $\mathcal{T}_j$ collides with Hayek (1945) on tacit knowledge \emph{and} with worker self-management. Hayek's point also applies to DPEP: coarse categories, an annual cycle and published technical recipes (REQ-DATA-03) still lose local know-how, and Algorithm 3c's federation MILP is a limited central solve over \emph{posted menus}, not over hidden $\mathcal{T}_j$. Treating 3c as a Phase-1 oracle and a ratified aid, not an authority, is the intended boundary; it remains a tension, not a resolution.
  \item \textbf{Laibman (2014), Kornai (1980) and Weitzman (1974, 1980) on equilibrium, soft budgets and the ratchet effect:}
  \begin{itemize}[leftmargin=1.4em,itemsep=1pt]
    \item \emph{The Walrasian resemblance and dynamic entry (Laibman 2014; Kuch 2024):} Laibman's abstract argues that the Albert--Hahnel model is designed in an ``abstractly speculative'' way, that its fear of authority and hierarchy has roots in classical anarchism, and that the result bears an ``uncanny resemblance to Walrasian competitive market equilibrium.'' Hahnel replied in the following issue. This monograph does not attribute to Laibman any argument beyond his abstract. One practical consequence of the Walrasian reading is that an iterative procedure over a fixed list of producers says nothing about how new producers enter. DPEP's proposal: an accredited worker group refused accreditation by an Industry Federation (Kuch's gatekeeping concern) can appeal to a \textbf{Cross-Industry Innovation Jury}, selected by sortition and co-ratified by the Consumer Federation, backed by a dedicated R\&D capital slice. This is untested (P7).
    \item \emph{Hardening Soft Budget Constraints (Kornai 1980, 1986):} If councils with $SB_j/SC_j < 1$ could obtain unlimited \texttt{requests\_appeal} waivers, budget discipline would collapse.\footnote{J\'anos Kornai, \emph{Economics of Shortage} (Amsterdam: North-Holland, 1980); Kornai, ``The Soft Budget Constraint,'' \emph{Kyklos} 39, no. 1 (1986): 3--30; Martin L.\ Weitzman, ``Prices vs.\ Quantities,'' \emph{Review of Economic Studies} 41, no. 4 (1974): 477--491; Weitzman, ``The `Ratchet Principle' and Performance Incentives,'' \emph{Bell Journal of Economics} 11, no. 1 (1980): 302--308.} DPEP limits appeal waivers to one planning cycle (requiring a physical retooling audit) and funds them from a capped \textbf{Industry Transition Pool}, so granting a waiver directly reduces peer councils' discretionary industry reserve.
    \item \emph{Ratchet effect (Weitzman 1980) --- open:} Basing endowment charges on physical assets rather than last year's output removes one Soviet-style ratchet. It does not remove the fear that revealed capacity will feed next year's benchmarks, public blueprints (REQ-DATA-03), or $SB/SC$ comparisons. Weitzman (1980) is cited at bibliographic level. Status: open.
  \end{itemize}
  \item \textbf{Effort rating, surveillance, ecology and limits (O'Neill 2024):} O'Neill makes three criticisms. First, distributing income by a metric of effort would involve \emph{pervasive surveillance and potential humiliation}; this bears directly on the effort default discussed in Section~\ref{sec:effort-default}. Second, a single price metric of opportunity cost and cost--benefit analysis cannot address value pluralism and incommensurability; he defends rational choice without commensurability and in-kind analysis in the tradition of Neurath and Kapp. Third, environmental limits must be taken more seriously than Hahnel does: substitution is necessary but not sufficient, and he proposes consumption and production ``corridors'' above sufficiency for all but below excess. DPEP's intended responses: (a) \textbf{strict one-sided biophysical ceilings ($r^*, E_z^{\max}$)}, which are in-kind constraints rather than prices, set by long-run environmental planning, with an \textbf{intergenerational trusteeship} mandate that living cohorts cannot relax $E_z^{\max}$ without a supermajority plus a sortitioned youth/future-persons panel; (b) a \textbf{per-capita exposure floor} so that no neighbourhood's modelled dose exceeds a published health threshold even if the regional CAP accepts the aggregate $E_z$; (c) a \textbf{plume sub-zone veto} for the most-impacted $CC_h\subset CAP_z$; (d) earmarking $\ge 50\%$ of $p_z^{\text{pol}} E_z$ to remediation and preventative health. Rules (c) and (d) create hold-out and earmark-gaming problems (a sub-zone can veto to extract transfers; 50\% earmarks invite recategorising ordinary spending), and DPEP has no counterpart yet to an upper consumption corridor. Status: working rules, not solved mechanism design.
  \item \textbf{Governance of solver defaults and deliberative fatigue (Devine 1988; Albert 2003, Ch.~18):} Even without AI, local QP solver defaults could become de facto decisions if citizens rubber-stamp them. DPEP bans silent auto-submission of $>5\%$ effort or category shifts, requires active choice among $\ge 3$ options at milestone rounds, and measures rubber-stamping, comprehension and the distribution of influence in Phase~1H (Table~\ref{tab:gates}).
  \item \textbf{Black-market arbitrage and wage labour (Albert 2003, Ch.~20):} \texttt{Prio3} is intended to keep personal bundles private and to reject reports that violate the encoded budget predicate \emph{if} that circuit is specified and sound. An unauthorised firm still cannot draw intermediate inputs through the planning tree. Off-protocol barter remains a social question. No ``mathematical prevention'' of wage labour is claimed.
\end{enumerate}

% ==============================================================================
% CHAPTER 9: REORGANISED THREE-PHASE EXPERIMENTAL PROGRAMME AND BENCHMARKS
% ==============================================================================
\chapter{Three-Phase Experimental Programme, Kill Gates, and Benchmarks}

\chapterepigraph{The easiest `experiments' to implement are a variety of computer simulations in which various features of the system and assumptions about actors' behaviours can be altered to see how outcomes change\dots\ Conceivably one could run a `shadow' participatory economy within an existing economy.}{Michael Albert and Robin Hahnel, \emph{The Political Economy of Participatory Economics} (1991), Ch.~6, ``Feasibility,'' pp.~108, 129\protect\footnotemark}
\footnotetext{Albert and Hahnel, \emph{Political Economy of Participatory Economics}, Ch.~6, pp.~108, 129.}

\section{Why DPEP Uses Gated Phases Instead of a Single Monolithic Build}

Attempting to build non-convex solvers, QUIC, aggregation, WASM, PQC, \texttt{Prio3}, and mobile UIs in one monolithic prototype would conflate mechanism validation with distributed-systems engineering. Each phase therefore ends in a \textbf{kill gate}: a test with a pre-registered pass condition, a fail condition, and a stated consequence. A gate that cannot fail is not a gate.

\begin{table}[htbp]
\centering
\caption{Staged experimental programme (budgets as in Table~\ref{tab:exec-budget})}
\label{tab:poc-tests}
\scriptsize
\setlength{\tabcolsep}{3.5pt}
\begin{tabularx}{\textwidth}{L{1.20in} Y L{1.05in}}
\toprule
\textbf{Phase \& Stack} & \textbf{Core experiments} & \textbf{Duration \& staff} \\
\midrule
\textbf{Phase 1: Economic coordination}\newline {\tiny (Python, \texttt{NumPy}, \texttt{HiGHS}, \texttt{Clarabel})} & Pre-registered instance set (below); Algorithms 3a--3d against a central LP/MILP oracle; scaling from 50 to 1,000 councils. & 4--6 months; 1.0 FTE OR/economics \\
\addlinespace[3pt]
\textbf{Phase 1H: Low-tech human probe}\newline {\tiny (disclosed trusted server + web UI)} & 30--60 volunteers in 3 worker and 3 consumer councils run three milestone rounds of a simplified plan. Ethics review and consent first. & 3--4 months; 1.0 FTE UX/dev; in parallel with Phase~1 \\
\addlinespace[3pt]
\textbf{Phase 2: Distributed coordination}\newline {\tiny (Rust \texttt{dpep-node}, CBOR, QUIC, \texttt{Wasmtime})} & WAN testbed at $\ge 10^4$ simulated participants; fault injection (censorship, equivocation, 20\% partition); \texttt{Prio3} verification and signature checks in the round. & 4--6 months; 2.5 FTE \\
\addlinespace[3pt]
\textbf{Phase 3: PQC, \texttt{Prio3} \& pilot}\newline {\tiny (FIPS 203--205, \texttt{Prio3}, mobile PWA)} & Device benchmarks; external audit; 10--30 cooperative field pilot, starting with the council-level trust model (Section~\ref{sec:privacy-baseline}). & 9--12 months; 4.0 FTE + audit \\
\bottomrule
\end{tabularx}
\end{table}

\subsection{Kill Gates as Tests}
\label{sec:gates}

Table~\ref{tab:gates} states each gate. \textbf{All numeric thresholds are proposed, for author ratification}, and each carries its reason. Consequences use four outcomes: \emph{proceed}; \emph{narrow scope} (continue only for the problem class or setting that passed); \emph{redesign and rerun} (once, on a freshly pre-registered set); \emph{stop} (publish the negative result and fund no later phase on that basis).

{\scriptsize\setlength{\tabcolsep}{3pt}\renewcommand{\arraystretch}{1.15}
\begin{longtable}{L{0.45in} L{1.50in} L{1.10in} L{1.50in}}
\caption{Kill gates: pass and fail conditions, consequences and reasons (thresholds proposed)}\label{tab:gates}\\
\toprule
\textbf{Gate} & \textbf{Pass (all required)} & \textbf{Fail (any)} & \textbf{Consequence and reason} \\
\midrule
\endfirsthead
\multicolumn{4}{l}{\scriptsize\itshape Continued from previous page}\\
\toprule
\textbf{Gate} & \textbf{Pass (all required)} & \textbf{Fail (any)} & \textbf{Consequence and reason} \\
\midrule
\endhead
\midrule\multicolumn{4}{r}{\scriptsize\itshape Continued on next page}\\
\endfoot
\bottomrule
\endlastfoot
\textbf{1}\newline Economic & For the best candidate: (a) $\ge 90\%$ of convex and Leontief instances \emph{cleared} (Definition~\ref{def:success}) within 25 rounds; (b) $\ge 70\%$ of staggered-lumpy instances cleared; (c) on Scarf and identical-lumpy instances, either cleared or correctly reported \emph{feasible, not cleared} or infeasible with diagnostics in 100\% of runs; (d) 100\% of infeasible instances reported infeasible; (e) $\ge 70\%$ of endogenous-income instances cleared by 3d; (f) median welfare gap to the oracle $\le 5\%$, 90th percentile $\le 15\%$. & Fewer than 50\% of convex and Leontief instances cleared; \emph{or} any infeasible instance reported feasible; \emph{or} median welfare gap $> 15\%$; \emph{or} endogenous-income clearance $< 40\%$. & \emph{Proceed} if all pass. \emph{Narrow scope} if (a), (d) and (f) pass but (b) or (e) does not: Phase~2 then covers only the passing classes, with lumpy decisions left to federation MILP and ratification. \emph{Redesign and rerun} once if between pass and fail. \emph{Stop} if the fail condition holds after the rerun. \emph{Reasons:} convex instances have known convergent methods, so failure there signals a design or implementation fault; 70\% for lumpy instances reflects the one-good test, in which 57\% of rounds clear; 25 rounds is about twice the 11.85-round \texttt{pequod-cljs} baseline; reporting an infeasible plan as feasible is never acceptable. \\
\addlinespace[3pt]
\textbf{1H}\newline Human probe & (a) $\ge 80\%$ of participants complete all three milestones; (b) median score $\ge 70\%$ on a pre-registered comprehension quiz (prices, budgets, what an override does); (c) $\le 25\%$ of solver drafts accepted without opening the change view; (d) of coded override reasons, $\ge 50\%$ are substantive (need, preference, local knowledge) rather than confusion; (e) the most active 20\% of participants originate $\le 60\%$ of adopted amendments. & Completion $< 60\%$; \emph{or} median comprehension $< 50\%$; \emph{or} $> 50\%$ of drafts accepted unopened; \emph{or} the most active 20\% originate $> 80\%$ of adopted amendments. & \emph{Proceed} if all pass. \emph{Redesign and rerun} (interface, facilitation, schedule) once if between. \emph{Stop} human-facing claims and keep simulator research only if the fail condition holds after the rerun. \emph{Reasons:} override rate alone rewards bad defaults, so it is paired with comprehension, reasons and the spread of influence; (e) is a direct test for an emerging technical or articulate elite. \emph{Limits:} 30--60 self-selected volunteers cannot show generalisation to a population, long-run fatigue, or power dynamics at scale; with $n=40$ a 95\% interval on a proportion near 0.5 is about $\pm 15$ points. \\
\addlinespace[3pt]
\textbf{2}\newline Distributed & (a) 100\% of injected omissions and tamperings detected within $T_{\text{challenge}}$ over $\ge 1{,}000$ fault-injection trials; (b) automated round time (excluding human deliberation), including \texttt{Prio3} verification and signature checks, p95 $\le 10$~min at the tested scale ($\ge 10^4$ participants), and $\le 60$~min when extrapolated to $10^6$ from measured per-item costs with $\le 64$ cores per aggregator; (c) total bytes and CPU $\le 10\times$ an Architecture~A baseline on the same workload. & Any undetected tampering; \emph{or} p95 $> 30$~min at the tested scale; \emph{or} overhead $> 50\times$ the A baseline. & \emph{Proceed} if all pass. \emph{Narrow scope} (smaller committees, coarse reports only, regional rather than national deployment) or \emph{redesign and rerun} if between. \emph{Stop} the distributed build and pilot on Architecture~A if the fail condition holds. \emph{Reasons:} detection is the point of Two-Root aggregation, so one miss fails it; 10~min fits the 1--10~min target of Table~\ref{tab:consistency}; a tenfold overhead is a reasonable price for removing a trusted operator, fiftyfold is not. \\
\addlinespace[3pt]
\textbf{3}\newline Crypto, mobile, pilot & (a) External audit with no unresolved critical or high findings before field use; (b) on a 3~GB-RAM Android device: $\le 150$~kB per ordinary round and $\le 450$~kB in a full-verification round (Section~\ref{sec:sortition}), peak RAM $\le 150$~MB, coarse report shard $\le 1$~s; (c) pilot: $\ge 70\%$ of member councils still participating at the end, no unresolved privacy incident, and members' rated legitimacy of the plan not below their rating of the pre-pilot process. & Unresolved critical audit finding; \emph{or} any device metric more than $2\times$ its cap; \emph{or} pilot retention $< 40\%$. & \emph{Proceed} to wider pilots if all pass. \emph{Narrow scope} to the council-level trust model (no household cryptography) if (b) fails but (c) passes. \emph{Stop} field use of custom circuits on an unresolved critical finding. \emph{Reasons:} the bandwidth caps follow from the coarse 55~kB report plus the light-client options; 150~MB leaves room on entry-level phones; retention and legitimacy test whether people want to keep using the process, which no simulation can show. \\
\end{longtable}}

\paragraph{Pre-registration.}
Before each gate, the fast-moving rows of the component maturity table (Table~\ref{tab:maturity}) are re-checked and the result is recorded with the gate's pre-registration.
Before any candidate solver is run on them, the instance generators, random seeds, instance counts, metrics, thresholds and analysis scripts are committed to a pre-registration directory in the project repository,\footnote{\path{research/preregistration/}.} tagged with a signed Git tag, and the tag's commit hash is deposited with an independent registry (for example the Open Science Framework). Results are reported against that tag. Deviations are logged with reasons in the same directory, and any rerun uses a freshly registered set. Phase~1H registers its quiz, coding scheme for override reasons, and influence measure in the same way, together with its ethics approval.

\paragraph{Pre-registered Phase~1 instance set.}
At least 20 instances in each family, at 50, 200 and 1,000 councils: (1) convex baselines (Cobb--Douglas and CES with fixed budgets); (2) Leontief input--output economies; (3) Scarf-style cycles; (4) identical lumpy plants; (5) staggered lumpy plants; (6) deliberately infeasible plans (caps below minimum requirements); (7) endogenous income through the Society Account; (8) binding ecological caps; (9) strategic misreporting, in which a fixed share of councils understate capacity by a fixed fraction, measuring the resulting price distortion and the misreporters' gain (reported as a diagnostic; no pass threshold is set until a baseline exists).

\paragraph{What a 1,000-council result can and cannot show.}
Phase~1 can show whether rounds-to-clear grow with the number of councils between 50 and 1,000 (a published log--log slope) and whether failure modes depend on scale. A flat or falling slope makes extrapolation plausible; a rising slope is evidence against national-scale feasibility. A pass at 1,000 councils does not show convergence at $10^6$ participants: heterogeneity, the number of distinct goods, strategic behaviour and communication costs all change with scale, and Phase~2 measures the last of these separately.

\subsection{Ethics, Consent and Data Protection}
\label{sec:ethics}
Phases~1H and~3 involve human participants and personal data. Before recruitment each requires: independent ethics review (a university research ethics committee if an academic partner is involved, otherwise an independent ethics panel whose membership is published); informed, written consent with a plain-language description of what is recorded, who can see it, and the right to withdraw and have data deleted; data minimisation and a retention schedule; a data-protection impact assessment where GDPR or UK GDPR applies (Article~35); and fair compensation for participants' time. Phase~3 adds a data-processing agreement with each pilot cooperative and a named contact for complaints. These add 4--10 weeks of elapsed time before each phase and the non-labour costs shown in Table~\ref{tab:exec-budget}.

\subsection{The Effort-Rating Default (proposed)}
\label{sec:effort-default}
Equation~(\ref{eq:effort-cap}) has two parameters with different meanings. $\bar\epsilon_{\max}$ caps a council's \emph{mean} effort multiplier and so prevents council-wide grade inflation; the individual spread $\epsilon_{\max}$, with $|\epsilon_u - 1| \le \epsilon_{\max}$, decides whether two people working the same hours in the same council can be paid differently for intensity. Maxim~3 asks for that difference. O'Neill (2024) argues that rating it means pervasive surveillance and potential humiliation, and cooperatives often prefer flat hourly pay. This is a political choice, not a technical one. \textbf{Proposed default (for author ratification):} $\bar\epsilon_{\max} = 1.0$ in every council; individual spread $\epsilon_{\max} = 0$ (flat pay per accredited hour, with onerousness handled through job-complex balancing and hours) unless a council votes to allow peer-rated spread, capped at $\epsilon_{\max} = 0.15$ and reviewed after one plan year. Phase~1 simulates both settings; P6 tests them with cooperatives. Under this default DPEP honours Maxim~3 through hours and job balancing but not, by default, through intensity, and the monograph says so.

\section{Multi-Scale Analytical Projections ($10^2$ to $10^6$ Participants)}

{\small\color{SlateDark} \textbf{Methodological note:} Only the \texttt{pequod-cljs} CSV convergence counts, the Python stress tests and the \texttt{Prio3} benchmark have been executed.\footnote{Scripts and benchmark code are in \path{research/analysis/} in the project repository.} Table~\ref{tab:scale-bench} uses a 10-round annual plan with three fresh household-report rounds; Kill Gate~1's 25-round bound is a halt criterion, not the table's round count. Archive rows use the coarse single-key proposal sizes of Table~\ref{tab:proposal-sizes}; multi-signature councils multiply them by up to 3.3. All rows assume an $n=128$ committee; small federations could use smaller committees.}

\begin{table}[htbp]
\centering
\caption{Analytical scaling projections across five population scales (10-round plan; SI units)}
\label{tab:scale-bench}
\scriptsize
\setlength{\tabcolsep}{3.0pt}
\begin{tabularx}{\textwidth}{Y R{0.52in} R{0.55in} R{0.58in} R{0.62in} R{0.72in}}
\toprule
\textbf{Metric (Category B analytical estimate)} & \textbf{100} & \textbf{1,000} & \textbf{10,000} & \textbf{100,000} & \textbf{1,000,000} \\
\midrule
Total councils \& federations ($M$) & 24 & 51 & 430 & 4,120 & 40,600 \\
Priced categories ($G = |\Omega|$) & 30 & 80 & 500 & 2,000 & 10,000 \\
Council proposals / round & 20 & 40 & 400 & 4,000 & 40,000 \\
Raw leaf archive / round & 119 kB & 238 kB & 2.4 MB & 24 MB & 237 MB \\
Raw leaf archive / year & 1.2 MB & 2.4 MB & 24 MB & 237 MB & 2.4 GB \\
Mobile traffic / round (\texttt{Role-LP}), target & \multicolumn{5}{r}{$\le 150$ kB ordinary round; $\le 450$ kB with full quorum check (Kill Gate 3)} \\
Layer-1 checkpoints / year (425 blocks, $n=128$) & 124 MB & 125 MB & 132 MB & 158 MB & 298 MB \\
Household-report CPU / aggregator / fresh round (desktop, coarse) & $<0.1$ s & 0.25 s & 2.5 s & 25 s & 4 core-min \\
Mobile peak RAM target (\texttt{Role-LP}) & \multicolumn{5}{r}{$\le 150$ MB (unmeasured)} \\
Automated round time & \multicolumn{5}{r}{Not estimated; measured in Phase~2 against Kill Gate~2} \\
Server hosting (excl.\ operations and audits) & \$0--\$5/mo & \$10--\$25/mo & \$80--\$150/mo & \$400--\$800/mo & \$3,000--\$6,000/mo \\
\bottomrule
\end{tabularx}
\end{table}

% ==============================================================================
% CHAPTER 10: OPEN THEORETICAL, INSTITUTIONAL, AND TECHNICAL QUESTIONS
% ==============================================================================
\chapter{Open Theoretical, Institutional, Economic, and Technical Questions}

\chapterepigraph{A simulation incorporates more detail than a differential economic model, but still misses the richness of living interaction. As a result, the full characteristics of participatory economics and the diverse solutions to institutional problems associated with it will ultimately only reveal themselves in actual operations.}{Michael Albert and Robin Hahnel, \emph{The Political Economy of Participatory Economics} (1991), Ch.~6, p.~129\protect\footnotemark}
\footnotetext{Albert and Hahnel, \emph{Political Economy of Participatory Economics}, Ch.~6, p.~129.}

We conclude our study by ranking the 15 most important unresolved questions across theoretical economics, institutional sociology, and distributed systems engineering, specifying how each will be investigated through either our simulation programme (\texttt{[SIM]}) or human-in-the-loop cooperative pilots (\texttt{[HUMAN]}):

{\scriptsize\setlength{\tabcolsep}{3.5pt}
\begin{longtable}{C{0.28in} L{1.50in} L{0.72in} L{1.90in}}
\caption{Ranked Research Agenda of Unresolved Questions (Priorities 1--15)}
\label{tab:open-questions} \\
\toprule
\textbf{Rank} & \textbf{Unresolved Question} & \textbf{Domain} & \textbf{Experimental Investigation Protocol} \\
\midrule
\endfirsthead
\multicolumn{4}{l}{\scriptsize\itshape Continued from previous page} \\
\toprule
\textbf{Rank} & \textbf{Unresolved Question} & \textbf{Domain} & \textbf{Experimental Investigation Protocol} \\
\midrule
\endhead
\midrule
\multicolumn{4}{r}{\scriptsize\itshape Continued on next page} \\
\endfoot
\bottomrule
\endlastfoot
\textbf{P1} & \textbf{Convergence under non-convexities, setup costs \& Leontief complementarities} & \texttt{[ECON / OR]} (\texttt{[SIM]}) & Replace unconstrained Cobb-Douglas in a new simulator (not by patching \texttt{pequod-cljs} FOCs); benchmark Algorithms 3a--3c on pre-registered instances. \\
\addlinespace[3pt]
\textbf{P2} & \textbf{Stability of coupled endogenous income} & \texttt{[ECON]} (\texttt{[SIM]}) & Couple Society Account updates into active rounds with Algorithm~3d (damped budget fixed point; damping $\alpha\in(0,1]$ is a relaxation, not a second num\'eraire, which remains $w_{\text{base}}\equiv 1$). Report cycling, multiple fixed points and the effect of the one-round budget lag. \\
\addlinespace[3pt]
\textbf{P3} & \textbf{Deliberative Complexity \& Participation Fatigue in Multi-Round Planning} & \texttt{[INST / HCI]} (\texttt{[HUMAN]}) & Test the ``3-Milestone Hybrid Schedule'' (human rounds at $t=1,4,8$ + automated intermediate solver rounds); measure NASA-TLX load. \\
\addlinespace[3pt]
\textbf{P4} & \textbf{Biophysical Incommensurability vs.\ PDRM Pollution Pricing (O'Neill 2024)} & \texttt{[ECOLOGY]} (\texttt{[SIM/GOV]}) & Verify hard barrier convergence when aggregate demand hits absolute biophysical caps ($r^*, E_z^{\max}$) above local CAP PDRM pricing. \\
\addlinespace[3pt]
\textbf{P5} & \textbf{Coarse-to-Granular SKU Disaggregation \& Slack Buffer Sizing} & \texttt{[ECON / OR]} (\texttt{[SIM]}) & Simulate 200 coarse reporting categories (the proposed privacy baseline, Section~\ref{sec:privacy-baseline}) mapped to 3,000 planning categories and 20,000 SKUs under stochastic preference drift, against slack $\eta \in [1\%, 10\%]$; measure the unmet-demand share it causes. \\
\addlinespace[3pt]
\textbf{P6} & \textbf{Workplace Effort-Rating Dynamics: Solidarity vs.\ Surveillance or Parity} & \texttt{[INST]} (\texttt{[HUMAN]}) & Pilot with worker cooperatives comparing scaled effort multipliers vs.\ flat hourly parity ($\epsilon_u \equiv 1.0$) under BJC rotation. \\
\addlinespace[3pt]
\textbf{P7} & \textbf{Preventing Industry Federation Cartelization \& Gatekeeping (Kuch 2024)} & \texttt{[INST / ECON]} (\texttt{[HUMAN]}) & Test ``Dual-Path Startup Accreditation'' allowing rejected innovators to appeal to a sortition Innovation Jury + Consumer Federation. \\
\addlinespace[3pt]
\textbf{P8} & \textbf{Calibrating Balanced Job Complexes Across Disparate Industries} & \texttt{[INST / OR]} (\texttt{[SIM/HUM]}) & Model cross-workplace BJC rotation as a bipartite matching problem in EFB accounting for travel \& skill constraints. \\
\addlinespace[3pt]
\textbf{P9} & \textbf{Inter-Regional Investment Apportionment Across Unequal Endowments} & \texttt{[ECON]} (\texttt{[SIM]}) & Compare per-capita vs.\ Porto Alegre need-weighted regional investment matrices over 10-year horizons. \\
\addlinespace[3pt]
\textbf{P10} & \textbf{Reproductive \& In-Home Caring Labour Integration (DEP Ch.~10)} & \texttt{[FEMINIST ECON]} (\texttt{[HUMAN]}) & Convene Women's and Caregivers' Caucuses to define care-allowance rules and household BJC norms. \\
\addlinespace[3pt]
\textbf{P11} & \textbf{Byzantine Resilience Under Prolonged Regional Network Partitions} & \texttt{[DIST-SYS]} (\texttt{[SIM]}) & Benchmark Regional Proxy Extrapolation accuracy when 10\% of councils are partitioned for multiple rounds. \\
\addlinespace[3pt]
\textbf{P12} & \textbf{Standardisation of Compact Post-Quantum Threshold Signatures} & \texttt{[CRYPTO]} (\texttt{[SW-ARCH]}) & Councils use $k$-of-$n$ multi-signatures (3,437~B per signer) until a standard post-quantum threshold signature exists; \texttt{CryptoSuiteId} agility allows the switch. NIST opened a call for multi-party threshold schemes in January 2026 (NIST IR 8214C), a reference process rather than a standards competition. Proposal sizes depend on $k$ (Table~\ref{tab:proposal-sizes}). \\
\addlinespace[3pt]
\textbf{P13} & \textbf{Privacy anonymity sets in small hamlets ($m < 30$)} & \texttt{[PRIVACY]} (\texttt{[SW-ARCH]}) & Apply $m_{\min}=30$ per coarse reporting category and aggregate rare categories at Ward level with $m_{\text{rare}}=250$; measure differencing leakage between milestone releases under the carry-forward rule of Section~\ref{sec:carry-forward}. \\
\addlinespace[3pt]
\textbf{P14} & \textbf{Endpoint Security vs.\ Open Modifiability on User Smartphones} & \texttt{[SECURITY]} (\texttt{[SW-ARCH]}) & Use local hardware keystore unlock + social anomaly verification without mandatory proprietary OS attestation lock-in; split hybrid custody until Android holds post-quantum keys in hardware (Section~\ref{sec:threat-model}). \\*
\addlinespace[3pt]
\textbf{P15} & \textbf{Bootstrapping a ``Parallel Economy'' Inside Surrounding Markets} & \texttt{[ADOPTION]} (\texttt{[HUMAN]}) & Implement a Dual-Price Fiat Reserve Bridge in Sandstr\"om's Society Account for federated cooperatives (PEPE Ch.~6 p.~129). \\
\end{longtable}}

% ==============================================================================
% CHAPTER 11: PRACTICAL ADOPTION STRATEGY AND THEORY OF TRANSITION
% ==============================================================================
\chapter{Practical Adoption Strategy, Cooperative Pilots, and Theory of Transition}

\chapterepigraph{Conceivably one could run a `shadow' participatory economy within an existing economy. This could occur in a subset of the whole economy, whether geographic or by industrial sector\dots\ not only [to] serve as a laboratory to test participatory planning, but also as a school, both educating the public and developing the means to implement participatory planning without undue disruption.}{Michael Albert and Robin Hahnel, \emph{The Political Economy of Participatory Economics} (1991), Ch.~6, ``Experiment 2: Developing a Parallel Economy,'' pp.~129--130\protect\footnotemark}
\footnotetext{Albert and Hahnel, \emph{Political Economy of Participatory Economics}, Ch.~6, pp.~129--130.}

\section{Distinguishing Research Simulation from Practical Cooperative Coordination}

To avoid equating small-scale cooperative software with a full post-capitalist economy, DPEP separates two distinct applications:
\begin{enumerate}[leftmargin=1.8em,itemsep=2pt]
  \item \textbf{The Scientific Research Platform (Phases 1--2):} A computational laboratory for testing whether multi-tier iterative planning algorithms converge under realistic industrial non-convexities, input-output networks, and strategic incentives.
  \item \textbf{The Practical Cooperative Coordination Suite (Phase 3):} Modular local-first tools allowing existing worker cooperatives, consumer buying clubs, community land trusts (CLTs), and participatory budgeting municipalities to coordinate procurement, shared equipment, and inter-cooperative production \emph{within} a surrounding capitalist market economy.
\end{enumerate}

\begin{table}[htbp]
\centering
\caption{Six Early Practical Applications and Their Institutional Requirements}
\label{tab:coop-apps}
\scriptsize
\setlength{\tabcolsep}{3.5pt}
\begin{tabularx}{\textwidth}{L{1.15in} L{1.35in} L{1.05in} Y}
\toprule
\textbf{Application} & \textbf{DPEP Modules Used} & \textbf{Viable Inside Market Economy?} & \textbf{Institutional / Ownership Changes Needed for Full Parecon Semantics} \\
\midrule
\textbf{1. Cooperative Joint Procurement} & Layer 2/3 demand aggregation across food co-ops \& worker co-ops. & \textbf{Yes (Immediately):} Legal B2B purchasing cooperative. & Requires pooled working-capital reserves to buffer seasonal supplier price swings. \\
\addlinespace[3pt]
\textbf{2. Shared Equipment \& Tool Trusts} & Stand-Alone Capital Asset ledger charging user-right fees ($P_k$). & \textbf{Yes:} Via a Cooperative Equipment Trust or CLT. & Member co-ops transfer title or long-term lease rights of heavy machinery/kitchens to the shared trust. \\
\addlinespace[3pt]
\textbf{3. Community Mutual-Aid Distribution} & Layer 3/4 Maxim-4 need allocation \& non-transferable credits. & \textbf{Yes:} Non-commercial mutual-aid association. & Bounded by the volume of goods produced by volunteer or grant-funded community labour. \\
\addlinespace[3pt]
\textbf{4. Inter-Cooperative Production Planning} & Quarterly ADMM/t\^atonnement clearing across 10--50 federated co-ops. & \textbf{Partially (Dual-Price Regime):} Needs Fiat Reserve Bridge. & External raw materials, rent, debt service, and taxes must still be settled in fiat currency (Eq.~\ref{eq:fiat-bridge}). \\
\addlinespace[3pt]
\textbf{5. Municipal Participatory Budgeting} & Layer 4 ``What-If'' explorer, Need Matrix \& \texttt{Prio3} voting. & \textbf{Yes:} Municipal public capital budgeting. & Vulnerable to municipal coalition reversals and fiscal austerity (as in Porto Alegre's 2017 suspension). \\
\addlinespace[3pt]
\textbf{6. Regional Co-op Supply Networks} & Multi-year investment ($\sigma_{\text{inv}}$) \& EFB worker reallocation. & \textbf{Partially} (e.g., Mondrag\'on / Cecosesola). & Full Parecon requires abolishing enterprise profit retention and equalising BJC empowerment across plants. \\
\bottomrule
\end{tabularx}
\end{table}

\section{Theory of Transition: Market Boundaries, Voluntary Exit, and Legal Exposure}

Any voluntary federation of cooperatives experimenting with participatory planning inside a surrounding market economy must navigate four structural realities of a \emph{pilot}, plus two that remain open even for a universal economy:
\begin{enumerate}[leftmargin=1.8em,itemsep=3pt]
  \item \textbf{Market-priced inputs, external debts, and taxes (the dual-price boundary):}
  Participating cooperatives cannot pay external steel suppliers, commercial landlords, bank mortgages, or national tax authorities in internal effort credits. Every pilot federation must enforce an explicit \textbf{External Fiat Clearing Constraint} alongside its internal Society Account:
  \begin{equation}
    \sum_{j \in \mathcal{J}} P_{\text{ext}}^{\text{fiat}} \cdot \bar{a}_j^{\text{export}} + D_{\text{fiat}}^{\text{dues}} \ge \sum_{j \in \mathcal{J}} P_{\text{ext}}^{\text{fiat}} \cdot \underline{a}_j^{\text{import}} + F^{\text{taxes}} + F^{\text{debt}} \label{eq:fiat-bridge}
  \end{equation}
  where external imports $\underline{a}_j^{\text{import}}$ carry an internal shadow price $\lambda_{\text{fiat}} P_{\text{ext}}^{\text{fiat}}$ reflecting the scarcity of external fiat revenue.
  \item \textbf{Strict Voluntarism, Right of Exit, and Non-Participants (Learning from Aragon 1936):}
  Informed by the conflict between Aragonese collectives and ``individualist'' smallholders (Casanova 2005), membership in any DPEP pilot or transitional federation must be strictly voluntary with a constitutionally codified \textbf{Right of Exit}: exiting individuals or cooperatives settle their net savings credits and refundable capital shares according to pre-agreed cooperative bylaws without confiscation, and non-participating artisans/households trade with the federation across transparent external interfaces.
  \item \textbf{Graduated sanctions without coercive labour discipline:}
  Following Ostrom (1990), failure to meet council commitments triggers: (a) automatic $SB_j/SC_j$ accounting adjustment, (b) UFB peer mediation, and (c) for persistent bad-faith breach, revocation of federation accreditation and access to shared federation inputs/credit---never forced labour. Losing access to federation inputs is mild in a voluntary pilot with an external market to exit into, and severe in a universal economy where exit is hard. The same rule therefore has two political meanings; a full-economy constitution needs a due-process and subsistence-access clause that a pilot can omit.
  \item \textbf{Jurisdiction-Specific Legal Exposure (Informational Analysis, Not Legal Advice):}
  \begin{itemize}[leftmargin=1.4em,itemsep=1pt]
    \item \emph{Competition / Antitrust Law (US Sherman Act \S 1; EU TFEU Art.\ 101; UK Competition Act 1998):} Independent corporations that coordinate prices and output quotas through an IFB risk prosecution as a horizontal price-fixing cartel unless structured under statutory \textbf{cooperative/agricultural exemptions} (e.g., US Capper-Volstead Act, EU producer-organisation rules) or incorporated as operating divisions of a \textbf{single secondary cooperative / common-ownership consortium}.
    \item \emph{Taxation of Multilateral Barter \& Clearing (US IRS Form 1099-B; EU/UK VAT):} Multilateral credit clearing is generally taxable at fair market value; pilot deployments must generate statutory fiat-equivalent tax reports where required by domestic law.
    \item \emph{Data Protection \& Right to Erasure (EU/UK GDPR Art.\ 17):} DPEP is designed to avoid immutable personal ledger data by keeping personal profiles local, deleting \texttt{Prio3} input shares when the challenge window closes (Section~\ref{sec:carry-forward}), and storing qualitative \texttt{Lists} behind mutable council pointers with tombstone erasure. Pilots additionally need the ethics and data-protection steps of Section~\ref{sec:ethics}.
  \end{itemize}
\end{enumerate}

\paragraph{Open even for a universal economy.}
\textbf{Foreign trade} (how a DPEP polity prices and rations extra-federation goods, manages a foreign-exchange constraint, and avoids reproducing centre--periphery extraction) and \textbf{defence} (who commands organised coercion, how it is paid, and how it is made compatible with sortition and dual-federation ratification) are not specified for either a pilot or a full economy. They remain open political questions, not protocol features.

% ==============================================================================
% CHAPTER 12: SIX-DISCIPLINE CRITICAL EXAMINATION
% ==============================================================================
\chapter{Six-Discipline Critical Examination and Remaining Weaknesses}

Table~\ref{tab:six-critics} lists the strongest remaining criticisms, grouped by discipline. For each it gives a severity, the evidence that would settle it, the phase that would produce that evidence, and whether it could stop the project. Severity has three levels: \textbf{Critical} (if confirmed, the programme as conceived should stop or be fundamentally redesigned), \textbf{Major} (forces narrowing of scope or a significant redesign), and \textbf{Moderate} (a known kind of fix exists, but it is not yet specified or tested). Nothing in the table is marked resolved.

{\scriptsize\setlength{\tabcolsep}{3pt}\renewcommand{\arraystretch}{1.15}
\begin{longtable}{L{1.45in} L{0.48in} L{1.55in} L{0.42in} L{0.42in}}
\caption{Strongest remaining criticisms: severity, settling evidence, and phase}\label{tab:six-critics}\\
\toprule
\textbf{Criticism} & \textbf{Severity} & \textbf{Evidence that would settle it} & \textbf{Phase} & \textbf{Could stop?} \\
\midrule
\endfirsthead
\multicolumn{5}{l}{\scriptsize\itshape Continued from previous page}\\
\toprule
\textbf{Criticism} & \textbf{Severity} & \textbf{Evidence that would settle it} & \textbf{Phase} & \textbf{Could stop?} \\
\midrule
\endhead
\midrule\multicolumn{5}{r}{\scriptsize\itshape Continued on next page}\\
\endfoot
\bottomrule
\endlastfoot
\multicolumn{5}{l}{\textbf{Participatory economics and democratic theory}}\\
Participation at scale: ordinary people will not sustain multi-round planning; turnout and attention will collapse after the novelty fades. & Critical & Completion, comprehension and fatigue (NASA-TLX) in Phase~1H; retention over a full plan year in pilots. & 1H, 3 & Yes \\
A technical elite forms around the software: facilitators, solver experts and articulate members set defaults and dominate amendments, re-creating a coordinator class. & Critical & Share of adopted amendments originated by the most active 20\%; who sets solver defaults; rotation of facilitation roles (Kill Gate~1H(e)). & 1H, 3 & Yes \\
Federations become central planning authorities in disguise (Kuch 2024). & Major & Count and scope of federation decisions in pilots; outcomes of Innovation Jury appeals (P7). & 3 & No (narrows) \\
Flat effort pay by default ($\epsilon_{\max}=0$) honours Maxim~3 only through hours and job balancing; rating intensity invites surveillance (O'Neill 2024). & Moderate & Cooperative pilots comparing flat and peer-rated pay (P6). & 3 & No \\
Cryptographic invisibility removes the 1991 channel of collective criticism of consumption. & Moderate & Members' stated preferences in pilots that offer both modes. & 3 & No \\
\addlinespace[2pt]
\multicolumn{5}{l}{\textbf{Operations research}}\\
No candidate solver is shown to clear Leontief, lumpy and endogenous-income economies; 3b's theorem covers only fixed-budget convex cases, and 3d has no convergence result. & Critical & Kill Gate~1 on the pre-registered set, against an oracle. & 1 & Yes \\
One-good smoothing results do not transfer to interdependent markets. & Major & Lumpy Leontief families in the instance set. & 1 & No (narrows) \\
Rounds-to-clear may grow with the number of councils. & Major & Log--log scaling slope from 50 to 1,000 councils. & 1 & Yes \\
\addlinespace[2pt]
\multicolumn{5}{l}{\textbf{Distributed systems}}\\
Round cost at $10^6$ (0.24--1.46~GB archives, 40k--200k signature checks, $10^6$ report verifications per aggregator) may not fit a usable round time. & Major & Kill Gate~2 measurements at $\ge 10^4$ with extrapolation from measured per-item costs. & 2 & Yes (for D; A remains) \\
Annual committee capture lasts a year; a 30\% corrupt pool gives a 21\% chance of an unsafe draw. & Major & Rotation frequency; operation of the interim threshold beacon (Section~\ref{sec:pq-urgency}); a post-quantum beacon (open). Not settled by an experiment. & 2 & No \\
\addlinespace[2pt]
\multicolumn{5}{l}{\textbf{Cryptography}}\\
Aggregate correctness needs all aggregators; the public budget check limits only gross overspend. & Major & Pilot use of non-nested aggregators; audit. & 3 & No \\
No construction for budget credentials and one-report-per-person; the budget circuit is unspecified and price-dependent. & Critical (for household crypto only) & A specified, audited construction: an interim classical credential with threshold issuance covers one report per person (Section~\ref{sec:pq-urgency}); budget binding and the circuit remain open. Until then the council-level trust model. & 3 & Yes (household mode only) \\
Carry-forward cannot re-aggregate reports; milestone-only revision freezes household demand between milestones. & Moderate & Planning-accuracy cost in simulation (P5) and pilots. & 1, 3 & No \\
\addlinespace[2pt]
\multicolumn{5}{l}{\textbf{Cooperative practice}}\\
30--60 self-selected volunteers say little about a general population. & Major & Larger, more representative pilots after Phase~1H. & 3 & No \\
Sanctions that cut federation inputs are mild in a pilot and severe in a universal economy. & Moderate & Constitutional due-process and subsistence clause (open). & -- & No \\
\addlinespace[2pt]
\multicolumn{5}{l}{\textbf{Sceptical economist}}\\
Ratchet effects and incentive limits (Weitzman, Hurwicz) remain; public blueprints do not remove them. & Major & Strategic-misreporting family in Phase~1 (diagnostic); pilots. & 1, 3 & No \\
The labour floor $p^{\text{lab}}\ge w_{\text{base}}$ blocks surplus clearing. & Moderate & Phase~1 instances with labour surplus. & 1 & No \\
Foreign trade and defence are unspecified. & Major & Not settled by any experiment here; a political design question. & -- & No (scope limit) \\
\end{longtable}}

% ==============================================================================
% APPENDIX A: SCIENTIFIC CLAIMS, EVIDENCE, AND VALIDATION STATUS
% ==============================================================================
\appendix
\chapter{Scientific Claims, Evidence, and Validation Status Register}

Each claim has eight separate fields: identifier, claim, classification, status (first line) and evidence, assumptions, limitations and required validation (second line). \textbf{Classifications:} Established (peer-reviewed theorem or standard), Source-supported (checked in primary literature, audited code or a measurement), Analytical (derived from models or arithmetic in this monograph), Conjectural (hypothesis or judgement), Experimental target (a goal to be tested). \textbf{Statuses:} Verified, Partially verified, Unverified, Refuted.

{\scriptsize\setlength{\tabcolsep}{3.0pt}\renewcommand{\arraystretch}{1.05}
\begin{longtable}{L{0.92in} L{1.42in} L{1.07in} L{1.12in}}
\caption{Master Register of Scientific and Engineering Claims}
\label{tab:claims-reg} \\
\toprule
\textbf{Claim ID / Evidence} & \textbf{Claim / Assumptions} & \textbf{Classification / Limitations} & \textbf{Status / Validation} \\
\midrule
\endfirsthead
\multicolumn{4}{l}{\scriptsize\itshape Continued from previous page} \\
\toprule
\textbf{Claim ID / Evidence} & \textbf{Claim / Assumptions} & \textbf{Classification / Limitations} & \textbf{Status / Validation} \\
\midrule
\endhead
\bottomrule
\endlastfoot
\textbf{CLM-ECON-01} & Economic convergence (reference): the continuous-time FMPE gradient process converges. & \emph{Class:} Established & \emph{Status:} \textbf{Verified (as a cited theorem; not a DPEP result)} \\*
\emph{Evidence:} Arrow--Hurwicz--Uzawa (1958); Heal (1973), Tier~3. & \emph{Assumptions:} Strict convexity; gross substitutability. & \emph{Limitations:} Fails on Scarf-type and non-convex economies. & \emph{Validation:} Reference only. \\
\addlinespace[1pt]\midrule
\textbf{CLM-ECON-02} & \texttt{dep1ex61--100.csv}: Year~1 reaches 5\% in 11.85 rounds and 3\% in 19.20; Year~2 5\% in 6.58. & \emph{Class:} Source-supported & \emph{Status:} \textbf{Verified} \\*
\emph{Evidence:} Audit script\footnote{\path{research/analysis/verify_r2.py} in the project repository.} on \texttt{pequod-cljs} @ \texttt{8bd46b5c}. & \emph{Assumptions:} Fixed \$5{,}000 incomes; Cobb--Douglas FOCs. & \emph{Limitations:} Other subsets differ (01--50: 11.18); no wall-clock log. & \emph{Validation:} Reproduce in the Phase~1 simulator. \\
\addlinespace[1pt]\midrule
\textbf{CLM-ECON-03} & Eq.~(\ref{eq:pequod-rule}) plus trust regions does not clear Scarf's cycle or a single lumpy plant. & \emph{Class:} Analytical & \emph{Status:} \textbf{Verified} \\*
\emph{Evidence:} Simulation script\footnote{\path{research/analysis/verify_economics.py} in the project repository.}: min imbalance 10.5\% and 50.4\%. & \emph{Assumptions:} The stress-test economies. & \emph{Limitations:} Not a general impossibility proof. & \emph{Validation:} Instance families 3--4 in Phase~1. \\
\addlinespace[1pt]\midrule
\textbf{CLM-ECON-04} & With the toothache increment, $\sum b + c/k > 1$ for every one of 750,000 worker-council records in the 70 generators. & \emph{Class:} Source-supported & \emph{Status:} \textbf{Partially verified} \\*
\emph{Evidence:} Dataset audit script\footnote{\path{research/analysis/verify_pequod_datasets.py} in the project repository.}; LFS objects SHA-256 checked. & \emph{Assumptions:} 20\% of councils selected at run time. & \emph{Limitations:} Perverse supply predicted, not observed in CSVs; Clojure not run. & \emph{Validation:} Run \texttt{pequod-cljs} with toothaches. \\
\addlinespace[1pt]\midrule
\textbf{CLM-ECON-05} & One good, $K=20$ staggered plants: min imbalance 0.40\% over 500 rounds; last-100-round max 8.4\%. & \emph{Class:} Analytical & \emph{Status:} \textbf{Partially verified} \\*
\emph{Evidence:} Simulation script\footnote{\path{research/analysis/verify_r2.py} in the project repository.}; sweep median 6.0\%, worst 9.5\%. & \emph{Assumptions:} $a\in[60,80]$; thresholds in $[8.5,11.5]$. & \emph{Limitations:} One good; not Leontief. & \emph{Validation:} Lumpy Leontief family in Phase~1. \\
\addlinespace[1pt]\midrule
\textbf{CLM-ECON-06} & $\sum_u B_u = (1-\sigma_{\text{inv}})Y^{\text{GDP}}$, and $\sigma_{\text{inv}}Y^{\text{GDP}} = p\cdot D^{\text{inv}}$ at a cleared plan. & \emph{Class:} Analytical & \emph{Status:} \textbf{Partially verified} \\*
\emph{Evidence:} Algebra in Section~4.2 and Table~\ref{tab:macro-sheet}. & \emph{Assumptions:} $Y^{\text{allow}}$ is cash transfers only. & \emph{Limitations:} No numerical test on a full Society Account ledger. & \emph{Validation:} Phase~1 ledger test. \\
\addlinespace[1pt]\midrule
\textbf{CLM-ECON-07} & Algorithm~3b (Deng et al.\ Algorithm~4 with slack) converges on fixed-budget convex problems. & \emph{Class:} Established & \emph{Status:} \textbf{Partially verified} \\*
\emph{Evidence:} Deng et al.\ (2017) Thm.~2.3--2.4 (read in full); test~B in the analysis script.\footnote{\path{research/analysis/verify_r3.py} in the project repository.} & \emph{Assumptions:} Convex $\mathcal{T}_c$; fixed budgets; $P_i$ condition (2.11). & \emph{Limitations:} Not endogenous income, non-convexity or dual renormalisation. & \emph{Validation:} Phase~1 convex families. \\
\addlinespace[1pt]\midrule
\textbf{CLM-ECON-08} & Some candidate (3a--3d) clears full DPEP instances (Leontief, lumpy, endogenous income, strategic). & \emph{Class:} Experimental target & \emph{Status:} \textbf{Unverified} \\*
\emph{Evidence:} None yet. & \emph{Assumptions:} Pre-registered instance set. & \emph{Limitations:} Toy tests only; 3d has no convergence result. & \emph{Validation:} Kill Gate~1. \\
\addlinespace[1pt]\midrule
\textbf{CLM-ECON-09} & One-sided feasibility ($D\le S$) is an adequate success criterion. & \emph{Class:} Analytical & \emph{Status:} \textbf{Refuted} \\*
\emph{Evidence:} Single plant: feasible at round 3 with 42\% unused output (reference test script)\footnote{\path{research/analysis/reference_tests_r3.py} in the project repository.}. & \emph{Assumptions:} -- & \emph{Limitations:} -- & \emph{Validation:} Replaced by Definition~\ref{def:success}. \\
\addlinespace[1pt]\midrule
\textbf{CLM-TECH-01} & Leaf transport plus BFT takes 9--27~s per round at $N=10^6$. & \emph{Class:} Analytical & \emph{Status:} \textbf{Refuted} \\*
\emph{Evidence:} Assumed 5.2~kB proposals; reworked sizes give 0.24--1.46~GB per round (Table~\ref{tab:proposal-sizes}). & \emph{Assumptions:} -- & \emph{Limitations:} Omitted report verification and signature checks. & \emph{Validation:} Replaced by Kill Gate~2 thresholds. \\
\addlinespace[1pt]\midrule
\textbf{CLM-TECH-02} & Computational scalability: \texttt{Prio3} verification costs $\approx 0.25$~ms (coarse) to 3.5~ms (dense) per report per aggregator. & \emph{Class:} Source-supported & \emph{Status:} \textbf{Partially verified} \\*
\emph{Evidence:} libprio 0.18.1 benchmark (Table~\ref{tab:prio3-bench}). & \emph{Assumptions:} One desktop core; VDAF draft~18 code. & \emph{Limitations:} No phone or Raspberry Pi measurement; network excluded. & \emph{Validation:} Device benchmarks, Kill Gate~3. \\
\addlinespace[1pt]\midrule
\textbf{CLM-TECH-03} & \texttt{Prio3SumVec} report sizes follow the VDAF size formula. & \emph{Class:} Source-supported & \emph{Status:} \textbf{Verified} \\*
\emph{Evidence:} The size model\footnote{\path{research/analysis/parameters.py} in the project repository.} matches libprio byte-for-byte on 9 configurations. & \emph{Assumptions:} libprio chunk-length rule. & \emph{Limitations:} Custom budget circuit sizes unknown. & \emph{Validation:} Recompute when the circuit is specified. \\
\addlinespace[1pt]\midrule
\textbf{CLM-TECH-04} & Mobile feasibility: $\le 150$~kB per ordinary round and $\le 450$~kB with a full quorum check; RAM $\le 150$~MB. & \emph{Class:} Experimental target & \emph{Status:} \textbf{Unverified} \\*
\emph{Evidence:} Byte budget analytical (Section~\ref{sec:sortition}). & \emph{Assumptions:} Coarse 55~kB report; coordinator attestation. & \emph{Limitations:} No device test. & \emph{Validation:} Kill Gate~3. \\
\addlinespace[1pt]\midrule
\textbf{CLM-TECH-05} & Operating cost: hosting 100 servers $\approx$ \$3k--\$6k/month. & \emph{Class:} Analytical & \emph{Status:} \textbf{Unverified} \\*
\emph{Evidence:} \$30--\$60 per server-month (assumed). & \emph{Assumptions:} Commodity servers. & \emph{Limitations:} Excludes audits, facilitation, research and energy. & \emph{Validation:} Phase~2 and 3 cost accounting. \\
\addlinespace[1pt]\midrule
\textbf{CLM-TECH-06} & Architecture D outranks A, B and C at baseline weights. & \emph{Class:} Conjectural & \emph{Status:} \textbf{Partially verified} \\*
\emph{Evidence:} Expert scores; arithmetic and perturbation test in the analysis script.\footnote{\path{research/analysis/verify_r3.py} in the project repository.} & \emph{Assumptions:} Table~\ref{tab:arch-comp} scores. & \emph{Limitations:} A overtakes D when simplicity, throughput and cost carry $\ge 63\%$ of the weight. & \emph{Validation:} Phases 2--3 replace judgements with measurements. \\
\addlinespace[1pt]\midrule
\textbf{CLM-TECH-07} & Component maturity: every function except threshold signatures, anonymous credentials, the beacon, capability-certificate formats and post-quantum HPKE has a standardised, deployed primitive. & \emph{Class:} Source-supported & \emph{Status:} \textbf{Partially verified} \\*
\emph{Evidence:} Table~\ref{tab:maturity}; maturity dossier\footnote{\path{reports/STACK_MATURITY.md} in the project repository.} (web checks, November 2026). & \emph{Assumptions:} Ratings are judgements on the M0--M5 scale. & \emph{Limitations:} Primitive maturity is not library or combination maturity; nine combinations are original. & \emph{Validation:} Re-check before each gate; Kill Gate~3 audit. \\
\addlinespace[1pt]\midrule
\textbf{CLM-SEC-01} & Privacy: \texttt{Prio3} hides an individual report if at least one aggregator is honest. & \emph{Class:} Source-supported & \emph{Status:} \textbf{Partially verified} \\*
\emph{Evidence:} VDAF-22 \S 9 (full text). & \emph{Assumptions:} Non-colluding aggregators; $m_{\min}=30$. & \emph{Limitations:} Budget circuit unspecified; differencing between releases. & \emph{Validation:} Phase~3 audit. \\
\addlinespace[1pt]\midrule
\textbf{CLM-SEC-02} & Censorship resistance: Two-Root aggregation with receipts detects omitted or tampered proposals. & \emph{Class:} Analytical & \emph{Status:} \textbf{Unverified} \\*
\emph{Evidence:} Design argument (Section~\ref{sec:two-root}). & \emph{Assumptions:} $\ge 1$ honest relay; clients keep receipts. & \emph{Limitations:} Does not bind aggregator shares. & \emph{Validation:} Kill Gate~2(a). \\
\addlinespace[1pt]\midrule
\textbf{CLM-SEC-03} & Post-quantum security: the FIPS 203/204/205 primitives are standardised. & \emph{Class:} Established & \emph{Status:} \textbf{Verified} \\*
\emph{Evidence:} NIST FIPS 203--205 (2024). & \emph{Assumptions:} Module-LWE/SIS and hash assumptions. & \emph{Limitations:} No whole-system PQ claim; BLAKE3 128-bit; VDAF no PQ claim. & \emph{Validation:} Phase~3 review. \\
\addlinespace[1pt]\midrule
\textbf{CLM-SEC-04} & Cryptographic integrity: aggregate correctness requires every aggregator to be honest. & \emph{Class:} Source-supported & \emph{Status:} \textbf{Verified} \\*
\emph{Evidence:} VDAF-22 \S 9 quotation. & \emph{Assumptions:} Standard \texttt{Prio3}. & \emph{Limitations:} The public budget check limits only gross overspend. & \emph{Validation:} Non-nested aggregators in pilots. \\
\addlinespace[1pt]\midrule
\textbf{CLM-SEC-05} & Sybil resistance (partial): INV-4B bounds credited hours given a true roster. & \emph{Class:} Analytical & \emph{Status:} \textbf{Partially verified} \\*
\emph{Evidence:} Eq.~(\ref{eq:effort-cap}). & \emph{Assumptions:} Roster and $SB/SC$ honest. & \emph{Limitations:} Falsified rosters and carousel fraud need physical audit. & \emph{Validation:} Pilot audits. \\
\addlinespace[1pt]\midrule
\textbf{CLM-SEC-06} & Sybil resistance (household): one report per person without a population register. & \emph{Class:} Conjectural & \emph{Status:} \textbf{Unverified} \\*
\emph{Evidence:} No post-quantum construction (Section~\ref{sec:bu-cred}); interim classical credential adopted (Section~\ref{sec:pq-urgency}). & \emph{Assumptions:} An anonymous credential with threshold issuance and nullifiers. & \emph{Limitations:} Budget binding open; interim scheme not post-quantum. & \emph{Validation:} Research; council-level trust meanwhile. \\
\addlinespace[1pt]\midrule
\textbf{CLM-SEC-07} & An $n=128$ annual draw with $p=20\%$ corrupt gives a 0.022\% unsafe committee. & \emph{Class:} Analytical & \emph{Status:} \textbf{Verified} \\*
\emph{Evidence:} Exact binomial computation.\footnote{\path{research/analysis/verify_r2.py} in the project repository.} & \emph{Assumptions:} i.i.d.\ draws; failure at $k\ge 43$. & \emph{Limitations:} $p=30\%$ gives 21\%; capture lasts a year. & \emph{Validation:} Beacon design. \\
\addlinespace[1pt]\midrule
\textbf{CLM-SEC-08} & \texttt{Prio3} reports cannot be carried forward into later rounds' aggregates. & \emph{Class:} Source-supported & \emph{Status:} \textbf{Verified} \\*
\emph{Evidence:} VDAF-22 \S 7.2.3, \S 9.4. & \emph{Assumptions:} Standard \texttt{Prio3}. & \emph{Limitations:} Aggregate-level carry-forward is still possible. & \emph{Validation:} -- \\
\addlinespace[1pt]\midrule
\textbf{CLM-SEC-09} & An off-the-shelf BFT engine accepts post-quantum validator keys. & \emph{Class:} Source-supported & \emph{Status:} \textbf{Verified} \\*
\emph{Evidence:} Cosmos SDK documentation (CometBFT 0.40, Cosmos SDK 0.55). & \emph{Assumptions:} Releases as documented. & \emph{Limitations:} Pure ML-DSA-65, not hybrid; $\approx 331$~kB of signatures per block at 100 validators; release dates not confirmed. & \emph{Validation:} Kill Gate~2. \\
\addlinespace[1pt]\midrule
\textbf{CLM-SEC-10} & An interim classical credential keeps past showings unlinkable after a CRQC. & \emph{Class:} Conjectural & \emph{Status:} \textbf{Unverified} \\*
\emph{Evidence:} None for a chosen scheme; RFC~9474 states no assumption for blindness. & \emph{Assumptions:} Statistically zero-knowledge presentation and statistical blinding. & \emph{Limitations:} If false, a CRQC reveals who submitted, not what. & \emph{Validation:} Scheme selection before Phase~3. \\
\addlinespace[1pt]\midrule
\textbf{CLM-SEC-11} & Canary batches detect indiscriminate tampering with aggregate shares. & \emph{Class:} Analytical & \emph{Status:} \textbf{Unverified} \\*
\emph{Evidence:} Argument in Section~\ref{sec:prio3}. & \emph{Assumptions:} Per-batch aggregate shares; batches indistinguishable to aggregators. & \emph{Limitations:} Misses tampering aimed at batches known to be real. & \emph{Validation:} Fault injection in Phase~3. \\
\addlinespace[1pt]\midrule
\textbf{CLM-POL-01} & Democratic scalability: ordinary participants complete, understand and genuinely shape multi-round planning. & \emph{Class:} Experimental target & \emph{Status:} \textbf{Unverified} \\*
\emph{Evidence:} None yet. & \emph{Assumptions:} Phase~1H design. & \emph{Limitations:} $n=30$--60 self-selected volunteers. & \emph{Validation:} Kill Gate~1H. \\
\addlinespace[1pt]\midrule
\textbf{CLM-POL-02} & No technical elite dominates proposals and defaults. & \emph{Class:} Experimental target & \emph{Status:} \textbf{Unverified} \\*
\emph{Evidence:} None yet. & \emph{Assumptions:} Influence measure of Kill Gate~1H(e). & \emph{Limitations:} Small, short probe. & \emph{Validation:} Kill Gate~1H and pilots. \\
\addlinespace[1pt]\midrule
\textbf{CLM-POL-03} & Neighbourhood councils of $\sim 50$--200 are workable. & \emph{Class:} Conjectural & \emph{Status:} \textbf{Unverified} \\*
\emph{Evidence:} Speaking-time arithmetic. & \emph{Assumptions:} One shared meeting hour. & \emph{Limitations:} Not a finding of Mansbridge or Dahl \& Tufte. & \emph{Validation:} Pilots. \\
\end{longtable}}

% ==============================================================================
% BACK MATTER: COMPREHENSIVE SOURCE BIBLIOGRAPHY
% ==============================================================================
\backmatter
\chapter{Comprehensive Source Bibliography}

Each entry's bibliographic identity was checked in November 2026 against Crossref (by DOI or bibliographic query), a publisher or journal page, or the IETF Datatracker.\footnote{Checked by \path{research/analysis/audit_bibliography.py}; results in \path{research/SOURCE_AUDIT.md}, both in the project repository.} That is not a claim that every page citation was re-read; access levels are in the front-matter notice.

\begin{enumerate}[leftmargin=1.8em,itemsep=4pt,label={[\arabic*]}]
  \item \textbf{Albert, Michael, and Robin Hahnel (1991a).} \emph{The Political Economy of Participatory Economics}. Princeton, NJ: Princeton University Press. 132 pp. ISBN 978-0-691-00384-9. DOI: \url{https://doi.org/10.2307/j.ctv131bvr2}. Print TOC/pagination scan preserved at \texttt{research/sources/pepe\_dandelon.pdf}; full text analyzed via ZNetwork (\url{https://znetwork.org/wp-content/uploads/zbooks/htdocs/books/intro.htm} and \texttt{1.htm}--\texttt{6.htm}).
  \item \textbf{Albert, Michael, and Robin Hahnel (1991b).} \emph{Looking Forward: Participatory Economics for the Twenty First Century}. Boston, MA: South End Press. 154 pp. ISBN 978-0-89608-405-6. Full text analyzed via ZNetwork (\url{https://znetwork.org/wp-content/uploads/zbooks/htdocs/ParEcon/lookingforward/toc.htm}).
  \item \textbf{Hahnel, Robin, and Michael Albert (1990).} \emph{Quiet Revolution in Welfare Economics}. Princeton, NJ: Princeton University Press. 440 pp. DOI (digital edition): \url{https://doi.org/10.1515/9781400887057}.
  \item \textbf{Albert, Michael (2003).} \emph{Parecon: Life After Capitalism}. London \& New York: Verso Books. 311 pp. ISBN 978-1-85984-698-8. Full text analyzed via ZCommunications/ZNetwork (\url{http://books.zcommunications.org/books/pareconv/}) and PDF edition; analytical summary preserved at \texttt{research/sources/albert\_2003\_parecon\_analysis.md}.
  \item \textbf{Hahnel, Robin (2021).} \emph{Democratic Economic Planning} (with contributions by Peter Bohmer, Savvina Chowdhury, Allison Kerkhoff, Mitchell Szczepanczyk, and Michael Weisdorf). Abingdon \& New York: Routledge. 358 pp. ISBN 978-1-032-00332-0. DOI: \url{https://doi.org/10.4324/9781003173700}.
  \item \textbf{Hahnel, Robin (2022).} \emph{A Participatory Economy}. Chico, CA: AK Press. 160 pp. ISBN 978-1-84935-484-4.
  \item \textbf{Sandstr\"om, Anders (2020).} \emph{Anarchist Accounting: Accounting Principles for a Democratic Economy}. Abingdon \& New York: Routledge. 140 pp. ISBN 978-0-367-47703-5. DOI: \url{https://doi.org/10.4324/9781003036005}.
  \item \textbf{Sandstr\"om, Anders (2023).} ``An Accounting System for a Libertarian Socialist Economy.'' Participatory Economy Project, February 14, 2023. \url{https://participatoryeconomy.org/an-accounting-system-for-a-libertarian-socialist-economy/}.
  \item \textbf{Sandstr\"om, Anders (2025).} \emph{The Participatory Economy: Adding Realism to Democratic Planning}. Abingdon \& New York: Routledge. ISBN 978-1-041-22413-6. DOI: \url{https://doi.org/10.4324/9781003732174} (Crossref dates the online edition 2026; subtitle not confirmed there).
  \item \textbf{Szczepanczyk, Mitchell (2023).} ``Pseudocode and Algorithms for Computer Simulations of Democratically Planned Economies.'' \emph{Journal of Information Economics} 1, no. 3: 43--54. DOI: \url{https://doi.org/10.58567/jie01030004}. Open-access CC BY 4.0 full text and \texttt{pequod-cljs} source code (\url{https://github.com/msszczep/pequod-cljs}) preserved in \texttt{research/sources/}.
  \item \textbf{Participatory Economy Project (2023).} ``Participatory Planning: User Story Mapping.'' \url{https://participatoryeconomy.org/project/participatory-planning-user-story-mapping/}.
  \item \textbf{Arrow, Kenneth J., Leonid Hurwicz, and Hirofumi Uzawa (1958).} \emph{Studies in Linear and Nonlinear Programming}. Stanford, CA: Stanford University Press.
  \item \textbf{Heal, Geoffrey M. (1973).} \emph{The Theory of Economic Planning}. Amsterdam: North-Holland Publishing Co.
  \item \textbf{Henry, Claude (1972).} ``Differential Equations with Discontinuous Right-Hand Side for Planning Procedures.'' \emph{Journal of Economic Theory} 4, no. 3: 545--551. DOI: \url{https://doi.org/10.1016/0022-0531(72)90138-x}.
  \item \textbf{Hurwicz, Leonid (1973).} ``The Design of Mechanisms for Resource Allocation.'' \emph{American Economic Review} 63, no. 2: 1--30.
  \item \textbf{Hayek, Friedrich A. (1945).} ``The Use of Knowledge in Society.'' \emph{American Economic Review} 35, no. 4: 519--530.
  \item \textbf{Lange, Oskar, and Fred M. Taylor (1938/1964).} \emph{On the Economic Theory of Socialism}, ed.\ Benjamin E. Lippincott. New York: Monthly Review Press.
  \item \textbf{Nove, Alec (1983).} \emph{The Economics of Feasible Socialism}. London: George Allen \& Unwin.
  \item \textbf{Buchanan, Allen (1985).} \emph{Ethics, Efficiency, and the Market}. Totowa, NJ: Rowman and Littlefield.
  \item \textbf{Cockshott, W. Paul, and Allin Cottrell (1993).} \emph{Towards a New Socialism}. Nottingham: Spokesman Books.
  \item \textbf{Devine, Pat (1988).} \emph{Democracy and Economic Planning: The Political Economy of a Self-Governing Society}. Boulder, CO: Westview Press.
  \item \textbf{Schweickart, David (1993).} \emph{Against Capitalism}. Cambridge: Cambridge University Press.
  \item \textbf{Laibman, David (2014).} ``Horizontalism and Idealism in Socialist Imagination: An Appraisal of the Participatory Economy.'' \emph{Science \& Society} 78, no. 2: 207--234. DOI: \url{https://doi.org/10.1521/siso.2014.78.2.207}. Abstract consulted.
  \item \textbf{Hahnel, Robin (2014).} ``Response to David Laibman's `Appraisal of the Participatory Economy'.'' \emph{Science \& Society} 78, no. 3: 379--389. DOI: \url{https://doi.org/10.1521/siso.2014.78.3.379}. Bibliographic record only.
  \item \textbf{Hahnel, Robin, and Erik Olin Wright (2016).} \emph{Alternatives to Capitalism: Proposals for a Democratic Economy}. London: Verso Books.
  \item \textbf{O'Neill, John (2024).} ``Pluralism, Ecology and Planning.'' \emph{Erasmus Journal for Philosophy and Economics} 17, no. 2: 27--48, \url{https://doi.org/10.23941/ejpe.v17i2.884}; and ``Incommensurability, Environment and Planning: A Response to Hahnel's Reply,'' ibid.: 117--124, \url{https://doi.org/10.23941/ejpe.v17i2.933}. Issue PDF: \url{https://ejpe.org/journal/issue/download/35/67}.
  \item \textbf{Kuch, Hannes (2024).} ``Centralization Tendencies in Participatory Planning.'' \emph{Erasmus Journal for Philosophy and Economics} 17, no. 2: 49--72. DOI: \url{https://doi.org/10.23941/ejpe.v17i2.885}.
  \item \textbf{Dapprich, Jan Philipp (2024).} ``Democratic Theory and Economic Planning: A Critique of Participatory Economics.'' \emph{Erasmus Journal for Philosophy and Economics} 17, no. 2: 73--93. DOI: \url{https://doi.org/10.23941/ejpe.v17i2.867}.
  \item \textbf{Dolgoff, Sam, ed.\ (1974).} \emph{The Anarchist Collectives: Workers' Self-Management in the Spanish Revolution, 1936--1939}. New York: Free Life Editions. \url{https://theanarchistlibrary.org/library/sam-dolgoff-editor-the-anarchist-collectives}.
  \item \textbf{Leval, Gaston (1975).} \emph{Collectives in the Spanish Revolution}, trans.\ Vernon Richards. London: Freedom Press.
  \item \textbf{Mintz, Frank (2013 [1970/1977]).} \emph{Anarchism and Workers' Self-Management in Revolutionary Spain}, trans.\ Paul Sharkey. Oakland, CA: AK Press.
  \item \textbf{Marx, Karl (1871).} \emph{The Civil War in France}. London: General Council of the International Workingmen's Association.
  \item \textbf{Merriman, John (2014).} \emph{Massacre: The Life and Death of the Paris Commune of 1871}. New York: Basic Books.
  \item \textbf{Hansen, Mogens Herman (1991).} \emph{The Athenian Democracy in the Age of Demosthenes: Structure, Principles, and Ideology}. Norman: University of Oklahoma Press.
  \item \textbf{Freeman, Jo (1972).} ``The Tyranny of Structurelessness.'' \emph{The Second Wave} 2, no. 1 / \emph{Berkeley Journal of Sociology} 17: 151--164.
  \item \textbf{Graeber, David (2013).} \emph{The Democracy Project: A History, a Crisis, a Movement}. New York: Spiegel \& Grau.
  \item \textbf{Ostrom, Elinor (1990).} \emph{Governing the Commons: The Evolution of Institutions for Collective Action}. Cambridge: Cambridge University Press.
  \item \textbf{Whyte, William Foote, and Kathleen King Whyte (1991).} \emph{Making Mondrag\'on: The Growth and Dynamics of the Worker Cooperative Complex}, 2nd ed. Ithaca, NY: ILR Press.
  \item \textbf{Goulart, Paul J., and Yuwen Chen (2026).} ``Clarabel: An Interior-Point Solver for Conic Programs with Quadratic Objectives.'' \emph{Mathematical Programming Computation}. DOI: \url{https://doi.org/10.1007/s12532-026-00320-7}. Software: \url{https://github.com/oxfordcontrol/Clarabel.rs}.
  \item \textbf{O'Connor, Jack, Jean-Philippe Aumasson, Samuel Neves, and Zooko Wilcox-O'Hearn (2020).} \emph{BLAKE3 specification}. \url{https://github.com/BLAKE3-team/BLAKE3-specs}. ``BLAKE3 targets 128-bit security for all of its security goals.'' Accessed November 2026.
  \item \textbf{NIST (2024a, 2024b, 2024c).} \emph{FIPS 203 (ML-KEM)}, \emph{FIPS 204 (ML-DSA)}, and \emph{FIPS 205 (SLH-DSA)}. National Institute of Standards and Technology, August 13, 2024.
  \item \textbf{Kleppmann, Martin, et al.\ (2019).} ``Local-First Software: You Own Your Data, in spite of the Cloud.'' \emph{ACM Onward! '19}, pp.~154--178.
  \item \textbf{Shapiro, Marc, et al.\ (2011).} ``Conflict-Free Replicated Data Types.'' \emph{SSS 2011}, LNCS 6976, pp.~386--400.
  \item \textbf{Buchman, Ethan, Jae Kwon, and Zarko Milosevic (2018).} ``The Latest Gossip on BFT Consensus.'' \emph{arXiv:1807.04938}.
  \item \textbf{Boyd, Stephen, et al.\ (2011).} ``Distributed Optimization and Statistical Learning via the Alternating Direction Method of Multipliers.'' \emph{Foundations and Trends in Machine Learning} 3, no. 1: 1--122, \url{https://doi.org/10.1561/2200000016}; \textbf{Deng, Wei, Ming-Jun Lai, Zhimin Peng, and Wotao Yin (2017).} ``Parallel Multi-Block ADMM with $o(1/k)$ Convergence.'' \emph{Journal of Scientific Computing} 71, no. 2: 712--736, \url{https://doi.org/10.1007/s10915-016-0318-2} (read in the arXiv version, 1312.3040v2); \textbf{Hong, Mingyi, Zhi-Quan Luo, and Meisam Razaviyayn (2016).} ``Convergence Analysis of Alternating Direction Method of Multipliers for a Family of Nonconvex Problems.'' \emph{SIAM Journal on Optimization} 26, no. 1: 337--364, \url{https://doi.org/10.1137/140990309}.
  \item \textbf{Scarf, Herbert (1960).} ``Some Examples of Global Instability of the Competitive Equilibrium.'' \emph{International Economic Review} 1, no. 3: 157--172; \textbf{Sonnenschein, Hugo (1973).} ``Do Walras' Identity and Continuity Characterize the Class of Community Excess Demand Functions?'' \emph{Journal of Economic Theory} 6, no. 4: 345--354.
  \item \textbf{Starr, Ross M.\ (1969).} ``Quasi-Equilibria in Markets with Non-Convex Preferences (Appendix 2: The Shapley--Folkman Theorem).'' \emph{Econometrica} 37, no. 1: 25--38.
  \item \textbf{Hurwicz, Leonid (1972).} ``On Informationally Decentralized Systems.'' In \emph{Decision and Organization: A Volume in Honor of Jacob Marschak}, ed.\ C.~B.\ McGuire and Roy Radner, pp.~297--336. Amsterdam: North-Holland Publishing Co.
  \item \textbf{Weitzman, Martin L.\ (1974).} ``Prices vs.\ Quantities.'' \emph{Review of Economic Studies} 41, no. 4: 477--491; \textbf{Weitzman, Martin L.\ (1980).} ``The `Ratchet Principle' and Performance Incentives.'' \emph{Bell Journal of Economics} 11, no. 1: 302--308.
  \item \textbf{Kornai, J\'anos (1980).} \emph{Economics of Shortage}. Amsterdam: North-Holland Publishing Co.; \textbf{Kornai, J\'anos (1986).} ``The Soft Budget Constraint.'' \emph{Kyklos} 39, no. 1: 3--30.
  \item \textbf{Casanova, Juli\'an (2005).} \emph{Anarchism, the Republic and Civil War in Spain: 1931--1939}, trans.\ Andrew Dowling and Graham Pollok. London \& New York: Routledge.
  \item \textbf{Seidman, Michael (1991).} \emph{Workers Against Work: Labour in Paris and Barcelona during the Popular Fronts}. Berkeley: University of California Press; \textbf{Seidman, Michael (2002).} \emph{Republic of Egos: A Social History of the Spanish Civil War}. Madison: University of Wisconsin Press.
  \item \textbf{Finley, M.~I.\ (1981).} \emph{Economy and Society in Ancient Greece}, ed.\ Brent D.\ Shaw and Richard P.\ Saller. London: Chatto \& Windus; \textbf{Ste.\ Croix, G.~E.~M.\ de (1981).} \emph{The Class Struggle in the Ancient Greek World}. Ithaca, NY: Cornell University Press.
  \item \textbf{Mansbridge, Jane J.\ (1980).} \emph{Beyond Adversary Democracy}. New York: Basic Books; \textbf{Dahl, Robert A., and Edward R.\ Tufte (1973).} \emph{Size and Democracy}. Stanford, CA: Stanford University Press.
  \item \textbf{Corrigan-Gibbs, Henry, and Dan Boneh (2017).} ``Prio: Private, Robust, and Scalable Computation of Aggregate Statistics.'' \emph{USENIX NSDI '17}, pp.~259--282; \textbf{Boneh, Dan, et al.\ (2019).} ``Zero-Knowledge Proofs on Secret-Shared Data via Fully Linear PCPs.'' \emph{CRYPTO 2019}; \textbf{Barnes, Richard, et al.\ (2026).} \emph{Verifiable Distributed Aggregation Functions (VDAFs)}, IETF CFRG \texttt{draft-irtf-cfrg-vdaf-22} (14 August 2026; Informational; IRSG poll). \url{https://datatracker.ietf.org/doc/html/draft-irtf-cfrg-vdaf-22}.
  \item \textbf{Crites, Elizabeth, and Alistair Stewart (2025).} ``On Reed--Solomon Proximity Gaps Conjectures.'' IACR Cryptology ePrint Archive \texttt{2025/2046} (\url{https://eprint.iacr.org/2025/2046.pdf}).
  \item \textbf{Aas, Josh, Tim Geoghegan, Christopher Patton, Karl Tarbe, et al.\ (2026).} \texttt{prio} Rust crate, version 0.18.1 (implements \texttt{draft-irtf-cfrg-vdaf-18}). \url{https://crates.io/crates/prio}; source \url{https://github.com/divviup/libprio-rs}. Compiled and executed for Table~\ref{tab:prio3-bench}.
  \item \textbf{Davis, Hannah, Christopher Patton, Mike Rosulek, and Phillipp Schoppmann (2023).} ``Verifiable Distributed Aggregation Functions.'' \emph{Proceedings on Privacy Enhancing Technologies} 2023, no.~4: 578--592. DOI: \url{https://doi.org/10.56553/popets-2023-0126}. Full version with corrected proof of Theorem~1: IACR ePrint \texttt{2023/130} (\url{https://eprint.iacr.org/2023/130}).
  \item \textbf{Massimo, Jake, Panos Kampanakis, Sean Turner, and Bas Westerbaan (2025).} ``Internet X.509 Public Key Infrastructure -- Algorithm Identifiers for the Module-Lattice-Based Digital Signature Algorithm (ML-DSA).'' RFC~9881, October 2025. \url{https://www.rfc-editor.org/rfc/rfc9881}.
  \item \textbf{Ounsworth, Mike, John Gray, Massimiliano Pala, Jan Klau\ss{}ner, and Scott Fluhrer (2026).} ``Composite Module-Lattice-Based Digital Signature Algorithm (ML-DSA) for use in X.509 Public Key Infrastructure.'' IETF Internet-Draft \texttt{draft-ietf-lamps-pq-composite-sigs-19}, 21 April 2026. \url{https://datatracker.ietf.org/doc/html/draft-ietf-lamps-pq-composite-sigs-19}.
  \item \textbf{Sonnino, Alberto, Mustafa Al-Bassam, Shehar Bano, Sarah Meiklejohn, and George Danezis (2019).} ``Coconut: Threshold Issuance Selective Disclosure Credentials with Applications to Distributed Ledgers.'' \emph{NDSS 2019}. DOI: \url{https://doi.org/10.14722/ndss.2019.23272}.
  \item \textbf{Bootle, Jonathan, Vadim Lyubashevsky, Ngoc Khanh Nguyen, and Alessandro Sorniotti (2023).} ``A Framework for Practical Anonymous Credentials from Lattices.'' \emph{CRYPTO 2023}. \url{https://research.ibm.com/publications/a-framework-for-practical-anonymous-credentials-from-lattices}.
  \item \textbf{Argo, Sven, Tim G\"uneysu, Corentin Jeudy, Georg Land, Adeline Roux-Langlois, and Olivier Sanders (2024).} ``Practical Post-Quantum Signatures for Privacy.'' \emph{ACM CCS 2024}. DOI: \url{https://doi.org/10.1145/3658644.3670297}.
  \item \textbf{Leval, Gaston (1971).} \emph{Espagne Libertaire (1936--1939)}. Paris: \'Editions du Monde Libertaire. Cited through Dolgoff's (1974) translation.
  \item \textbf{Peirats, Jos\'e.} Contributions on pay scales and on the Aragon federation, in Dolgoff (1974), Chs.~6 and~9 (Ch.~9 read via the marxists.org edition).
\end{enumerate}

\vfill
\noindent\rule{\textwidth}{0.6pt}\\[0.4em]
\noindent\colorbox{NavyPrimary}{%
  \parbox{\dimexpr\textwidth-2\fboxsep\relax}{%
    \vspace{6pt}
    \centering
    {\footnotesize\bfseries\textcolor{GoldAccent}{POST-CAPITALIST LABS IMPRINT $\cdot$ WORKING DRAFT MONOGRAPH}}\\[3pt]
    {\scriptsize\color{white} Leveraging technology to amplify social movements $\cdot$ Decentralised consensus, mechanism design, and democratic economic coordination}\\[3pt]
    {\scriptsize\color{GoldAccent} \texttt{https://github.com/Post-Capitalist-Labs} $\quad\cdot\quad$ \texttt{https://github.com/titancassini}}
    \vspace{6pt}
  }%
}

\end{document}
